% This must be in the first 5 lines to tell arXiv to use pdfLaTeX, which is strongly recommended.
\pdfoutput=1
% In particular, the hyperref package requires pdfLaTeX in order to break URLs across lines.

\documentclass[11pt]{article}

% Remove the "review" option to generate the final version.
\usepackage{EMNLP2023}

% Standard package includes
\usepackage{times}
\usepackage{latexsym}
%\usepackage{subcaption}
% For proper rendering and hyphenation of words containing Latin characters (including in bib files)
\usepackage[T1]{fontenc}
\usepackage{listings}
\usepackage{amssymb}
\usepackage{tcolorbox}
% For Vietnamese characters
% \usepackage[T5]{fontenc}
% See https://www.latex-project.org/help/documentation/encguide.pdf for other character sets

% This assumes your files are encoded as UTF8
\usepackage[utf8]{inputenc}

% This is not strictly necessary, and may be commented out.
% However, it will improve the layout of the manuscript,
% and will typically save some space.
\usepackage{microtype}

% This is also not strictly necessary, and may be commented out.
% However, it will improve the aesthetics of text in
% the typewriter font.
\usepackage{inconsolata}
\usepackage{wrapfig}

% Custom packages
\usepackage{booktabs}
\usepackage[export]{adjustbox}
\usepackage{graphicx}
\usepackage{float}
\usepackage{subfig}
\usepackage{caption}
%\usepackage{subcaption}
%\usepackage{subfigure}
\usepackage{amsmath}
\usepackage{siunitx}
\usepackage{framed}
% reduce spacing in lists
\usepackage{enumitem}
% reduce spacing after figure captions
\usepackage{caption}
\usepackage{tabularx}
% \newcommand{TB}{\textbf}
% \newcommand{\R}{\textbf{R}}
% \newcommand\TBF[1]{\textbf{#1}}
% \newcommand*\TBF[1]{\textbf{#1}}
\newcommand{\TBF}[1]{\textbf{#1}}
%\usepackage{letltxmacro}% http://tex.stackexchange.com/q/88001/5764
%\LetLtxMacro{\bb}{\mathbb}
%\LetLtxMacro{\TBF}{\textbf}
% \newcommand{\TBF}[1]{\textbf{#1}}
% If the title and author information does not fit in the area allocated, uncomment the following
%
\usepackage{comment}
\usepackage{pifont}%
\newcommand{\cmark}{\ding{51}}%
\newcommand{\xmark}{\ding{55}}%
\newcommand\rwkvemoji{\raisebox{-2pt}{\includegraphics[width=1.2em]{images/rwkv_emoji.png}}}
%\setlength\titlebox{<dim>}
%
% and set <dim> to something 5cm or larger.

\setlist[2]{noitemsep} % sets the itemsep and parsep for all level two lists to 0
\setitemize{noitemsep} % sets no itemsep for enumerate lists only
\setenumerate{noitemsep} % sets no itemsep for enumerate lists only


\title{RWKV: Reinventing RNNs for the Transformer Era}

% Author information can be set in various styles:
% For several authors from the same institution:
% \author{Author 1 \and ... \and Author n \\
%         Address line \\ ... \\ Address line}
% if the names do not fit well on one line use
%         Author 1 \\ {\bf Author 2} \\ ... \\ {\bf Author n} \\
% For authors from different institutions:
% \author{Author 1 \\ Address line \\  ... \\ Address line
%         \And  ... \And
%         Author n \\ Address line \\ ... \\ Address line}
% To start a seperate ``row'' of authors use \AND, as in
% \author{Author 1 \\ Address line \\  ... \\ Address line
%         \AND
%         Author 2 \\ Address line \\ ... \\ Address line \And
%         Author 3 \\ Address line \\ ... \\ Address line}

\setlength\titlebox{7.5 cm}
%All other authors are listed alphabetically.
\author{Bo Peng\textsuperscript{1,2}\thanks{\quad Equal first authorship. Others listed alphabetically.} \;\; Eric Alcaide\textsuperscript{2,3,4}\footnotemark[1] \;\; Quentin Anthony\textsuperscript{2,5}\footnotemark[1]
\\ {
\bf Alon Albalak\textsuperscript{2,6}   Samuel Arcadinho\textsuperscript{2,7}  Stella Biderman\textsuperscript{2,8}   Huanqi Cao\textsuperscript{9}  Xin Cheng\textsuperscript{10} 
} \\ {
\bf Michael Chung\textsuperscript{11} Xingjian Du\textsuperscript{1}  Matteo Grella\textsuperscript{12}  Kranthi Kiran GV\textsuperscript{2,13}   Xuzheng He\textsuperscript{2}
} \\ {
\bf Haowen Hou\textsuperscript{14} Jiaju Lin\textsuperscript{1} Przemysław Kazienko\textsuperscript{15}   Jan Koco\'{n}\textsuperscript{15}   Jiaming Kong\textsuperscript{16} 
} \\  {
\bf Bart\l{}omiej Koptyra\textsuperscript{15}   Hayden Lau\textsuperscript{2}  Krishna Sri Ipsit Mantri\textsuperscript{17}  Ferdinand Mom\textsuperscript{18,19}
}  \\ {
\bf  Atsushi Saito\textsuperscript{2,20} Guangyu Song\textsuperscript{21}  Xiangru Tang\textsuperscript{22} Bolun Wang\textsuperscript{23} Johan S. Wind\textsuperscript{24} 
}  \\ {
\bf Stanis\l{}aw Wo\'{z}niak\textsuperscript{15}   Ruichong Zhang\textsuperscript{9}    Zhenyuan Zhang\textsuperscript{2} Qihang Zhao\textsuperscript{25,26} 
} \\ {
\bf Peng Zhou\textsuperscript{23}   Qinghua Zhou\textsuperscript{5} Jian Zhu\textsuperscript{27} Rui-Jie Zhu\textsuperscript{28,29}
} \\ 
 \small
 \textsuperscript{1}Generative AI Commons
 \textsuperscript{2}EleutherAI
 \textsuperscript{3}U. of Barcelona 
 \textsuperscript{4}Charm Therapeutics 
 \textsuperscript{5}Ohio State U.
 \textsuperscript{6}U. of C., Santa Barbara \\
 \small
 \textsuperscript{7}Zendesk
 \textsuperscript{8}Booz Allen Hamilton
 \textsuperscript{9}Tsinghua University
 \textsuperscript{10}Peking University
 \textsuperscript{11}Storyteller.io 
 \textsuperscript{12}Crisis24 
 \textsuperscript{13}New York U. \\
 \small
 \textsuperscript{14}National U. of Singapore
 \textsuperscript{15}Wroclaw U. of Science and Technology
 \textsuperscript{16}Databaker Technology
 \textsuperscript{17}Purdue U.
 \textsuperscript{18}Criteo AI Lab\\
 \small
 \textsuperscript{19}Epita 
 \textsuperscript{20}Nextremer
 \textsuperscript{21}Moves
 \textsuperscript{22}Yale U. 
 \textsuperscript{23}RuoxinTech
 \textsuperscript{24}U. of Oslo
 \textsuperscript{25}U. of Science and Technology of China\\
 \small
 \textsuperscript{26}Kuaishou Technology
 \textsuperscript{27}U. of British Columbia 
 \textsuperscript{28}U. of C., Santa Cruz
 \textsuperscript{29}U. of Electronic Science and Technology of China
}


\providecommand{\imwidth}{}

%
\newcommand{\evaluations}{
\begin{figure*}[htbp]
%
\centering
    \renewcommand
{\imwidth}{0.33\linewidth}
	\setlength{\tabcolsep}{1pt}
%
	\begin{tabular}{ccc}

	\multicolumn{3}{c}{\small{Evaluation Plots}} \\

	\includegraphics[align=c,width=\imwidth]{images/acc_arc_challenge.png} &
	\includegraphics[align=c,width=\imwidth]{images/acc_arc_easy.png} &
	\includegraphics[align=c,width=\imwidth]{images/acc_lambada.png} \\

	
	\includegraphics[align=c,width=\imwidth]{images/acc_piqa.png} &
	\includegraphics[align=c,width=\imwidth]{images/acc_sciq.png} &
	\includegraphics[align=c,width=\imwidth]{images/acc_winogrande.png} &
 \\

	\end{tabular}
	%
\caption{\label{fig:evaluations}Caption}
%
\end{figure*}
}





\begin{document}
\maketitle

\begin{abstract} \label{abstract}
Transformers have revolutionized almost all natural language processing (NLP) tasks but suffer from memory and computational complexity that scales quadratically with sequence length. In contrast, recurrent neural networks (RNNs) exhibit linear scaling in memory and computational requirements but struggle to match the same performance as Transformers due to limitations in parallelization and scalability. We propose a novel model architecture, Receptance Weighted Key Value (RWKV), that combines the efficient parallelizable training of transformers with the efficient inference of RNNs.

Our approach leverages a linear attention mechanism and allows us to formulate the model as either a Transformer or an RNN, thus parallelizing computations during training and maintains constant computational and memory complexity during inference. We scale our models as large as 14 billion parameters, by far the largest dense RNN ever trained, and find RWKV performs on par with similarly sized Transformers, suggesting future work can leverage this architecture to create more efficient models. This work presents a significant step towards reconciling trade-offs between computational efficiency and model performance in sequence processing tasks.
\footnote{Code at: \href{https://github.com/BlinkDL/RWKV-LM}{https://github.com/BlinkDL/RWKV-LM}}

\end{abstract}


%\section{Introduction} \label{introduction}
%\textcolor{blue}{Introduction is a placeholder, Alon will go back and re-do introduction once we have more concrete details of what is being done in this work.}
% What are we trying to do, and why is this relevant?
% Develop efficient sequential NN architecture (efficient in training and inference)
%Since the introduction of the transformer by Vaswani et al.~\cite{NIPS2017_3f5ee243}, its success has been attributed to the self-attention mechanism \cite{brown2020language, openai2022chatgpt,kocon2023chatgpt,openai2023gpt4,touvron2023llama,hoffmann2022training}. This powerful and flexible method allows models to capture long-range dependencies and complex interactions between input elements. The inherent parallelizability of self-attention combined with ever-increasing dataset sizes, and its ability to create rich contextual representations, has led to state-of-the-art models across a multitude of benchmarks and applications.
% Transformers have become dominant architectural paradigm, but have quadratic computational and space complexity due to self-attention mechanism
%However, despite the transformative impact of the Transformer, it is not without its shortcomings. One notable drawback is the quadratic complexity of its self-attention mechanism. It renders the architecture computationally expensive for tasks involving long input sequences and in situations with restricted memory budgets. These limitations have spurred a wealth of research aiming to improve the scaling properties of the Transformer, often at the expense of sacrificing some of the properties that make it so effective~\citep{wang2020linformer, zaheer2020bigbird, dao2022flashattention}.


% Quadratic complexity limits the use cases (on-device, conversations, etc.)
% SEE: 4.3 RNN-like Formulation for current explanation of advantages

% Why is developing efficient NN architectures hard?
%Addressing the issue of quadratic scaling has proven to be a challenging endeavor, as it requires navigating the intricate balance between the computational efficiency of the self-attention mechanism and the expressive capacity of the model. Furthermore, the non-local nature of the self-attention mechanism poses significant difficulty in devising more efficient alternatives that do not compromise the rich expressiveness of the Transformer.

% How do we solve it?
%In this paper, we propose a novel neural network architecture, dubbed the RWKV network, which serves as a variant of both the RNN and the Transformer. The RWKV network is designed to alleviate the quadratic scaling issue and replace it with a linear one while retaining the crucial properties that have made the Transformer the dominant architecture in the field
% Go more into depth on what these crucial properties are
%. By leveraging the strengths of both RNNs and Transformers, the RWKV network presents a compelling alternative for tackling a wide range of tasks that demand both computational efficiency and expressiveness. In doing so, the time and space complexity is sharply decreased. Our approach is very friendly to large model size and computation resources. 
\section{Introduction} \label{introduction}
Deep learning has greatly advanced artificial intelligence, impacting a range of scientific and industrial uses. These often involve complex sequential data processing tasks such as natural language understanding, conversational AI, time-series analysis, and indirectly sequential formats like images and graphs \cite{brown2020language, ismail2019deep, wu2020comprehensive,albalak-etal-2022-feta}. Predominant among these techniques include RNNs %convolutional neural networks (CNNs), 
and Transformers \cite{vaswani2017attention}, each with specific benefits and drawbacks. RNNs require less memory, particularly for handling long sequences.  However, they suffer from the vanishing gradient problem and non-parallelizability in the time dimension during training, limiting their scalability \cite{hochreiter1998vanishing,le2016quantifying}. %CNNs, on the other hand, are only adept at capturing local patterns, which limits their capacity to deal with long-range dependencies, crucial to many sequence processing tasks \cite{bai2018empirical}.


\begin{figure}[ht]
    \centering
    \includegraphics[width=\linewidth]{images/acc_ave.pdf}
    \caption{Average performance of RWKV models compared to transformers across twelve NLP tasks. For further details, see section \ref{evaluations}.}
    \label{fig:average}
\end{figure}

Transformers emerged as a powerful alternative, adept at managing local and long-range dependencies and supporting parallelized training \cite{tay2022efficient}. Models such as GPT-3 \cite{brown2020language}, ChatGPT \cite{openai2022chatgpt,kocon2023chatgpt}, LLaMA \cite{touvron2023llama}, and Chinchilla \cite{hoffmann2022training} showcase the potential of Transformers in NLP. However, the self-attention mechanism's  quadratic complexity makes it computationally and memory intensive for tasks involving long sequences and constrained resources. This has stimulated research to enhance Transformers' scalability, sometimes sacrificing some of their effectiveness~\cite{wang2020linformer,zaheer2020bigbird,dao2022flashattention}. 

\begin{table}
\centering
\setlength{\tabcolsep}{2pt}
\small
\begin{tabular}{lcc}
\toprule
\textbf{Model} & \textbf{Time} & \textbf{Space}\\
\midrule
Transformer & $O(T^2d)$ & $O(T^2 + Td)$ \\
Reformer & $O(T\log{T}  d)$ & $O(T\log{T} + Td)$ \\
Performer & $O(T d^2\log{d})$ & $O(Td\log{d} + d^2\log{d})$ \\
Linear Transformers & $O(T d^2)$ & $O(Td + d^2)$ \\
AFT-full & $O(T^2d)$ & $O(Td)$ \\
AFT-local & $O(T sd)$ & $O(Td)$ \\
MEGA & $O(cTd)$ & $O(cd)$\\\midrule
RWKV (ours) & $O(\mathbf{Td})$ & $O(\mathbf{d})$ \\
\bottomrule
\end{tabular}
\caption{Inference complexity comparison with different Transformers. Here $T$ denotes the sequence length, $d$ the feature dimension, $c$ is MEGA's chunk size of quadratic attention, and $s$ is the size of a local window for AFT.}
\label{tab:comparison}
\end{table}

To tackle these challenges, we introduce the Receptance Weighted Key Value (\textbf{RWKV}) model, combining the strengths of RNNs and Transformers while circumventing key drawbacks. RWKV alleviates  memory bottleneck and quadratic scaling associated with Transformers  \cite{katharopoulos2020transformers} with efficient linear scaling, while maintaining the expressive properties of the Transformer, such as  parallelized training and robust scalability. RWKV reformulates the attention mechanism with a variant of linear attention, replacing traditional dot-product token interaction with more effective channel-directed attention. This implementation, \emph{without approximation}, offers the lowest computational and memory complexity; see Table~\ref{tab:comparison}.

The motivation behind RWKV is to balance computational efficiency with expressive capacity in neural networks. It offers a solution for handling large-scale models with billions of parameters, exhibiting competitive performance at a reduced computational cost. Experiments suggest RWKV addresses  scaling and deployment challenges in AI, especially for sequential data processing, pointing towards more sustainable and efficient AI models.

Our contributions in this paper are as follows:
\begin{itemize}
 \item The introduction of RWKV, a novel architecture combining RNNs and Transformer advantages while mitigating their limitations.
 %\item A new linear attention mechanism, avoiding the quadratic complexity of Transformers.
 \item Detailed experiments demonstrating RWKV's performance and efficiency on benchmark datasets for large-scale models. %and long-range dependencies?
 \item The release of pretrained models, from 169 million to 14 billion parameters, trained on the Pile \cite{gao2020pile,biderman2022datasheet}.\footnote{\href{https://huggingface.co/RWKV}{https://huggingface.co/RWKV}}
\end{itemize}





% How do we verify that we solved it?
%We demonstrate the efficacy of our approach through a series of experiments on benchmark datasets and showcase the RWKV network's ability to \dots
%* Scale both in number of training tokens and model parameters
%* Handle long-range dependencies in text data, while keeping a linear cost on input size
%* 

% \section{Motivation} \label{motivation}

% [eric alcaide/hypnopump draft]
% TODO: cite 
% vanishing gradients DONE
% transformer effectiveness in parallelization vs RNN, CNN DONE
% attention cost DONE
% GPT / LLaMA (for popularity) DONE
% Sequential data processing covers a wide range of scientific and industrial applications, both directly (language, time, sequences) and indirectly (other modalities, potentially formulated as sequences including images, graphs). Models for sequential data have experienced an explosive rise in capabilities and popularity. Existing deep learning architectures for sequential data have known limitations: vanishing gradients for RNNs, slow training (RNNs, CNNs), memory and compute bottlenecks (transformers), lack of scaling (non-transformer models), difficulty handling long contexts (CNNs, transformers), lack of parallelization (RNNs), slow inference (transformers)...\\
% RWKV aims to combine strengths of previous techniques while minimizing drawbacks: parallel training and scaling like transformers but no compute/memory bottleneck ($O(n)$ vs $O(n^2)$); efficient inference like RNNs but no vanishing gradients in training. It also addresses the quadratic cost of attention in Transformer  with a reformulation to obtain "scalar attention" at linear cost without any approximation (i.e., low rank). While the approximation may be negligible with small models, it can become important when scaling up.
% Indeed, RWKV proves reliable for scaling up to billions of parameters, competitive with current SOTA at lower cost. As such, RWKV is likely positioned to address existing challenges in further scaling and deploying AI models.
% Replace positional embedding by recurrent architecture and custom weighting of past and present information
% Reformulation of recurrence for enhanced parallelization
% Layer stacking by token/time-mixing and channel-mixing, in residual blocks
% Increased training efficiency with custom compute kernels
% Improved training dynamics by custom embedding initialization

% The surge of deep learning techniques in recent years has led to impressive advancements in many scientific and industrial applications involving sequential data processing, such as natural language understanding \cite{brown2020language}, time-series analysis \cite{ismail2019deep}, and indirect modalities that can be reframed as sequences, including images and graphs \cite{wu2020comprehensive}. This progress has been catalyzed by deploying neural network architectures such as RNNs, convolutional neural networks (CNNs), and, more recently, Transformers \cite{NIPS2017_3f5ee243}. The development and success of advanced Transformer models such as GPT-3 \cite{brown2020language}, ChatGPT \cite{openai2022chatgpt,kocon2023chatgpt}, GPT-4 \cite{openai2023gpt4}, LLaMA \cite{touvron2023llama}, and Chinchilla \cite{hoffmann2022training}, underline the significance of this architecture in pushing the boundaries of what is possible in natural language processing and related fields. 

% However, these architectures have inherent limitations that may restrict their usage in specific scenarios. RNNs, for instance, suffer from the vanishing gradient problem \cite{hochreiter1998vanishing,le2016quantifying} and cannot parallelize during training, limiting their scalability. CNNs, while excellent at capturing local patterns, struggle with long-range dependencies, a crucial aspect of many sequence processing tasks \cite{bai2018empirical}. On the other hand, transformers have proven to be highly effective at capturing both local and long-range dependencies and allow parallelized training \cite{tay2022efficient}. Still, their self-attention mechanism comes with significant memory and computational cost, which scales quadratically with the length of the input sequence, making them unsuitable for long sequences and settings with constrained computational resources \cite{katharopoulos2020transformers}.

% To address these limitations, we introduce RWKV, a novel neural network architecture combining previous techniques' strengths while mitigating their drawbacks. RWKV offers parallelized training and robust scaling capabilities akin to Transformers but without their quadratic computational and memory bottleneck, reducing the complexity from $O(n^2)$ to $O(n)$. Moreover, it provides efficient inference, reminiscent of RNNs, without suffering from vanishing gradients during training. 

% The attention through dot-product token interaction has been eschewed in favor of channel-directed attention, meaning that a token’s attention is more meaningfully quantified by which data channels are prioritized over others. This is in contrast with Transformer architecture where attention is more directly affected by specific token interactions.

% Crucially, RWKV reformulates the attention mechanism resulting in \emph{linear attention}, which incurs a linear cost without approximation. While the effects of approximation in the attention mechanism might be negligible for small models, they become increasingly significant when scaling up.

% This work demonstrates that RWKV can handle large-scale models with billions of parameters and is competitive with state-of-the-art architectures at a significantly reduced computational cost. Our experimental results suggest that RWKV could be a powerful tool to tackle the ongoing challenges in scaling and deploying AI models in various domains, particularly those involving sequential data processing. Thus, RWKV paves the way toward more sustainable and efficient AI models for sequence processing tasks by bridging the gap between computational efficiency and expressive capacity.

\section{Background} \label{background}
%This section will discuss relevant background information on transformers/RNNs (including equations, motivations, and downsides of each) so that when we discuss the RWKV model in the next section, it will be obvious how our method is different.
Here we briefly review the fundamentals of RNNs and Transformers. 

%\begin{figure*}[hbt!]
%\centering
%\subfloat[RNN]{%
%  \includegraphics[height=2.22cm]{images/Comparison-RNNs-Cells-ver5.eps}%
%  \label{fig:rnns-cells}%
%}\quad
%\subfloat[QuasiRNN \cite{RNN_QuasiRNN17}]{%
%  \includegraphics[height=2.22cm]{images/Comparison-RNNs-QRNN-ver5.eps}%
%  \label{fig:rnns-qrnn}%
%}\quad
%\subfloat[RWKV]{%
%  \includegraphics[height=2.22cm]{images/Comparison-RNNs-RWKV-ver5.eps}%
%  \label{fig:rnns-rwkv}%
%}
%\caption{Computation structure of the RWKV in comparison to QRNN and RNN (Vanilla, LSTM, GRU, etc) architectures. Color codes: orange indicates time-mixing, convolutions or matrix multiplications, and the continuous block indicates that these computations can proceed simultaneously; blue signifies parameterless functions that operate concurrently along the channel or feature dimension (element-wise). Green indicates channel-mixing.}
%\label{fig:rnn_archs_background}
%\end{figure*}

\subsection{Recurrent Neural Networks (RNNs)}
%In this subsection, the reason why parallelizing RNNs is difficult and the way how QRNNs parallelize recurrences are explained.% RNNs are mentioned as follows.
% this RNN formula is not helpful, better start from LSTM
% \begin{align}
% o_{t}, s_{t} = RNN(x_{t}, s_{t-1}), \label{rnn-eq:general}
% \end{align}
% where $x_{t},o_{t}, s_{t}$ are the input, the output, and the state for each time step $t$, respectively.

Popular RNN architectures such as LSTM \cite{RNN_LSTM97} and GRU \cite{RNN_GRU2014} are characterized by the following formulation (shown for LSTM, others can be reasoned similarly):
% For LSTM \cite{RNN_LSTM97}, one can take $(h_t, c_t)$ as $s_{t}$ in equation \ref{rnn-eq:general} 
\begin{align}
f_t &= \sigma_g(W_{f} x_t + U_{f} h_{t-1} + b_f), \label{eq:lstm-1}\\
i_t &= \sigma_g(W_{i} x_t + U_{i} h_{t-1} + b_i), \\
o_t &= \sigma_g(W_{o} x_t + U_{o} h_{t-1} + b_o), \\
\tilde{c}_t &= \sigma_c(W_{c} x_t + U_{c} h_{t-1} + b_c), \\
c_t &= f_t \odot c_{t-1} + i_t \odot \tilde{c}_t,\\
h_t &= o_t \odot \sigma_h(c_t). \label{eq:lstm-6}
\end{align}
% For GRU \cite{RNN_GRU2014}, one can take $h_t$ as $s_{t}$ in equation \ref{rnn-eq:general} 
% \begin{align}
% z_t &= \sigma_g(W_{z} x_t + U_{z} h_{t-1} + b_z), \\
% r_t &= \sigma_g(W_{r} x_t + U_{r} h_{t-1} + b_r), \\
% \hat{h}_t &= \phi_h(W_{h} x_t + U_{h} (r_t \odot h_{t-1}) + b_h), \\
% h_t &=   (1-z_t) \odot h_{t-1} + z_t \odot  \hat{h}_t.
% \end{align}
% The data flow of RNNs is shown in Fig. \ref{fig:rnns-cells}. 
Although RNNs can be factored into two linear blocks ($W$ and $U$) and an RNN-specific block (\ref{eq:lstm-1})--(\ref{eq:lstm-6}), as noted by \citet{RNN_QuasiRNN17},
% [original sentence]: the computation at each time step still relies on the previous timestep's outcomes.
% [Reason for change]: Do not copy sentences from the QRNN paper
the data dependency relying on previous time steps prohibits parallelizing these typical RNNs.

% QRNN architecture breaks this dependency via convolutional and pooling operations. For a given input sequence $X = (x_1,...,x_T)$, a convolutional operator with filter $ W $ denoted as $ W \star X$ executes convolutions in the timestep dimension with a bank of $m$ filters and outputs a sequence $Z \in \mathrm{R}^{m \times T}$ of $m$-dimensional candidate vectors $z_t$. To be effective for tasks involving the prediction of the next token, the filters must prohibit the calculation at a particular time step to access data from future time steps. In particular, when utilizing filters with a width of $k$, the value of $z_t$ at time step $t$ depends solely on the input sequence $x_{t-k+1}$ through $x_t$. Let $W_z, W_f, W_o$ be convolutional filter banks of time-convolution. Quasi-RNN with \textit{fo-pooling} is formulated as follows:
% \begin{align}
% Z &= \tanh \left(W_z \star  X\right), \\
% F &= \sigma(W_{f} \star   X), \\
% O &= \sigma\left(W_{o} \star  X\right), \\
% c_t &= f_t \odot c_{t-1}+\left(1-f_t\right) \odot z_t \label{eq:qrnn-fo-pooling-c-is}, \\
% h_t &= o_t \odot c_{t}, \label{eq:qrnn-fo-pooling-h-is}
% \end{align}
% where $h_t, c_t, z_t$ are states, cells, and convolution gate's values, respectively. 
% Note that the convolutional operator $\star$ and \textit{fo-pooling} functions in (\ref{eq:qrnn-fo-pooling-c-is}) and (\ref{eq:qrnn-fo-pooling-h-is}) are all parallelizable  along timesteps with filter size $k$. Thus, recurrences are computed in parallel. 
% which is a parallelized version of RNNs through convolutional layers and pooling functions

% \begin{figure}[htp!]
% \subfloat[Vanilla RNN \cite{Rumelhart_Hinton_Williams_1985}, LSTM \cite{10.1162/neco.1997.9.8.1735}, GRU \cite{RNNMT_NIPS2014}, etc    .]{%  
%   \includegraphics[clip,width=\columnwidth]{images/Comparison-RNNs-Cells-ver3.png}%
% }\label{fig:rnns-cells}
% \subfloat[QuasiRNN \cite{RNN_QuasiRNN17}]{%
%   \includegraphics[clip,width=\columnwidth]{images/Comparison-RNNs-QRNN-ver3.png}%
% }\label{fig:rnns-qrnn}
% \subfloat[RWKV\label{fig:rnns-rwkv}]{%
%   \includegraphics[clip,width=\columnwidth]{images/Comparison-RNNs-RWKV-ver3.png}%
% }
% \caption{Depiction of the computation structure of the RWKV in comparison to QRNN and RNN such as LSTM and GRU architectures. }\label{fig:rnns-cells}

% \end{figure}
\subsection{Transformers and AFT} \label{sec:background-attn} % TODO: change the title
Introduced by \citet{vaswani2017attention}, Transformers are a class of neural networks that have become the dominant architecture for several NLP tasks. Instead of operating on sequences step-by-step like RNNs, Transformers rely on attention mechanisms to capture relationships between all input and all output tokens:
% The self-attention mechanism faces the challenge of loss of positional information that is crucial when dealing with sequential data, so it is combined with positional encoding to maintain the order of tokens in the sequence. The architecture is composed of stacked encoder and decoder layers which facilitates the efficient parallel processing of large amounts of data.
\begin{align}
\mathrm{Attn}(Q, K, V) = \mathrm{softmax}(QK^\top)V,
\end{align}
where the multi-headness and scaling factor $\frac{1}{\sqrt{d_k}}$ is omitted for convenience. The core $QK^\top$ multiplication is an ensemble of pairwise attention scores between each token in a sequence, which can be decomposed as vector operations:
\begin{align}
\mathrm{Attn}(Q, K, V)_t &= \frac{\sum_{i=1}^{T} e^{q^{\top}_t k_i} \odot v_i}{\sum_{i=1}^{T} e^{q^{\top}_t k_i}}.
% \text{(Attention mask applied)} &= \frac{\sum_{i=1}^{t} e^{q^{\top}_t k_i}v_i}{\sum_{i=1}^{t} e^{q^{\top}_t k_i}}.
\end{align}

AFT \cite{zhai2021attention}, alternately formulates
\begin{align}
% \text{Attn}^{+}(W,K,V) &= \text{softmax}(W+K)V, \\
\mathrm{Attn}^{+}(W,K,V)_t = \frac{ \sum_{i=1}^{t} e^{w_{t,i}+k_i} \odot v_i}{\sum_{i=1}^{t} e^{w_{t,i}+k_i}},
\end{align}
where $\{w_{t,i}\}\in R^{T\times T}$ is the learned pair-wise position biases, and each $w_{t,i}$ is a scalar.

Inspired by AFT, RWKV takes a similar approach. However, for simplicity, it modifies the interaction weights so that it can be transformed into an RNN. Each $w_{t,i}$ in RWKV is a channel-wise time decay vector multiplied by the relative position and traced backward from current time as it decays:
\begin{align}
    w_{t,i} = -(t-i)w,
\end{align}
where $w\in (R_{\geq 0})^{d}$, with $d$ the number of channels. We require $w$ to be non-negative to ensure that $e^{w_{t,i}} \le 1$ and the per-channel weights decay backwards in time.
% \subsection{Attention Free Transformers}
% [TODO???] Is this needed? @BlinkDL said that this model inspired RWKV
% better to be briefly touched in Related Works section. No need to focus on AFT, just say it's related work, and ours is better.





\section{RWKV} %\label{rwkv_model}
\label{rwkv_model}

% refer to : 
% https://johanwind.github.io/2023/03/23/rwkv_overview.html
% https://johanwind.github.io/2023/03/23/rwkv_details.html 

% refer to https://cdn.discordapp.com/attachments/1103039376184852622/1105118955917541486/Fr1Q5zDaUAAWRWd.png for technical details 
% refer to https://twitter.com/BlinkDL_AI/status/1655581735906598919

The RWKV model architecture is defined by four fundamental elements that are intrinsic to the time-mixing and channel-mixing blocks:
\begin{itemize}
    \item $R$: The \textbf{Receptance} vector acts as the receiver of past information.
    \item $W$: The \textbf{Weight} signifies the positional \\ weight decay vector, a trainable parameter within the model.
    \item $K$: The \textbf{Key} vector performs a role analogous to $K$ in traditional attention mechanisms. 
    \item $V$: The \textbf{Value} vector functions similarly to $V$ in conventional attention processes.
\end{itemize}
These core elements interact multiplicatively at each timestep, as depicted in Figure \ref{fig:blocks}.

% Moved figure further one up to improve readability. Otherwise, page 4 would have nothing but figures and equations

\begin{figure}
    \centering
    % \includegraphics[width=0.48\linewidth]{images/RWKV-block.eps}
    \raisebox{+0.135\height}{
        \includegraphics[width=0.5\linewidth]{images/RWKV-block.pdf}
    }
    % \includegraphics[width=0.4\linewidth]{images/RWKV-full.eps}
    \includegraphics[width=0.4\linewidth]{images/RWKV-full.pdf}
    \caption{Elements within an RWKV block (left) and the complete RWKV residual block, equipped with a final head for language modeling (right).}
    \label{fig:blocks}
\end{figure}

\subsection{Architecture}

The RWKV model is composed of stacked residual blocks. Each block consists of a time-mixing and a channel-mixing sub-block, embodying recurrent structures to leverage past information.

This model uses a unique attention-like score update process, which includes a time-dependent softmax operation improving numerical stability and mitigating vanishing gradients (for rigorous proof, see Appendix \ref{sec:appendix-gradstab}). It ensures that the gradient is propagated along the most relevant path. Additionally, layer normalization \cite{ba2016layer} incorporated within the architecture aids in stabilizing the gradients, effectively addressing both vanishing and exploding gradient issues. These design elements not only enhance the training dynamics of deep neural networks but also facilitate the stacking of multiple layers, leading to superior performance over conventional RNN models by capturing complex patterns across different levels of abstraction (see also Appendix \ref{sec-appen:vis}).

\begin{figure}[htp]\label{rwkv_arch}    
    \centering
    \includegraphics[width=\linewidth]{images/RWKV-arch-ver5.pdf}    
    \caption{RWKV architecture for language modeling.}
    \label{fig:arch}
\end{figure}

\subsubsection{Token Shift} \label {tokenshift}

% In this architecture, recurrence denotes a linear interpolation between current and previous timestep inputs, facilitating a token shift. This shift can adjust independently for each input embedding linear projection, including $R$, $K$, $V$ in time-mixing, and $R$, $K$ in channel-mixing.

In this architecture, all linear projection vectors ($R$, $K$, $V$ in time-mixing, and $R'$, $K'$ in channel-mixing) involved in computations are produced by linear interpolation between current and previous timestep inputs, facilitating a token shift.

The vectors for time-mixing computation are linear projections of linear combinations of the current and previous inputs of the block:
\begin{align}
\label{eq:time-mix1}
r_t &= W_r \cdot (\mu_r \odot x_t + (1 - \mu_r) \odot x_{t-1} ), \\
\label{eq:time-mix2}
k_t &= W_k \cdot (\mu_k \odot x_t + (1 - \mu_k) \odot x_{t-1} ), \\
\label{eq:time-mix3}
v_t &= W_v \cdot (\mu_v \odot x_t + (1 - \mu_v) \odot x_{t-1} ),
\end{align}
as are the channel-mixing inputs:
\begin{align}\label{eq:channel_mix_eqs1}
r'_t &= W'_r \cdot (\mu'_r \odot x_t + (1 - \mu'_r) \odot x_{t-1} ), \\
\label{eq:channel_mix_eqs2}
k'_t &= W'_k \cdot (\mu'_k \odot x_t + (1 - \mu'_k) \odot x_{t-1} ).
\end{align}

The token shift is implemented as a simple offset in the temporal dimension at each block using the PyTorch \citep{paszke2019pytorch} library as \texttt{nn.ZeroPad2d((0,0,1,-1))}.

\subsubsection{WKV Operator}

The computation of the $WKV$ operator in our model parallels the method used in Attention Free Transformer (AFT) \cite{zhai2021attention}. However, unlike AFT where $W$ is a pairwise matrix, our model treats $W$ as a channel-wise vector that is modified by relative position. In our model, this recurrent behavior is defined by the time-dependent update of the $WKV$ vectors, formalized in the following equation:

\begin{align}
    \label{eq:nom-denom}
    wkv_t &= \frac{ \sum_{i=1}^{t-1} e^{-(t-1-i)w+k_i} \odot v_i + e^{u+k_t} \odot v_t }{\sum_{i=1}^{t-1} e^{-(t-1-i)w+k_i} + e^{u+k_t}}.
\end{align}

% matrix as formalized in equation \ref{eq:nom-denom}.
% [Comment]: I suggest remove the sentence at all, because time-shift mixing is already explained (as a linear interpolation between the current input and the input at the previous time step), and leave the readers a little mystery temporarily
% Time-shift mixing allows the model to better predict the next token (obvious next token) and also helps to integrate better all previous information.
To circumvent any potential degradation of $W$, we introduce a vector $U$ that separately attends to the current token. More information about this can be found in Appendix \ref{sec-appen:vis}.
%\citet{zhai2021attention}, dealing with all previous information.

\subsubsection{Output Gating}

Output gating is implemented in both time-mixing and channel-mixing blocks using the sigmoid of the receptance, $\sigma(r)$. The output vector $o_t$ post the $WKV$ operator is given by:

\begin{align}
\label{eq:time-mix4}
o_t &= W_o \cdot (\sigma(r_t) \odot wkv_t).
\end{align}

In the channel-mixing block, a similar operation is performed:
\begin{align}
\label{eq:channel_mix_eqs3}
o'_t &= \sigma(r'_t) \odot (W'_v \cdot \max({k'_t}, 0)^2),
\end{align}
where we adopt the squared ReLU activation function \cite{DBLP:journals/corr/abs-2109-08668}.

% but we introduce a special vector $U$ for the current token to compensate for the potential degeneration of $W$ (Appendix \ref{sec-appen:vis}), which is now a vector instead of the pairwise matrix in \citet{zhai2021attention}, dealing with all previous information. % squared ReLU activation \cite{DBLP:journals/corr/abs-2109-08668}.

\subsection{Transformer-like Training} \label{transformer_parallel}

% RWKV's architecture enables efficient parallelization of operations, reminiscent of transformer-like parallelization.
RWKV can be efficiently parallelized using a technique called \textit{time-parallel mode}, reminiscent of Transformers. The time complexity of processing a batch of sequences in a single layer is $O(BTd^2)$, primarily consisting of matrix multiplications $W_\lambda$, where $\lambda\in \{r,k,v,o\}$ (assuming $B$ sequences, $T$ maximum tokens, and $d$ channels).
% First, we leverage that the channel-mixing block does not require any information from previous states, and thus can be applied in parallel to all time steps, greatly reducing the total execution time when the sequence is known in advance.
% (benefiting both the training phase and the processing of the sequence before autoregressive generation)
% The time-mixing block is also in this sense parallelizable to the computation of \textit{key}, \textit{value}, and \textit{receptance} vectors but then requring a sequential scan in updating attention scores $wkv_t$, $a_t$, $b_t$, $p_t$ (see Appendix \ref{sec:appendix-rnn}).
In contrast, updating attention scores $wkv_t$ involves a serial scan
%of $a_t$ and $b_t$
(see Appendix \ref{sec:appendix-rnn} for more detail) and has complexity $O(BTd)$.
% We call this approach \textit{time-parallel mode} as the model's temporal context bypasses the inherently sequential nature of a recurrent network that in theory precludes parallelization.
% This mode can be leveraged to train LLMs efficiently for next token prediction analogous to GPTs \cite{Radford2019LanguageMA_gpt2}.

%In detail, for efficient use of hardware during training and prompt processing in inference, we perform the computation of nearly all components in the model in parallel along the time axis, forming what we call ``time-parallel mode''. As the first recurrence,
The matrix multiplications can be parallelized similarly to $W_{\lambda}$, where $\lambda \in \{Q,K,V,O\}$ in conventional Transformers. The element-wise $WKV$ computation is time-dependent but can be readily parallelized along the other two dimensions \cite{lei-etal-2018-simple}\footnote{For extremely long sequences, more sophisticated methods such as \citet{Martin2017ParallelizingLR} that parallelize over sequence length could be used.}. %Additionally, to facilitate efficient training, we also implement token shifting as mentioned in \ref{tokenshift}.
% The other recurrent component $wkv$, which is sequentially carried out as described in eq.~\ref{eq:nom-denom}, is done as a custom CUDA kernel for reduced overhead. All the remaining operations, including the heaviest linear projections in time-mixing and channel-mixing blocks, are timestep-wise and can be implemented in parallel along the time axis, similar to well-known transformer building blocks.

\subsection{RNN-like Inference} \label{rnn_inference}

Recurrent networks commonly utilize the output at state $t$ as input at state $t+1$. This usage is also observed in the autoregressive decoding inference of language models, where each token must be computed before being passed to the next step. RWKV takes advantage of this RNN-like structure, known as \textit{time-sequential mode}. In this context, RWKV can be conveniently formulated recursively for decoding during inference, as demonstrated in Appendix~\ref{sec:appendix-rnn}.

% The state (or \textit{context}) acts as a summary or compressed representation of the prior inputs and captures relevant information needed to generate the subsequent output.
%RWKV incorporates three mechanisms to capture and propagate sequential information: recurrence, time decay, and token shift. These elements allow the model to efficiently process longer sequences, bypassing the limitations of traditional self-attention mechanisms that typically require a KV cache growing linearly with sequence length \cite{Beltagy2020Longformer, reformer, zaheer2021big}.


\subsection{Additional Optimizations} \label{adding_opts}

% In general, see https://github.com/BlinkDL/RWKV-LM/tree/main#how-it-works
% explained below

\paragraph{Custom Kernels} \label{custom_kernel}

% FROM https://johanwind.github.io/2023/03/23/rwkv_details.html
% My simplified code processes the tokens one by one, which is much slower than processing them in parallel, especially when running on GPUs. For inference, there is no way around this, as we need to sample a token before we can use it to calculate the next one. However, for training, all the text is already available. This lets us parallelize across tokens. Most of the code is fairly straightforward to parallelize like this, as there is little dependence through time. For example, all the expensive matrix multiplications work on each token independently, leading to good performance.

% However, the RWKV attention is inherently sequential. Fortunately, it has very little computation (on the order of 1024 times less than the matrix multiplications), so it should be fast. Sadly, PyTorch does not have a good way of handling this sequential task, so the attention part becomes slow (even compared to the matrix multiplications). Therefore, I wrote optimized CUDA kernels for computing the RWKV attention, which has been my main contribution to the RWKV project.

% JAX has jax.lax.scan and jax.lax.associative_scan, which allows a pure JAX implementation to perform better than pure pytorch. However, I still estimate that JAX would lead to about 40% slower training compared to CUDA (that estimate may be outdated, as it was made for training a relatively small 1.5B model).

To address inefficiencies in the $WKV$ computation arising from the sequential nature of the task when using standard deep learning frameworks, we have developed a custom CUDA kernel.
This kernel enables the execution of a single compute kernel on training accelerators, while all other parts of the model, such as matrix multiplications and point-wise operations, are already inherently parallelizable and efficient. 
% see for single kernel: https://github.com/BlinkDL/RWKV-LM/blob/main/RWKV-v4/cuda/wkv_cuda.cu#L7 
% for call to kernel, other ops in torch: https://github.com/BlinkDL/RWKV-LM/blob/main/RWKV-v4/src/model.py#L230


% \paragraph{FFN with $R$ gate}

% % from : cyclekiller#6762 https://discord.com/channels/729741769192767510/1103039376184852622/1106926765760331827
% Prior research \cite{DBLP:journals/corr/abs-2105-01601, liu2021pay, yu2022metaformer} suggests that the role of self-attention in Transformer-based vision tasks may not be as crucial as previously believed. While these findings provide valuable insights, completely replacing self-attention in natural language tasks could be an excessively radical approach. In our study, we take a more nuanced approach by partially deconstructing the attention mechanism. We replace the fixed $QKV$ formula with $KV$ and introduce a new time-decaying factor $W$.
% This approach enables us to incorporate token and channel-mixing components inspired by MLP-mixer \cite{DBLP:journals/corr/abs-2105-01601} and a gating unit $R$ similar to gMLP \cite{liu2021pay}. These additions enhance the performance of our RWKV model.

\paragraph{Small Init Embedding} \label{smallinit}

During the initial stage of training a transformer model \cite{vaswani2017attention}, we observe that the embedding matrix undergoes slow changes, presenting a challenge for the model to move away from its initial noisy embedding state. To address this issue, we propose an approach that involves initializing the embedding matrix with small values and subsequently applying an additional LayerNorm operation. This accelerates and stabilizes the training process, allowing for the training of deep architectures with post-LN components. The effectiveness of this approach is demonstrated in Figure \ref{fig:small_init_emb}, illustrating improved convergence by enabling the model to quickly transition away from the initially small embedding. This is achieved through small changes occurring in a single step, which subsequently lead to substantial alterations in directions and further notable changes after the LayerNorm operation.

% check: https://github.com/BlinkDL/SmallInitEmb

% Embeddings are usually initialized with a normal (0, 1) distribution
% LRs are small values, ~1e-4. Embeddings move by that much in the beginning. This results in time spent initially in moving out of suboptimal configurations
% We propose initializing the embedings to uniform values of the same order, ~1e-4 and adding a LayerNorm afterwards, which speeds up training in the first iterations, especially for Base Pair Encoding (BPE) models. 
% We find it even allows for training of postLN models, but for the rest of this paper we focus on the preLN mode. 


\paragraph{Custom Initialization}

Building on principles from previous works \citep{he2016identity, af2_2021_nature}, we adopt an initialization strategy where parameters are set to values resembling an identity mapping while breaking symmetry to establish a clear information flow. The majority of weights are initialized to zero, and linear layers do not employ biases. Detailed formulas are given in Appendix \ref{sec:initialization_formulas}. We observe that the choice of initialization plays a crucial role in both the speed and quality of convergence (refer to Appendix \ref{appendix_smallinit} for further details).
% TODO: Ablations or remove last sentence. 

\subsection{Implementation}\label{detailed_implementation}

RWKV is implemented using the PyTorch Deep Learning Library \cite{paszke2019pytorch}. We integrate additional optimization strategies inspired by DeepSpeed \cite{10.1145/3394486.3406703} into the system, improving its efficiency and scalability.

% Removed dupliate sentence on CUDA kernel
% and a custom CUDA kernel for the $WKV$ computation, as further detailed in Section \ref{custom_kernel}. Additionally, optimization strategies inspired by DeepSpeed \cite{10.1145/3394486.3406703} are integrated into the system, increasing its efficiency and scalability.

% Although the structure of RWKV resembles a general recurrent network, it is specifically oriented towards language modeling tasks. This focus is influenced by the language modeling optimizations in DeepSpeed \cite{10.1145/3394486.3406703}, which has proven effective for training Large Language Models (LLMs)

The model begins with an embedding layer, as detailed in Section \ref{smallinit}. Following this are several identical residual blocks arranged sequentially. These are depicted in Figures \ref{fig:blocks} and \ref{fig:arch} and  adheres to the principles outlined in Section \ref{tokenshift}. After the last block, a simple output projection head, consisting of a LayerNorm \cite{ba2016layer} and a linear projection, is employed for logits generation for next-token prediction and  computation of the cross-entropy loss during training.

%The embeddings obtained from the final residual block, alongside the logits, can be utilized for various downstream NLP tasks.
%We perform training in time-parallel mode (Section \ref{transformer_parallel}), while autoregressive inference and an optional chat interface
%\footnote{\href{https://github.com/BlinkDL/ChatRWKV}{https://github.com/BlinkDL/ChatRWKV}}
%make use of the time-sequential mode (Section \ref{rnn_inference}), as described by Vaswani et al. \cite{vaswani2017attention}.

% Commented out as they have been merged with 4.1 and 4.3
% \subsection{Gradient Stability and Layer Stacking}
% \label{grad-stability-layer-stacking}
% % added references, working or experiments to show the stability of gradients with ablations...

% The RWKV architecture has been designed as a fusion of both Transformers and RNNs, offering the advantage of stable gradients and deeper architectures of Transformers compared to traditional RNNs while being efficient in inference.
% % [original text]: The RWKV architecture has been designed with several key elements that effectively mitigate the vanishing and exploding gradient problems often encountered in traditional RNN architectures \cite{pmlr-v28-pascanu13}. These design choices also facilitate the stacking of multiple RWKV layers—an essential aspect of the invention of Transformers, allowing the models to achieve varying levels of generalization or abstraction of the input \cite{vaswani2017attention}.
% % [reasons for changing the original text]: LSTM also can be stacked and capture more complex patterns of the input data. See for example https://stats.stackexchange.com/questions/163304/what-are-the-advantages-of-stacking-multiple-lstms

% Prior research has attempted to address the issue of gradient stability in RNNs using various techniques, including non-saturated activation functions \cite{Chandar2019TowardsNR}, gating mechanism \cite{Gu2019ImprovingTG}, gradient clipping \cite{Pascanu2012OnTD}, and adding constraints \cite{Kanai2017PreventingGE, Miller2018StableRM}. However, these techniques have had limited success. RWKV addresses this challenge effectively by integrating a softmax function with RNN-style updates, thereby circumventing issues with gradient stability.
% % Previous work has sought to tackle the problem of gradient stability in RNNs with a variety of techniques including using non-saturated activation functions \cite{Chandar2019TowardsNR}, gating mechanism \cite{Gu2019ImprovingTG}, gradient clipping \cite{Pascanu2012OnTD}, and adding constraints \cite{Kanai2017PreventingGE, Miller2018StableRM}. 
% % While these techniques have seen little success, RWKV avoids the problem inherently by utilizing softmax in conjunction with RNN-style updates.
% % [original text]: Incorporating the alternation of time-mixing and channel-mixing blocks, the RWKV architecture adeptly blends the current input with past inputs and focuses on channel interactions. Notably, the channel-mixing operation utilizes the ReLU activation function. The ReLU function, by providing gradients that are either 0 or 1, ensures the smooth propagation of gradients across different time steps, minimizing the risk of vanishing gradients, and thereby contributing to the overall stability of the learning process.
% % [reasons for changing the original text]: ReLU activation does not help address vanishing gradients. (Update: sorry, ReLU does help, but other architectures also use ReLU. We should contrast with them, not identify with them.)

% % [original text]: This results in a smoother flow of gradients across different time steps, reducing the likelihood of vanishing or exploding gradients \cite{Kolen2001GradientFI}. Furthermore, the model employs Sigmoid-based receptance gate to dynamically adjust the influence of key and value components, selectively attending to pertinent information, and contributing to the stabilization of gradients \cite{chung2015gated}.
% % [reasons for changing the original text]: Other architectures (LSTM and GRU) also use a gating mechanism.

% The RWKV model utilizes a single-step process for updating attention-like scores, incorporating a time-dependent softmax operation that improves numerical stability and mitigates vanishing gradients (for rigorous proof, see Appendix \ref{sec:appendix-gradstab}). This operation ensures that the gradient is propagated along the most relevant path. Additionally, layer normalization \cite{ba2016layer} is a crucial aspect of the architecture that enhances the training dynamics of deep neural networks by stabilizing gradients, effectively addressing both vanishing and exploding gradient issues.

% These design elements not only contribute to the stability and learning capabilities of the RWKV architecture but also facilitate the stacking of multiple layers, surpassing the capabilities of existing RNN models.
% As a result, the model can capture more complex patterns across various levels of abstraction (see also Appendix \ref{sec-appen:vis}).
 
% \subsection{Harnessing Temporal Structure for Sequential Data Processing} % Harnessing Inherent Recurrence of RWKV for Sequential Data Processing
% \label{sec-appen:vis}

% % from: @mgrellahttps://discord.com/channels/729741769192767510/1103039376184852622/1106949316284788806

% RWKV incorporates three mechanisms to capture and propagate sequential information: recurrence, time decay, and token shift.
% % The RWKV architecture inherently captures and propagates locality information. This is achieved through the implementation of three key mechanisms: recurrence, time decay, and time-shift mixing.

% % One of the main challenges faced by the Transformer architecture is the loss of positional information happening in the self-attention crucial in dealing with sequential data processing such as natural language. This necessitates the addition of explicit positional encoding to the input ~\cite{vaswani2017attention}. Instead, the RWKV architecture inherently captures and propagate locality information. This is achieved through the implementation of three key mechanisms: recurrence, time decay, and time-shift mixing.

% The \textbf{recurrence} mechanism in the time-mixing block of RWKV forms the foundation for the model's ability to capture complex relationships among sequence elements and propagate temporal locality information.
% % The \textbf{recurrence} in the time-mixing block of RWKV is the basis for the model's capacity to capture intricate relationships between sequence elements and to propagate locality information through time.

% The \textbf{time decay} mechanism (represented by $e^{-w}$ and $e^u$ in Equation \ref{eq:nom-denom}), maintains sensitivity to the positional relationship between sequence elements. By gradually diminishing the influence of past information over time, the model preserves a sense of temporal locality and progression, crucial for sequential processing. This treatment of positional information in sequential data shares similarities with the Attention with Linear Biases (ALiBi) model \cite{ALiBi}, where the linear biases facilitate extrapolation of input length. In this context, the RWKV architecture can be perceived as a trainable version of ALiBi, seamlessly incorporating positional information without the necessity for explicit encoding. It can also be viewed as an extension of the gated convolution introduced in \citet{zhai2021attention} to the entire sequence length until a given step.  

% \label{tokenshift}The \textbf{token shift} or \textbf{time-shift mixing} (depicted by diagonal arrows in Figure \ref{fig:arch}), further contributes to the model's adaptation to sequential data.
% By linearly interpolating between the current input and the input from the previous time step, the model naturally aggregates and gates information in the input channels.
% % where linear interpolation between the current input and the previous time step input, as shown in equations \ref{eq:time-mix1}-\ref{eq:time-mix3}, enables the adjustment of past information influence on the present input. This mechanism facilitates information flow between consecutive positions in the sequence, allowing the model to efficiently handle both short- and long-range dependencies in sequential data.
% % Intuitively, by assigning each token the dual tasks of (1) aggregating all previous information and (2) predicting the next token, shifted channels can focus on the former task, enhancing information propagation. Time-shift mixing multiplies every feature dimension of the token which makes it easier for the model to store different characteristics for different dimensions of the embedding, giving it more granularity.
% The overall structure of time-shift mixing bears a resemblance to the causal convolution with no dilations in WaveNet \cite{Oord2016WaveNetAG}, a classical architecture employed for forecasting time series data. %; as both share the incorporation of both the current input and the previous input into the subsequent layer.
% % This mechanism strenA defining characteristic of the RWKV architecture, as a recurrent network, is its ability to maintain a \textit{state} (or \textit{context}) across time. This feature allows the model to efficiently retain information from prior inputs while utilizing a fixed, constant memory capacity regardless of the input sequence length. This \textit{context} acts as a summary or compressed representation of the prior inputs and captures relevant information needed to generate the subsequent output. In other words, RWKV relies solely on the \textit{context} at position \textit{t} to compute the output at position \textit{t+1}.

% % For a more comprehensive and formal description, please refer to Appendix \ref{sec:appendix-rnn}. % saves space

% % I would probably suggest to remove the following detailed explaination from here. Commenting:
% %Technically, the context consists has shape $\textit{(n\_embd, 5 $\times$ %nb\_layer)}$, where each block gets 5 vectors \textit{(xx, aa, bb, pp)} where \textit{xx} refers to previous token $x_{t-1}$ used in channel-mixing (\ref{eq:channel_mix_eqs}) and also to previous token $x_{t-1}$ used in time-mixing (\ref{eq:time-mix1}) (\ref{eq:time-mix2}) (\ref{eq:time-mix3}). \textit{aa} and \textit{bb} refers to running sum in numerator and denominator of (\ref{eq:nom-denom}) and \textit{pp} is subtracted before computing the exponential in order to stabilize the softmax computation in (\ref{eq:nom-denom}).  

% % Comment the following as not clear:
% % In the context of LLM applications, injecting the \textit{context} into the % model is equivalent to prompt engineering or p-Tuning\cite{liu2022ptuning}. % This feature enables one copy of RWKV to serve multiple domains or purposes % with an implementation of state cache, minimizing computation overhead.

% % Future work and speculation: -The state is composed by 5 vectors for each block of each stacked layer. Those vectors xx, aa, bb, pp. Empirical experiments show that those vectors have very different meanings and ranges (reference ???). -From https://github.com/BlinkDL/RWKV-LM: final hidden state（.xx .aa .bb) as a faithful sentence embedding for other tasks. Probably you should begin with .xx and .aa/.bb (.aa divided by .bb). We should probably mention the hidden state originally was 4 vectors (mathematically accurate version) and required a shift to 5 vectors for numerical stability, because during time-mixing K could get very large, so some factorization was done.
% %

\section{Trained Models and Computing Costs}\label{appendix:flop_count}

To demonstrate the scalability of RWKV, we train six models ranging from 169 million to 14 billion parameters as shown in Table \ref{tab:model_flop_count}. All models are trained for one epoch (330 billion tokens) on the Pile \citep{gao2020pile,biderman2022datasheet}.

\begin{table}[htb]
    \centering
    \begin{adjustbox}{max width=\linewidth}
        \begin{tabular}{ccccc} 
 \toprule
Name & Layers & Model Dimension & Parameters & FLOP per token \\
\midrule
\SI{169}{M} & 12 & 768 & \num{1.693e+08} & \num{2.613e+08} \\
\SI{430}{M} & 24 & 1024 & \num{4.304e+08} & \num{7.573e+08} \\
\SI{1.5}{B} & 24 & 2048 & \num{1.515e+09} & \num{2.823e+09} \\
\SI{3}{B} & 32 & 2560 & \num{2.985e+09} & \num{5.710e+09} \\
\SI{7}{B} & 32 & 4096 & \num{7.393e+09} & \num{1.437e+10} \\
\SI{14}{B} & 40 & 5120 & \num{1.415e+10} & \num{2.778e+10} \\
\bottomrule
\end{tabular}
\end{adjustbox}
\caption{RWKV model architectures and FLOP counts. Further details of these hyperparameters are elaborated upon in Appendix~\ref{appendix_hyperparam}.}
\label{tab:model_flop_count}
\end{table}

The number of parameters for each model is computed using the formula:  $\text{\# parameters} = 2VD + 13D^2L + D(11L+4)$ where $V$ = 50277 is the vocabulary size, $D$ represents the Model Dimension and $L$ corresponds to the number of layers. FLOPs is for a forward pass for one token. It was calculated as $2(2VD + 13D^2L)$, which is the twice (add and multiply) the number of parameters in linear layers. The backwards pass FLOPs can be approximated as twice that of the forward pass, giving a total of $6(2VD + 13D^2L)$ FLOP per token. Notably, this matches the standard formula for FLOP calculations in transformers \citet{kaplan2020scaling}: $\text{FLOP} = 6\cdot [\text{\# tokens}]\cdot [\text{\# parameters}]$.

\subsection{Additional Training Details}

% New paragraph on learning rate/optimizer/...
For training, we use the standard Adam optimizer without weight decay, use bfloat16 precision, and train with a context length of 1024 tokens. Further details on hyperparameters are in Appendix~\ref{appendix_hyperparam}. Diverting from standard practice for transformers, we apply exponential decay to our learning rate. We also incorporate the auxiliary loss introduced by PaLM \cite{PaLM}, supplementing the standard cross-entropy loss function. This auxiliary loss encourages the softmax normalizer to approximate zero closely. As for the learning rate schedule, it remains constant for the initial iterations, and subsequently decays exponentially. 

% TODO: Add the appendix and ref here

\subsection{Scaling Laws}

Scaling laws \cite{kaplan2020scaling,henighan2020scaling,hoffmann2022training,muennighoff2023scaling} in language models refer to the mathematical relationships that describe how the performance of a language model changes with respect to various factors. These factors can include the model size ($N$), dataset size ($D$), or the optimally allocated compute budget ($C_{\rm min}$). Scaling laws are important for two primary reasons: they allow us to make predictions and plans regarding the costs and performance of large models before they are trained via interpolation and extrapolation \cite{black2022gpta,le2022language} and the contexts in which they fail provides rich feedback on important areas for future research \cite{Wei2022EmergentAO,biderman2023emergent}.

Previous work on scaling laws for RNNs has claimed that LSTMs do not strictly follow the same log-log linear scaling that transformers do \cite{kaplan2020scaling}.
We train 45 RWKV models for a variety of pairs (dataset, parameters) and find that RWKV \textit{does} follow the same general form of the scaling law that is well established for transformers. Figure \ref{fig:scaling} shows our results for loss as a function of compute, with the linear fit to the Pareto optimal points holding an $r^2$ value of $0.994$. Even when we extrapolate our curve an additional order of magnitude (blue), we find an extremely good fit with an $r^2$ of $0.875$.

\begin{figure}[hb]
    \includegraphics[width=\columnwidth]{images/scaling_compute_loss_ood.pdf}
    \caption{Scaling laws curves for RWKV models}\label{fig:scaling}
\end{figure}

\section{Evaluations}\label{evaluations}

Having demonstrated the scalability of RWKV models in the previous section, we now turn our attention to their competitiveness with traditional transformers. We focus on two questions:

\paragraph{Competitiveness} Is RWKV competitive against quadratic transformer architectures with the same amount of compute?

\paragraph{Long Context} Does increasing the context length of RWKV yield better language modeling loss when RWKV models are trained for context lengths that most open-sourced quadratic transformers \textit{cannot} efficiently process?

\subsection{NLP Evaluations}

\begin{figure*}[!ht]
\centering
\subfloat[ARC (Challenge)]{%
\includegraphics[width=0.32\textwidth]{images/acc_arc_challenge_main.pdf}
}
\subfloat[HellaSwag]{%
\includegraphics[width=0.32\textwidth]{images/acc_hellaswag_main.pdf}
}
\subfloat[LAMBADA (OpenAI)]{%
\includegraphics[width=0.32\textwidth]{images/acc_lambada_openai_main.pdf}
}\\
\subfloat[OpenBookQA]{%
\includegraphics[width=0.32\textwidth]{images/acc_openbookqa_main.pdf}
}
\subfloat[ReCoRD]{%
\includegraphics[width=0.32\textwidth]{images/acc_record_main.pdf}
}
\subfloat[Winogrande]{%
\includegraphics[width=0.32\textwidth]{images/acc_winogrande_main.pdf}
}
\caption{Zero-Shot Performance of RWKV on common language modeling evaluation benchmarks. Additional plots can be found in Appendix \ref{sec-appen:eval-detail}.} 
\label{fig:eval}
\end{figure*}


To demonstrate that RWKV is competitive with traditional transformers at NLP tasks, we compare with similarly sized models trained for a similar number of tokens (Pythia \cite{biderman2023pythia}, OPT \cite{zhang2022opt} and BLOOM \cite{scao2022bloom}). All RWKV models were trained for one epoch on the Pile (330B tokens), which is close but not identical to the amount of tokens the Pythia, OPT, and BLOOM models were trained for. Consequently, we compare our models on a \textit{FLOP-matched basis}. We avoid comparing with model trained in the Chinchilla-optimal regime \citep{hoffmann2022training} or the overtrained regime \citep{touvron2023llama} to ensure the most equitable comparison.

We report results on ARC (both Easy and Challenge) ~\cite{ARC2018}, BoolQ ~\cite{clark2019boolq}, COPA ~\cite{COPA2011}, HeadQA ~\cite{HeadQA2020}, HellaSwag ~\cite{HellaSwag2019}, LAMBADA~\cite{LAMBADAdataset}, OpenBookQA ~\cite{OpenBookQA2018}, PIQA~\cite{Bisk2020}, ReCoRD ~\cite{ReCord}, SciQ ~\cite{SciQ2017}, and Winogrande ~\cite{Wino2020}. Figure \ref{fig:average} shows the average results across all benchmarks. Some individual benchmarks are shown in Fig \ref{fig:eval}, with the rest in Appendix \ref{sec-appen:eval-detail}.

Additionally, we carried out comparative studies on RWKV and ChatGPT / GPT-4, see Appendix \ref{appendix:promptImportanceAndRWKVvsGPT}. They revealed that RWKV is very sensitive to prompt engineering. When the prompts were adjusted (re-ordered) from the ones used for GPT to more suitable for RWKV, the performance (F1) increased even from 44.2\% to 74.8\%. For sarcasm detection, RWKV outperformed ChatGPT, but was still slightly worse than the SOTA solution.


\subsection{Extended Context Finetuning}

Unlike transformers, RNNs do not have a pre-defined sequences length when they are created. However in order to efficient make use of compute we nevertheless need to preprocess the training data into contexts of the same length. We find that we are able to teach the model how to efficiently handle substantially larger batch sizes by finetuning with progressively increasing sequence length. Specifically, we first double the sequence length from 1024 to 2048 and finetune for 10B tokens from the original pretraining corpus, then we double again to 4096 for 100B tokens from the same corpus, and finally double to 8192 tokens for another 100B tokens from the same corpus. In  Fig.~\ref{fig:ctxlen_rwkv_loss} we show that increasing context length leads to lower test loss on the Pile, an indication that RWKV can make effective use of long contextual information.

\begin{figure}[ht]
    \centering
    \includegraphics[width=0.9\columnwidth]{images/rwkv_ctxlen_loss.pdf}
    \caption{RWKV shows decreasing mean test loss as a function of context length on the Pile \cite{gao2020pile} }
    \label{fig:ctxlen_rwkv_loss}
\end{figure}

\subsection{Long Context Benchmarks}

Additionally, we evaluate our model's ability to handle very long sequences by comparing to state-of-the-art long sequence models on the Long-Range Arena (LRA) benchmark \cite{tay2021long}. LRA is designed to assess the performance of models in handling lengthy context situations. It includes a collection of tasks with sequences ranging from 1,000 to 16,000 tokens, covering various types of data like text, natural language, synthetic images, and mathematical expressions. We apply RWKV on the LRA benchmark and the results are in Appendix \ref{lorfa}. The results show that RWKV performs second only to the S4 model in five datasets.


%\begin{figure*}[htp!]
%     \centering
%     \subfloat[]{%
%     \includegraphics[scale=0.5, width=\columnwidth]{images/acc_arc_challenge.png}
%          \caption{Caption}
%     \label{fig:eval_arc_challenge}
%     }

%     \centering
%     \includegraphics[scale=0.5, width=\columnwidth]{images/acc_arc_easy.png}
     % \caption{Caption}
     % \label{fig:my_label}

%     \centering
%     \includegraphics[scale=0.5, width=\columnwidth]{images/acc_lambada.png}
%     \caption{Caption}
%     \label{fig:my_label}

%      \centering
%      \includegraphics[scale=0.5, width=\columnwidth]{images/acc_piqa.png}
%      \caption{Caption}
%      \label{fig:my_label}

%     \centering
%     \includegraphics[scale=0.5, width=\columnwidth]%{images/acc_sciq.png}
%     \caption{Caption}
 %    \label{fig:my_label}

  %   \centering
 %    \includegraphics[scale=0.5, width=\columnwidth]{images/acc_winogrande.png}
%      \caption{Caption}\label{fig:my_label}
%\end{figure*}






%Focus on equi-token equi-parameter comparisons probably?

%Compare to:
%\begin{itemize}
%    \item Pythia
%    \item OPT
%    \item BLOOM
%\end{itemize}

% refer to: https://github.com/BlinkDL/RWKV-LM/blob/main/RWKV-eval.png
% with source: https://discord.com/channels/992359628979568762/1084725482035097611/1104117503967375451

% And softer comparisons to other OpenSource models like RedPajama, LLaMA, etc. 
\begin{comment}
\section{Scaling Laws} \label{scaling laws}

% TODO: add a paragraph to explain what is scaling law
Scaling laws \cite{kaplan2020scaling,henighan2020scaling,hoffmann2022training,muennighoff2023scaling} in language models refer to the mathematical relationships that describe how the performance of a language model changes with respect to various factors. These factors can include the model size ($N$), dataset size ($D$), or the optimally allocated compute budget ($C_{\rm min}$). Scaling laws are important for two primary reasons

\begin{enumerate}
    \item They allow us to make predictions and plans regarding the costs and performance of large models before they are trained via interpolation and extrapolation \cite{black2022gpta,le2022language}
    \item The contexts in which they fail provides rich feedback on important areas for future research \cite{Wei2022EmergentAO,biderman2023emergent}
\end{enumerate}

We have observed that the RWKV model also follows scaling laws and the test loss $L$ can be estimated using a power-law relationship. This observation is depicted in Figure \ref{fig:scaling-laws}.

We conducted experiments on a range of RWKV models with different model sizes, ranging from 169M to 14B parameters. In doing so, we discovered the first scaling law:
\begin{align}
L(N) = \left(N_{\mathrm{c}}/N\right)^{\alpha_N}, \label{eq:scale_law_N}
\end{align}
where $\alpha_N \sim 0.072$ and $N_{\mathrm{c}} \sim 3.1 \times 10^{13}$ (non-embedding parameters). This scaling law demonstrates that as the model size increases, the test loss decreases without encountering any bottlenecks or plateaus, indicating that the RWKV architecture exhibits excellent scalability.

When the dataset size increases by an order of magnitude, we also observe a linear decrease in loss, which represents the second scaling law:
\begin{align}
L(D) = \left(D_{\mathrm{c}}/D\right)^{\alpha_D}, \label{eq:scale_law_N_2}
\end{align}
where $\alpha_D \sim 0.066$ and $D_{\mathrm{c}} \sim 1.4 \times 10^{15}$ (tokens).

The amount of computation is determined by two factors: the dataset size and the model size. Therefore, when there is an ample amount of data and an appropriate model size, the computational cost also follows the scaling law:
\begin{align}
L(C) = \left(C_{\mathrm{c}}/C\right)^{\alpha_C} \label{eq:scale_law_N_3},
\end{align}
where $\alpha_C \sim 0.053$ and $C_{\mathrm{c}} \sim 1.1 \times 10^{7}$ (PF-days). This indicates that with sufficient computational resources, combined with an ample amount of data and an appropriate model size, the RWKV model can continue to scale.

The next question we attempt to answer is how to make the trade-off between model size and dataset size when only a limited computing budget is available.


% TODO: add equation to scaling laws

% TODO: this figure needs commentary
\begin{figure*}
    \centering
    \includegraphics[width=0.32\linewidth]{images/loss_compute.png}
    \includegraphics[width=0.32\linewidth]{images/loss_tokens.png}
    \includegraphics[width=0.32\linewidth]{images/loss_params.png}
    \caption{Scaling Laws plot for RWKV models}
    \label{fig:scaling-laws}
\end{figure*}

% mgrella: I propose to cut off section 8. To make it effective I would insert comparison graphs for each experimented task but not bringing significant value at the end. 
% \section{Fundamental Experiments}
% While the RWKV model was primarily developed to compete with Transformer architectures, it is essential to evaluate its performance on typical benchmarks for recurrent networks. This section compares our RWKV architecture with well-established RNN architectures from the current literature, including LSTM and GRU.
% We used synthetic datasets with sequence lengths from 1,000 to 10,000 randomly generated floating-point numbers, challenging for LSTM and GRU due to the vanishing gradient problem.
% \textbf{Copy Task}: Models reproduce the same sequence as input.
% \textbf{Adding Task}: Given a sequence and two indices, models output the sum of the two corresponding numbers.
% \textbf{Temporal Order Verification}: Models determine if sequence elements are in ascending order.
\end{comment}
\section{Inference Experiments} \label{infer_experiments}

We benchmark inference requirements according to size and family. Specifically, we evaluate text generation speed and memory requirements on typical compute platforms including CPU (x86) and GPU (NVIDIA A100 80\,GB). For all of our inference experiments we use float32 precision and the HuggingFace Transformers~\cite{hf_transformers}. We include all model parameters in the parameter count, including both embedding and non-embedding layers. Performance under different quantization setups is left to further work. See Appendix \ref{appendix:inference_results} for more results.

\begin{figure}[H] % btp]
    \centering
    \includegraphics[width=0.8\linewidth]{images/cuda_cumtime.pdf}
    \caption{Cumulative time on text generation for LLMs. Unlike transformers, RWKV exhibits linear scaling.}
    \label{fig:inference_time}
\end{figure}

% This property saves RWKV from needing $n$ forward passes of $n_i^2$ each, which scales like $O(n^3)$, during autoregressive text generation, and it is replaced with just $n$ passes of a single (the current) token, which scales like $O(n)$. The improved computational efficiency is linked to faster text generation and can be further leveraged to lower the hardware requirements and/or to handle larger sequences.

\section{Future Work} \label{future_work}

There are several promising directions for future work on the RWKV architecture. Work can be done to increase model expressivity by enhancing the time-decay formulations and exploring initial model states while maintaining efficiency.

The RWKV computational efficiency can be further improved by applying a parallel scan in the $wkv_t$ step to reduce the computational cost to $O(B\log(T)d)$.

The mechanisms used in RWKV can be applied to encoder-decoder architectures, potentially replacing the cross-attention mechanism. This could be applicable in seq2seq or multimodal settings, thereby enhancing efficiency during both training and inference. 

RWKV's state (or \textit{context}) can be leveraged for interpretability, predictability in sequence data, and safety. Manipulating the hidden state could also guide behavior and allow greater customizability through prompt tuning.

The RWKV architecture is not perfect, and can be improved via many aspects, such as modifying the formulae or implementing larger internal states. Larger states can enhance the model's memory to previous context and improve performance over various tasks.

%\paragraph{Downstream Uses}
%The RWKV model has several promising applications to explore in downstream use cases, such as task-specific fine-tuning~\cite{albalak2023improving} and enhanced interaction with humans~\cite{ouyang2022training}. Recently popular parameter-efficient fine-tuning methods such as LoRA~\cite{hu2022lora} can be adapted for RWKV, and would be interesting to characterize the model's behavior under different quantization schemes. Furthermore, exploring the transferability of learned representations across different tasks, including benchmarks such as FETA~\cite{albalak-etal-2022-feta}, can also be a direction worth investigating.



\section{Conclusions} \label{conclusions}
% agreed by [@mgrella - Matteo Grella, @hypnopump - Eric Alcaide] pls reach out before changing
We introduced RWKV, a new approach to RNN models exploiting the potential of time-based mixing components. RWKV introduces several key strategies that allow it to capture locality and long-range dependencies while addressing limitations of current architectures by: (1) replacing the quadratic QK attention with a scalar formulation at linear cost, (2) reformulating recurrence and sequential inductive biases to enable efficient training parallelization and efficient inference, and (3) enhancing training dynamics using custom initializations. 

We benchmark the proposed architecture in a wide variety of NLP tasks and show comparable performance to SoTA with reduced cost. Further experiments on expressivity, interpretability, and scaling showcase the model capabilities and draw parallels in behavior between RWKV and other LLMs. 

RWKV opens a new route for scalable and efficient architectures to model complex relationships in sequential data. While many alternatives to Transformers have been proposed with similar claims, ours is the first to back up those claims with pretrained models with tens of billions of parameters.

\section{Limitations}\label{limitations}
While our proposed RWKV model has demonstrated promising results regarding training and memory efficiency during inference, some limitations should be acknowledged and addressed in future work.

% linear attention may suffer on tasks that require long context
First, the linear attention of RWKV leads to significant efficiency gains but still, it may also limit the model's performance on tasks that require recalling minutiae information over very long contexts. This is due to the funneling of information through a single vector representation over many time steps, compared with the full information maintained by the quadratic attention of standard Transformers. In other words, the model's recurrent architecture inherently limits its ability to ``look back'' at previous tokens, as opposed to traditional self-attention mechanisms.
While learned time decay helps prevent the loss of information, it is mechanistically limited compared to full self-attention. 

% concrete examples ane needed to demonstrate the sensitivity of RWKV to prompts. also, please provide examples of prompts that requires look-back.
Another limitation of this work is the increased importance of prompt engineering in comparison to standard Transformer models. The linear attention mechanism used in RWKV limits the information from the prompt that will be carried over to the model's continuation. As a result, carefully designed prompts may be even more crucial for the model to perform well on tasks.

The above RWKV property was confirmed by studies on prompt engineering presented in Appendix \ref{appendix:promptImportanceAndRWKVvsGPT}. By changing the order of the information pieces, we were even able to almost double the RWKV performance for some tasks. 

\section{Ethics Statement}

In this paper, we present a novel architecture for sequential data processing and prove its effectiveness by building a series of LLMs trained on publicly released pretraining data \cite{gao2020pile,biderman2022datasheet} and later fine-tuned on publicly available instructions \cite{alpaca, codealpaca, guanaco, gpt4all, sharegpt, Firefly, BELLE,belle2023exploring}.

As a novel architecture for sequential data, RWKV has the potential to improve sequence-based models across different applications ranging from natural language processing to biomedical data processing or climate modelling. Since the training code is released open source, RWKV contributes to the democratization of AI, levels the playing field, and empowers members of the Open Source community to inspect, study, and finetune RWKV in particular tasks. Moreover, it contributes to advancing the understanding of LLMs capabilities and limitations. A significant amount of work has been devoted to increasing the efficiency of RWKV training  so as to minimize its cost and promote accessibility.
% see for ref: https://discord.com/channels/729741769192767510/1103039376184852622/1110868583308873768 

As LLMs trained on public data, RWKV's lower inference cost compared to Transformer alternatives makes it more suitable for deployment in consumer and edge hardware, which is a step towards the democratization and distribution of LLMs to the general public, creating better privacy and ownership incentives. It also lowers the resource barrier to Chat assistants and text generation for small and/or underrepresented communities. PreTrained model weights for different sizes ranging from 0.1B to 14B parameters trained on multiple languages are released to increase ease of adoption and allow for the study of emergent phenomena. 

On the other hand, with lower resource barriers, the spreading of AI-generated text might become more prevalent. Current RWKV LLMs may exhibit and/or reproduce biases and potentially harmful content present in the data used for training. Nonetheless, mitigation and finetuning strategies discussed for other, large Transformer models should be applicable to RWKV as well. 

% Finally, although current RWKV models show sensitivity to prompt choices and potential adversarial prompts, they mostly underperform when the ordering of information does not favor the RNN architecture. 

% Our RWKV model was developed using "The Pile" \cite{gao2020pile}, an open-source dataset, reflecting our commitment to transparency in AI research. We put significant emphasis on fairness, potential bias, and the careful interpretation of experimental results to prevent overgeneralization. The importance of respecting privacy rights, particularly in scenarios involving sensitive data, is underscored in our approach. We actively mitigate against potential misuse and adverse effects. By adhering to these ethical considerations, we aim to uphold the highest standards of AI research and application, aligning with the values upheld by EMNLP.  

\section*{Acknowledgements}

We thank StabilityAI for the compute used to train our models and for technical support in development of RWKV. We also thank the members of the RWKV and EleutherAI Discord servers for their help and work on further extending the applicability of RWKV to different domains.

% Entries for the entire Anthology, followed by custom entries
\bibliography{custom}
\bibliographystyle{acl_natbib}

\clearpage
\onecolumn

\appendix

% This might be both easier to read and take up much less space. The current version is somehat obnoxious to look at IMO.
\section{Author Contributions}
All authors contributed to the drafting of this paper. Eric Alcaide and Quentin Anthony organized the paper and its experiments and were involved in all phases of the development process.

\paragraph{Model Design and Development} Bo Peng (lead), Matteo Grella, Xuzheng He, Haowen Hou, Jiaming Kong, Johan S. Wind

\paragraph{Model Training} Bo Peng

\paragraph{Scaling Laws Analysis} Stella Biderman, Bo Peng

\paragraph{Benchmark Evaluations} Stella Biderman (lead), Kranthi Kiran GV, Krishna Sri Ipsit Mantri, Atsushi Saito, Qihang Zhao, Peng Zhou, Rui-Jie Zhuåç

\paragraph{Long Context Experiments} Xingjian Du, Rui-Jie Zhu, Bolun Wang, Ruichong Zhang, Jian Zhu, Rui-Jie Zhu

\paragraph{Inference Speed Experiments} Samuel Arcadinho, Przemys\l{}aw Kazienko, Qinghua Zhou

\paragraph{Information Flow Experiments} Huanqi Cao, Michael Chung, Matteo Grella, Ferdinand Mom, Zhenyuan Zhang

\paragraph{Chat Experiments} Jan Kocoń (lead), Przemys\l{}aw Kazienko, Bart\l{}omiej Koptyra, Hayden Lau, Xiangru Tang, Stanis\l{}aw Wo\'{z}niak, Zhenyuan Zhang

\paragraph{Ethics and Broader Impacts} Stella Biderman, Guangyu Song

\clearpage
\section{Author Contributions}

\paragraph{Bo Peng} Original RWKV idea, original code, performance optimizations, original experiments, and trained RWKV models from 0.1B to 14B.

\paragraph{Eric Alcaide} Manuscript (initial draft sections \ref{introduction}, \ref{related_work}; sections \ref{rwkv_model}, \ref{future_work} and \ref{conclusions}; revision and proofreading; final version ). Figures (\ref{fig:blocks},  \ref{fig:arch}, \ref{rwkv_model}, \ref{fig:rwkv-rnn-arch}). Experiments section \ref{infer_experiments}. Appendices \ref{sec:initialization_formulas}, \ref{appendix:inference_results}. Contributions to Appendix \ref{sec:appendix-cases}.

% \paragraph{Quentin Anthony} Led writing the paper
\paragraph{Quentin Anthony}  Manuscript (organization, initial draft sections \ref{introduction}, \ref{related_work}, \ref{background}; revision and proofreading; final version). 

% \paragraph{Alon Albalak - University of California, Santa Barbara} Wrote the abstract, introduction, and limitations. Sporadically helped in other sections: fixing typos, improving formatting and proofreading.
\paragraph{Alon Albalak} Manuscript (abstract and sections \ref{introduction}, \ref{limitations}, and \ref{future_work}; proofreading and revision). 

% \paragraph{Samuel Arcadinho} Inference benchmark and added related work.
\paragraph{Samuel Arcadinho} Contributions to Figures \ref{fig:inference_time}, \ref{fig:total_inference_mem}, and \ref{fig:total_inference_time}. Contributions to Appendix \ref{appendix:inference_results}.

\paragraph{Stella Biderman} Performed the scaling laws analysis and evaluated competitor models on benchmark tasks.

% \paragraph{Huanqi Cao - Tsinghua University} Proofreading, advises about figure drawing, discovered the permanent channels, prototyped Figure~\ref{fig:visualizations} (a), and rewrite a version of section \ref{transformer_parallel} and \ref{rnn_inference}.
\paragraph{Huanqi Cao} Manuscript (contributions to \ref{transformer_parallel} and \ref{rnn_inference}; proofreading and revision). Experiments for Appendix \ref{sec-appen:vis}.

% \paragraph{Xin Cheng} Proofreading on Section~\ref{background} and Section~\ref{rwkv_model}, adding references and working on the cases section in the Appendix~\ref{appendix:cases}.
\paragraph{Xin Cheng} Manuscript (proofreading and revision). Contributions to Appendix \ref{sec:appendix-cases}, \ref{sec-appen:eval-detail}.

% \paragraph{Michael Chung} Proofreading. Added some notes about time-mixing.
\paragraph{Michael Chung} Manuscript (contributions to section \ref{sec-appen:vis}; proofreading and revision).

\paragraph{Xingjian Du} Evaluation on Long Range Arena Benchmark (TBD until 5.31).

%\paragraph{Matteo Grella} Wrote section \ref{grad-stability-layer-stacking} about gradient stability and layer stacking. Wrote section \ref{sec-appen:vis} about key mechanisms enabling information propagation and the (absence of) positional encoding. Wrote a section to introduce the state (or context) at a high level. Wrote the section \ref{conclusions}. Contributed on section \ref{future_work}. Revised section \ref{rwkv_model} on high level architecture including \ref{transformer_parallel},  \ref{rnn_inference} and detailed implementation \ref{detailed_implementation}. Contributed on section \ref{limitations}. Contributed on section \ref{introduction} emphasizing differences between self-attention and channel-directed attention.
\paragraph{Matteo Grella} Manuscript (sections \ref{sec:appendix-gradstab},  \ref{sec-appen:vis}, \ref{conclusions}; contributions to sections \ref{introduction}, \ref{future_work} and \ref{limitations}; proofreading and revision). Contributions to Appendix \ref{sec:appendix-rnn}.

% \paragraph{Kranthi Kiran GV} Wrote related works section (including FlashAttention, LazySoftmax, Linformer, Perceiver, S4, RMT, AFT), Evaluation section (including Table \ref{tab:commonsense_reasoning_results} and \ref{tab:opencloseqa}), FLOP count section (restored and rewrote), Background section (Transformers). Prepared and fixed formatting of several tables and equations.
\paragraph{Kranthi Kiran GV} Manuscript (sections \ref{related_work} and \ref{evaluations}; contributions to section \ref{background}; revision and proofreading).  Tables \ref{appendix:inference_results} and \ref{appendix:inference_results}. Appendix \ref{appendix:flop_count}.

% \paragraph{Xuzheng He} Wrote two paragraphs for contrasting existing works with RWKV (FlashAttention, etc. and MLPMixers, etc). Added appendix \ref{sec-appen:vis}. Revised section \ref{sec:background-attn}, \ref{transformer_parallel}, \ref{grad-stability-layer-stacking}, \ref{sec-appen:vis}, Appendix \ref{sec:appendix-gradstab}. Made cleaner versions of Figure \ref{fig:rnn_archs_background}, \ref{fig:rwkv-rnn-arch}. Proofreading and editing.
\paragraph{Xuzheng He} Manuscript (contributions to section \ref{background}; proofreading and revision). Contributions to Figure%s \ref{fig:rnn_archs_background}
\ref{fig:rwkv-rnn-arch}. Appendix \ref{sec-appen:vis}. Contributions to appendix \ref{sec:appendix-gradstab}.

% \paragraph{Haowen Hou}
%Wrote the text and equations for the "Scaling Laws" section \ref{scaling laws}, reploted Scaling Laws figure (currently Figure \ref{fig:scaling-laws}) for better resolution. Wrote appendix \ref{sec-appen:compute-optimal} for estimating the optimal allocation of parameters and training tokens and made figure \ref{fig:isoFLOP} to show the isoFLOP curves of RWKV model.
% Wrote appendix \ref{appendix_smallinit} for discussion and experiment of small init embedding. Made the figure \ref{fig:small_init_emb} to show the experiment results of small init embedding.
\paragraph{Haowen Hou} Figure \ref{fig:small_init_emb}. Appendix \ref{appendix_smallinit}.

\paragraph{Jiaju Lin} RWKV on LRA benchmarking

% \paragraph{Przemys\l{}aw Kazienko - Wroclaw University of Science and Technology} Section \ref{infer_experiments}. Proofreading the whole paper. Supervising the comparison of RWKV and ChatGPT/GPT-4 
\paragraph{Przemys\l{}aw Kazienko} Manuscript (proofreading and revision). Contributions to Section \ref{infer_experiments}, \ref{limitations}, and Appendix~\ref{appendix:promptImportanceAndRWKVvsGPT}.

% \paragraph{Jan Kocon - Wroclaw University of Science and Technology} Fully rewritten  Section \ref{introduction} (\emph{Introduction}), additionally including motivation, comparison with RNNs/CNNs, and paper contributions. Proofreading the whole paper. Supervising the comparison of RWKV and ChatGPT/GPT-4 (Appendix~\ref{appendix:promptImportanceAndRWKVvsGPT}).
\paragraph{Jan Kocon} Manuscript (Section \ref{introduction}; proofreading and revision). Contributions to Appendix~\ref{appendix:promptImportanceAndRWKVvsGPT}.

% \paragraph{Jiaming Kong} Wrote the appendix \ref{sec:appendix-gradstab} for proving the gradient stability in RWKV mathematically. Ran some experiments and prepared a few proof-of-concept cases. Revised section \ref{sec-appen:vis} and appendix \ref{sec:appendix-rnn}, \ref{sec:appendix-cases}.
\paragraph{Jiaming Kong} Manuscript (revision and proofreading). Appendix  \ref{sec:appendix-gradstab}.

% \paragraph{Bart\l{}omiej Koptyra - Wroclaw University of Science and Technology} Additional experiments on more tasks using Raven and comparison to chatGPT. Working on the Appendix \ref{appendix:promptImportanceAndRWKVvsGPT}.
\paragraph{Bart\l{}omiej Koptyra} Manuscript (revision and proofreading) Contributions to Appendix \ref{appendix:promptImportanceAndRWKVvsGPT}.

% \paragraph{Hayden Lau} Proofread and copy-edited the entire paper, corrected sentence structures and typos. Provided cases in Appendix~\ref{appendix:cases} and made small fixes in \ref{introduction}. Revised section \ref{future_work} and \ref{conclusions}.
\paragraph{Hayden Lau} Manuscript (contributions to section \ref{introduction} and \ref{limitations}; proofreading and revision). Contributions to Appendix \ref{sec:appendix-cases}.

% \paragraph{Krishna Sri Ipsit Mantri} Wrote scripts for figures 4 and 6 (Scaling Laws plots and Evaluation benchmark plots)
\paragraph{Krishna Sri Ipsit Mantri} Figure \ref{fig:eval_all}


%\paragraph{Ferdinand Mom} Contribute to section \ref{introduction}, \ref{related_work}, \ref{rnn_inference}, \ref{sec-appen:vis} and appendix \ref{sec:appendix-rnn}. Proofreading + editing
\paragraph{Ferdinand Mom} Manuscript (contributions to section \ref{introduction}, \ref{related_work}, \ref{rnn_inference}, \ref{sec-appen:vis}; proofreading and revision). Contributions to Appendix \ref{sec:appendix-rnn}.

% \paragraph{Atsushi Saito}   Wrote background section (part of RNNs and  Fig:1-a,b,c), evaluation section (RQs and the other paragraphs), appendix evaluation details section(part of datasets), modified related-work section.
\paragraph{Atsushi Saito}  Manuscript (sections \ref{background} and \ref{evaluations}; contributions to section \ref{related_work}). %Figures \ref{fig:rnns-cells} , \ref{fig:rnns-qrnn}, \ref{fig:rnns-rwkv}.
Contributions to Appendix \ref{sec-appen:eval-detail}

\paragraph{Guangyu Song} Manuscript (rewrote section \ref{rwkv_model}; final version). Initial draft Ethics Statement). 

% \paragraph{Xiangru Tang} Authored the related works and background sections of the paper. He polished the abstract and added two results figures related to Zero-Shot Performance. Performed various small adjustments to the manuscript such as resizes, reorders, and citation modifications for improved coherence and cohesion. Adjusted the paper format, helped with table creation, and resolved numerous LaTeX errors. Provided supplementary explanations and advised on figure drawing and intermittently assisted with other sections by fixing typos, improving formatting, and proofreading. Added references and conducted extensive proofreading for the entire paper, which involved correcting sentence structures and typos.
\paragraph{Xiangru Tang} Manuscript (sections \ref{related_work} and \ref{background}; contributions to abstract; revision and proofreading). Contributions to Appendix \ref{sec:appendix-cases}. 

\paragraph{Bolun Wang} Contributions to Tables \ref{tab:comparison}.

% \paragraph{Johan S. Wind} Wrote numerically stable and efficient CUDA kernels used for training RWKV models. Made formulas for parameter and FLOP counts.
\paragraph{Johan S. Wind} RWKV performance optimizations (CUDA), Contributions to Appendix \ref{appendix:flop_count}. 


% \paragraph{Stanis\l{}aw Wo\'{z}niak - Wroclaw University of Science and Technology} Experiments with optimizing prompts and comparison to GPT models. Working on the Appendix \ref{appendix:promptImportanceAndRWKVvsGPT}.
\paragraph{Stanis\l{}aw Wo\'{z}niak} Contributions to Appendix \ref{appendix:promptImportanceAndRWKVvsGPT}.

% \paragraph{Ruichong Zhang - Tsinghua University} Proofreading and typo corrections; Advices on Figure \ref{fig:ctxlen_rwkv_loss}.
\paragraph{Ruichong Zhang} Manuscript (proofreading and revision); Contributions to Figure \ref{fig:ctxlen_rwkv_loss} and Appendix \ref{sec:appendix-cases}.

% \paragraph{Zhenyuan Zhang} Made Fig. \ref{fig:arch}. Wrote the recursive $WKV$ formula. Formatted all equations to be aligned and fix punctuation. Did the information propagation experiment in Appendix \ref{sec-appen:vis}. Provided the multi-round example case.
\paragraph{Zhenyuan Zhang} Manuscript (revision and proofreading). Figure \ref{fig:arch}. Experiments Appendix \ref{sec-appen:vis}. Contributions to Appendices \ref{sec:appendix-rnn} and \ref{sec:appendix-cases}.

% \paragraph{Qihang Zhao} Adjust the paper format, help wrote Tab. \ref{tab:enwik8}, and plan to add some schematic figures. Provide supplementary explanations on the experimental configuration.
\paragraph{Qihang Zhao} Manuscript (proofreading and revision). Contributions to Table \ref{tab:enwik8}.

% \paragraph{Peng Zhou} Calculate the computational complexity on Table \ref{tab:enwik8}.
\paragraph{Peng Zhou} Contributions to Tables \ref{tab:comparison} and Table \ref{tab:enwik8}.

\paragraph{Qinghua Zhou} Manuscript (Proofreading and revision of section \ref{rwkv_model}; Add missing citations in \ref{rnn_inference}). Revision of Figures \ref{fig:blocks} and \ref{fig:eval_all}.

% \paragraph{Jian Zhu - University of British Columbia} Wrote and polished the Related Works section (adding references and condensing the whole section). Remade Figure \ref{fig:arch} and \ref{fig:ctxlen_rwkv_loss} into (most of) its current form. Proofreading and various small fixes across the manuscript. 
\paragraph{Jian Zhu} Manuscript (section \ref{related_work}; proofreading and revision). Figures \ref{fig:arch} and \ref{fig:ctxlen_rwkv_loss}.

% \paragraph{Rui-Jie Zhu} Added the Table \ref{tab:enwik8} and run the experiments on Enwik8. Also added Table 1 for illustrating the complexity of the RWKV.
\paragraph{Rui-Jie Zhu} Tables \ref{tab:comparison} and \ref{tab:enwik8}. Experiments for table \ref{tab:enwik8}. 

\section{Additional Related Work} \label{related_work}
%TODO: cite https://arxiv.org/pdf/2302.13939.pdf, 
Recently, a number of techniques have been proposed to address the limitations of transformers. 
\paragraph{Optimizing Attention Mechanism} % just a draft name
%FlashAttention \citep{dao2022flashattention} proposes an IO-aware attention mechanism which uses tiling to reduce the number of memory reads/writes, enabling faster training with increased context length. %\citet{rabe2022selfattention} achieves constant memory for a single query and logarithmic ($O(\log n)$) memory for self-attention by computing the attention values in chunks and using the Lazy Softmax \citep{jang2019mnnfast}, which delays the division in Softmax until the end of attention operation.
Many transformer variants (``x-formers'') have been introduced to reduce the complexity of transformers \cite{tay2022efficient}, including sparse attention \cite{Beltagy2020Longformer,reformer,guo2022longt5}, approximating the full attention matrix \cite{wang2020linformer,ma2021luna,performer}, combining chunked attention with gating \cite{MEGA2023} and other efficient methods \cite{katharopoulos2020transformers,jaegle2021perceiver}. 

Some recent works like FlashAttention \cite{dao2022flashattention} and others \cite{rabe2022selfattention, jang2019mnnfast} share similarities with RWKV's chunked computation scheme. %, as they all iterate over the decomposed denominator and numerator of the Softmax, and only perform the division at the end. 
Despite being memory-efficient, their time complexity remains quadratic or contains chunk size as a hidden factor. In contrast, RWKV achieves better space and time complexity during inference by formulating a linear attention as an RNN.

% However, despite the memory efficiency, their time complexity remains quadratic. In contrast, by reformulating a linear attention mechanism as RNN, our work achieves both the best space and time complexity during inference. % I don't think the O(Td) time complexity and O(d) space complexity can be improved further.

%Recently, equipped with an exponential moving average which transmits position-aware dependencies between chronological  adjacent two chunks, 

%Contrary to MEGA \cite{MEGA2023}, which results in a linear complexity for a fixed sequence chunk size $c$ and has a drawback of missing contextual information from  other chunks, RWKV achieves linear complexity with a quadratic-attention free approach and controls contextual information via \textit{receptance} parameters.

% that incorporates the inductive bias of position-aware local dependencies into the position-agnostic attention mechanisms. For fixed chunk size $c$ and an arbitrary sequence length $T$, resulting in a linear complexity of $O(kc^{2})$ with respect to $T$, where $T=kc$. To avoid a drawback of missing contextual information from  other chunks, RWKV achieves linear complexity with a quadratic-attention-free approach and controls contextual information via \textit{receptance} parameters.

%\paragraph{Altering the Attention} % just a draft name
%Some approaches modify the attention mechanism while preserving most of its characteristics \cite{tay2022efficient}. Linformer \cite{wang2020linformer} achieved linear ($O(n)$) complexity in both time and space by approximating the self-attention matrix as low-rank matrix through linear projections.

%The Perceiver \cite{jaegle2021perceiver} leverages an asymmetric attention mechanism and iterative cross-attention modules to scale the model efficiently to large sequences. Linear Transformer \cite{katharopoulos2020transformers} reformulates the self-attention mechanism as a linear dot-product of kernel feature maps, revealing an inherent relationship between RNNs and Transformers while demonstrating an acceleration in autoregressive inference.

\paragraph{Attention Free Models} % just a draft name
Another line of research replaces the attention mechanism with other modules to scale to long sequences. MLP-Mixer and others \cite{DBLP:journals/corr/abs-2105-01601, liu2021pay} propose replacing attention by Multi-Layer Perceptrons (MLPs) in computer vision tasks. The Attention Free Transformer (AFT) \citep{zhai2021attention} and HrrFormer \citep{alam2023recasting} replaces dot-product self-attention with a computationally efficient alternative. None of these models have been successfully scaled to the point where drawing comparisons with transformer-based large language models makes sense.

There has also been substantial research into state space models (SSM) \citep{gu2021efficiently} and its variants \citep{dao2022hungry,gupta2022diagonal,poli2023hyena}. In contrast to the preceding models, SSM and its successors have shown substantial progress towards efficient scaling. Simultaneously with this work, \citet{poli2023hyena} train SSM-based models with 125 million and 355 million parameters and show that the performance is on-par with a transformer that uses a mix of local and global attention \citep{black2021gpt}.

%Recurrent neural networks (RNNs) \cite{RNN_LSTM97, RNN_GRU2014} are models capable of selectively transmitting information across sequence steps while processing sequential data one element at a time.  While RNNs accomplished various progresses for sequential modeling tasks such as language modeling \cite{RNNLM2010, RNNLM2011, RNNLM2015}, speech recognition \cite{RNNASR2013, RNNASR2015}, and machine translation \cite{RNNMT_ACL2015, RNNMT_ICLR2015, RNNMT2014, RNNMT_NIPS2014}, inherent difficulty to parallelize training remains. 
\paragraph{Advances in RNNs}

Inspired by the success of transformers, RNN-style \cite{RNN_LSTM97, RNN_GRU2014} recursive components have also been modified to increase context length, such as the Recurrent Memory Transformer \cite{bulatov2022recurrent,bulatov2023scaling} and Linear Recurrent Units \cite{orvieto2023resurrecting}.
Most similar to our work, the Quasi-Recurrent neural network (QRNN)  \cite{RNN_QuasiRNN17} uses both convolutional layers and recurrent pooling functions across timesteps and channels. While QRNN utilizes convolutional filters with fixed sizes, RWKV employs a time-mixing module as an attention mechanism with time-decaying factors. Different from the element-wise pooling in QRNN, RWKV includes a parametrized channel-mixing module  that is parallelizable.


\section{Time-Mixing Block as an RNN Cell}
%\section{Time-Mixing and Channel-Mixing Blocks as RNN Cells}
\label{sec:appendix-rnn}
As stated in \ref{rnn_inference}, the RWKV time-mixing block can be formulated as an RNN, as the $WKV$ computation can be written in such a recursive form:
\begin{align}
    a_0, b_0 &= 0, \\
    \label{eq:wkv-recur}
    wkv_t &= \frac{a_{t-1} + e^{u + k_t} \odot v_t}{b_{t-1} + e^{u + k_t}}, \\
    a_t &= e^{-w} \odot a_{t-1} + e^{k_t} \odot v_t, \label{eq:statea}\\
    b_t &= e^{-w} \odot b_{t-1} + e^{k_t}. \label{eq:stateb}
\end{align}

\begin{wrapfigure}{r}{0.5\textwidth}
\centering
    \includegraphics[width=0.4\textwidth]{images/rwkv_as_rnn.pdf}
    \caption{RWKV time-mixing block formulated as an RNN cell. Color codes: yellow ($\mu$) denotes the token shift, red (1) denotes the denominator, blue (2) denotes the numerator, and pink (3) denotes the fraction computations in \ref{eq:nom-denom}. $h$ denotes the numerator-denominator tuple.}
    \label{fig:rwkv-rnn-arch}
\end{wrapfigure}


The dataflow of the RNN-like time-mixing is shown in Fig. \ref{fig:rwkv-rnn-arch}, where the hidden states $h$ is the numerator-denominator tuple $(a, b)$. To avoid overflow in calculating $e^{k_t}$, a numerical trick is used in the official implementation. Noticing that $a_1 = e^{k_1} \odot v_1$ and $b_1 = e^{k_1}$, we set $a'_1 = v_1, b'_1 = 1, p_1 = k_1$, where $p_t$ stores the shared exponents of $a_t$ and $b_t$. Now the above recursion can be converted into a numerical safe version, for each time step $t > 1$:
\begin{align}
    q &:= \max(p_{t-1}, u+k_t), \\
    % a_t^* &= e^{p_{t-1}-q} \odot a'_{t-1} + e^{u+k_{t} - q} \odot v_t, \\
    % b_t^* &= e^{p_{t-1}-q} \odot b'_{t-1} + e^{u+k_{t} - q},\\
    wkv_t &=\frac{e^{p_{t-1}-q} \odot a'_{t-1} + e^{u+k_{t} - q} \odot v_t}{e^{p_{t-1}-q} \odot b'_{t-1} + e^{u+k_{t} - q}}.
\end{align}
The update to $a'_t, b'_t$, and their shared exponent is also carried out in a similar fashion:
\begin{align}
    q' &:= \max(p_{t-1} - w, k_t), \\
    a'_{t} &= e^{p_{t-1} - w - q'} \odot a'_{t-1} + e^{k_t - q'} \odot v_t, \label{eq:stateaa}\\
    b'_{t} &= e^{p_{t-1} - w - q'} \odot b'_{t-1} + e^{k_t - q'}, \label{eq:statebb}\\
    p_t &= q'.\label{eq:statepp}
\end{align}
% TODO: add channel-mixing as RNN cell here
The RWKV model has an internal state that stores some previous information. In each layer, the internal state consists five parts, each of which is a vector with $D$ numbers, where $D$ is the model dimension. The five parts are:
\begin{itemize}
    \item The current input of the Time-mix block $x_t$;
    \item The current input of the Channel-mix block $y_t$;
    \item The numerator of the $WKV$ value $a'_t$, as defined in equation \eqref{eq:stateaa};
    \item The denominator of the $WKV$ value $b'_t$, as defined in equation \eqref{eq:statebb};
    \item An auxiliary state $p_t$ in \eqref{eq:statepp}, which is used for $WKV$ computation to maintain numerical precision.
\end{itemize}
Which yields a total size of $5DL$ parameters. It is worth noting that in an algebraic context with infinite precision, the helper state $p_t$ can be ignored, and the $WKV$ numerator and denominator can be computed directly using equations \eqref{eq:statea} and \eqref{eq:stateb}, reducing the size of the internal state to $4DL$.

\section{Parameter initializations}
\label{sec:initialization_formulas}

We describe the specific parameter initializations below and motivate the design choices. Parameters belonging to residual blocks are often adjusted by layer depth and total number of layers. Let $\#$ denote the vocabulary size, $s$ denote the embedding dimension, $d$ denote the hidden size (we use $d=4s$), $L$ the number of layers, $l$ the layer index (from 0 to $L-1$), we use the following initializations:
\begin{itemize}
    \item Embeddings are initialized to $\mathcal{U}$ ($\pm$\num{1e-4}) as explained in \ref{smallinit}
    \item For the time-mixing blocks (\ref{eq:time-mix1}, \ref{eq:time-mix2}, \ref{eq:time-mix3}), initializations are
    $\mu_{k_i} = (\frac{i}{s})^{1-\frac{l}{L}}$,
    $\mu_{v_i} = (\frac{i}{s})^{1-\frac{l}{L}} + \frac{0.3l}{L-1}$ and $\mu_{r_i} = \frac{1}{2} \cdot (\frac{i}{s})^{1-\frac{l}{L}}$
    \item For the channel-mixing blocks (\ref{eq:channel_mix_eqs1}, \ref{eq:channel_mix_eqs2}), $\mu_{k_i}$ and $\mu_{r_i}$ are initialized to $(\frac{i}{s})^{1-\frac{l}{L}}$
    \item $w_i$ (\ref{eq:nom-denom}), also known as ``time decay'', is initialized to $-5 + 8 \cdot (\frac{i}{d-1})^{0.7+\frac{1.3l}{L-1}}$. Intuitively, it is the discount factor applied to previous tokens over time.  
    % for u: https://github.com/BlinkDL/RWKV-LM/blob/main/RWKV-v4/src/model.py#L184
        \item $u_i$ (\ref{eq:nom-denom}), also known as ``bonus'', is set to $0.5 \cdot (((i+1) \mod 3) - 1) + \log0.3$. It is the special weighting applied to the current token in equation \ref{eq:nom-denom}. The alternating zigzag pattern initially creates subtle variations in the tensor elements, which are intended to help the model treat different dimensions of the embedding distinctively.
    \item $W_o$ (\ref{eq:time-mix4}) (time-mixing) and $W_v$ (channel-mixing) are initialized to $\mathcal{N}(0, \sqrt{\frac{d}{s}}=2)$ % should be W_k here?
    \item All other $W_r, W_k, W_v$ weights are initialized to 0 so the model can start learning from the beginning without noisy signals. 
    \item All LayerNorm weights start from 1 and biases from 0.
\end{itemize}

\section{Small Init Embedding}  \label{appendix_smallinit}

This section presents the experimental validation of small initialization embedding. The experimental setup is as follows. In the baseline configuration, the parameters are initialized using a normal distribution with a mean of 0.0 and a standard deviation of 0.02, which is a commonly used initialization method in models like BERT and GPT. On the other hand, in the small initialization of the embedding (small init emb) experiment, the parameters are initialized using a uniform distribution with a range of 1e-4, which is slightly different from RWKV where a normal distribution with a standard deviation of 1e-4 is used. However, this difference is negligible and does not affect our conclusions. The experiments were conducted with a batch size of 400. As depicted in Figure \ref{fig:small_init_emb}, the loss curve for the small init emb exhibits a faster rate of decrease and convergence compared to the traditional initialization using a normal distribution.
\begin{figure}[htp]\label{small_init_emb}
    \centering
    \includegraphics[height=0.5\linewidth]{images/small_init_emb.pdf}    
    \caption{Effect of small initialization embedding.}
    \label{fig:small_init_emb}
\end{figure}

\section{Hyperparameters}\label{appendix_hyperparam}

\begin{table*}[t]
\centering
\begin{tabular}{lcccccc}\toprule
Model              & 169M    & 430M    & 1.5B    & 3B      & 7B      & 14B      \\ \midrule
Init LR         & 0.0006  & 0.0004  & 0.0003  & 0.00015 & 0.00015 & 0.0001   \\
Warmup Mini-Epochs & 361     & 411     & 443     & 451     & 465     & 544      \\
End LR          & 0.00001 & 0.00001 & 0.00001 & 0.00001 & 0.00001 & 0.000007\\ \bottomrule
\end{tabular}
\caption{Hyperparameters for our learning rate (LR) schedule of the pretrained models.}
\label{tab:hyperparam}
\end{table*}

To train the models mentioned, we use $\epsilon=(0.9,0.99)$ without weight decay for the Adam optimizer, and switch batch size dynamically between 128 or 256 sequences, each of 1024 tokens.
We further organize the training into multiple mini-epochs, each of 40320 samples, to guide our learning rate schedule.
The training process takes 8043 mini-epochs to make one pass over the Pile.
The initial warming up mini-epochs have a constant learning rate of ``Init LR''.
After the warming up mini-epochs, the learning rate exponentially decays until in the last mini-epoch, in which the model finishes training on the entire Pile, the learning rate arrives at the ``End LR''.
The related hyperparameters are shown in Table~\ref{tab:hyperparam}.

\section{Gradient Stability in RWKV} 
\label{sec:appendix-gradstab}

In this section, we present a mathematical description of the gradient stability property in RWKV, focusing specifically on the time-mixing block. By gradient stability we mean that if the inputs $x_t$ are bounded and the model parameters are fixed, then the gradients with respect to $W_k$ and $W_v$ are uniformly bounded for all $T$ (thus not exploding). Consequently, we can control the amount each $x_t$ contributes to the gradient at $T$ in a naturally decaying fashion by the weight decay mechanism $w$ (thus not vanishing unless desired).
% Specifically we want to address how the novel design of RWKV's layers alleviated the vanishing gradient problem.

First, we make the simplification that there are no token shifts, this will not affect the final conclusion. In this scenario, $wkv_T$ can be written as
\begin{align} wkv_T = \frac{\sum _{t=1}^T K^e_t \odot v_t} {\sum _{t=1}^T K^e_t} = \mathrm{E} (v_t) = \frac{\mathrm{S}(v_t)}{\mathrm{S}(1)},
\end{align}
where
$$ v_t = W_v x_t, \quad\frac{\partial (v_t)_i}{\partial(W_v)_{i,j}} = (x_t)_j, $$
$$ K^e_t = e^{W_k x_t+w_{T,t}}, \quad\frac{\partial (K^e_t)_i}{\partial (W_k)_{i,j}} = (x_t)_j (K^e_t)_i, $$
and $\mathrm{S}(\cdot)$ and $\mathrm{E}(\cdot)$ are shorthand for denoting sums and averages over weights $K^e_t$.

The loss function at position $T$ can be written as
\begin{align}
L_T =l(f(wkv_T), y_T).
\end{align}
Because $wkv_T$ relates to $(W_k)_{i,j}$ and $(W_v)_{i,j}$ only through the $i$-th channel $(wkv_T)_i$, we have
\begin{align} \frac{\partial L_T}{\partial (W_v)_{i,j}} = \frac{\partial L_T}{\partial (wkv_T)_i} \frac{\partial (wkv_T)_i}{\partial (W_v)_{i,j}}. \end{align}
The first part of the above equation contains trivial operations like output layers, and other layers of time-mixing, which can be proven inductively.
The second part of the above equation can be bounded as
\begin{align}
\left|\frac{\partial (wkv_T)_i}{\partial (W_v)_{i,j}}\right| &= \left| \frac{\partial \mathrm{E}_i[(v_t)_i]}{\partial (W_v)_{i,j}} \right| \nonumber\\
&= |\mathrm{E}_{i}[(x_t)_j]| \leq \max_t |(x_t)_j|,
\end{align}
which is irrelevant to $T$. Similarly,
\begin{align}
\frac{\partial (wkv_T)_i}{\partial (W_k)_{i,j}} &= \partial \frac{\mathrm{S}_i[(v_t)_i]}{\mathrm{S}_i(1)}/ \partial(W_k)_{i,j} \nonumber \\
&= \frac{\mathrm{S}_i[(x_t)_j (v_t)_i]}{\mathrm{S}_i(1)} - \frac{\mathrm{S}_i[(x_t)_j]\mathrm{S}_i[(v_t)_i]}{\mathrm{S}_i(1)^2} \nonumber\\
&= \mathrm{E}_{i}[(x_t)_j(v_t)_i] - \mathrm{E}_{i}[(x_t)_j]\mathrm{E}_{i}[(v_t)_i] \nonumber\\
&= \mathrm{cov}_i((x_t)_j, (v_t)_i)
\end{align}
can also be bounded. Note that $wkv$'s softmax operation contains at least two non-zero terms ($u$ and $w$), so the above ``covariance'' will not degenerate into 0.
% Take the partial derivative of loss $L_T$ at time T %(cross entropy between \(O_T\) and \(\text{label}_T\))
% to \(W_v\) from one time-mixing layer,
% as expressed in formula \ref{eq:pd_to_wk}. %Note that \(O_T\) only relate to \(W_v\) through \(wkv_{T}\),
% Let $\alpha$ be the normalizing coefficient introduced by the LayerNorm layer before time-mixing:

% \begin{align}
% \notag
% \frac{\partial L}{\partial W_v} &= \alpha\left(\frac{\partial L}{\partial wkv_T}\sum_{t=1}^T
% \frac{\partial {wkv_T}}{\partial v_t}\frac{\partial {v_t}}{\partial W_v}\right) \\
% \label{eq:pd_to_wk}
% &\propto \left(\sum_{t=1}^T  \frac{\partial {wkv_T}}{\partial v_t} \frac{\partial v_t}{\partial W_v}\right).
% \end{align}

% In the context of vanishing gradient problem for recurrent network, the risk of RWKV\textquotesingle s design would be that for some relatively large $T$ values and small $t$ values, $\frac{\partial {wkv_T}}{\partial v_t}$ \emph{could be near zero}. And the model would not learn anything useful inside the whole context window. What we want to prove here is that under mild assumptions, and the use of LayerNorms, this will not be the case. All \(J(i)\) for any $i$ will be bounded at some threshold.

% Now we can simplify the above expression, ignoring $\alpha$ and $\frac{\partial L}{\partial wkv_T} = w_0 \sigma(r_T)$ which are considered external coefficients:
% \begin{align}
% \label{eq:contri_to_grad}
% \frac{\partial {wkv_T}}{\partial v_t} &= \frac{e^{-(T-1-t)w + k_t}}{{\sum_{i=1}^{t-1} e^{-(t-1-w)w+k_i} + e^{u+k_t}}},\\
% \frac{\partial v_t}{\partial W_v} &=\mu_v x_t + (1-\mu_v) x_{t-1}.
% \end{align}
% % Note that equation \ref{eq:contri_to_grad} is very similar to that in \cite{zhai2021attention}.
% The contributions of gradient from token $x_t$ or $v_t$ to a parameter $W_v$, are mostly controlled by a small decaying factor $w$, and the learned key \(k_t\). The Softmax operation guarded any $k_t$ from total control of the final gradient. Under some mild assumptions about the initial embedding, we can safely assume all tokens in the context window are accounted for. Also, it is trivial to see that the total gradient of \(W_v\) in one update will also be bounded:

% \begin{align}
% \notag
% \sum_{t=1}^T \frac{\partial {wkv_T}}{\partial v_t} &=\sum_{t=1}^T \frac{\partial wkv_T}{\partial v_t}\\
% \notag
% &\leq \frac{\sum_{t=1}^T e^{-(T-1-t)w + k_t}}{\sum_{i=1}^{t-1} e^{-(t-1-w)w+k_i} + e^{u+k_t}}\\
% &\leq 1.
% \end{align}

% Likewise, the partial derivative of $L$ w.r.t. $W_k$ can be derived as:

% \begin{align}
% \notag
% \frac{\partial L}{\partial W_k} &= \alpha\left(\frac{\partial L}{\partial wkv_T}\sum_{t=1}^T
% \frac{\partial {wkv_T}}{\partial k_t}\frac{\partial {k_t}}{\partial W_k}\right)\\
% \notag
% &\propto \frac{e^{u+W_k X_K^T}X_K^T \sum(ev) W_v}{(e^{u+W_k X_K^T} + \sum(e))^2}\\
% &\propto \frac{\lambda \sum(ev)}{(\lambda + \sum(e))^2},
% \end{align}
% where
% \begin{align}
%     \sum(ev) &= \sum_{i=1}^{t-1}e^{(1+i-t)w+k_i} v_i, \\
%     \sum(e) &= \sum_{i=1}^{t-1}e^{(1+i-t)w+k_i}, \\
%     X_K^T &= \mu_k x_T + (1-\mu_k)x_{T-1}, \\
%     X_V^T &= \mu_v x_T + (1-\mu_v)x_{T-1}, \\
%     \lambda &= e^{u+W_k X_K^T}X_K^T.
% \end{align}
% We observe that the contribution of each token $x_i$ to the gradient of $W_k$ is also controlled in a softmax weighted manner and learned from training data.

% The partial gradients of $L$ w.r.t. $W_k, W_v$ in channel-mixing layers have very similar behavior except for the squared ReLU activation function. The stability of these gradients can be shown using similar techniques.

\section{Model Behavior Visualization}\label{sec-appen:vis} % just a draft name

\begin{wrapfigure}{r}{0.5\textwidth}
\centering
    \includegraphics[width=0.45\textwidth]{images/visualize-decay.pdf}
    \caption{Model behavior visualizations of RWKV.}
    \label{fig:visualizations-decay}
\end{wrapfigure}

The right plot illustrates the time decays ($e^{-w}$) in each layer of the RWKV-169M model, sorted along the channel axis. Notably, several decays in the last layers are very close or equal to one, implying that certain information is preserved and propagated throughout the model's temporal context. Meanwhile, many decays in the initial layer are close to zero, which corresponds to local operations in $wkv$ (\ref{eq:nom-denom}), likely to be associated with tasks such as text parsing or lexical analysis.
(Note that the local operations in $wkv$ are due to the extra parameter $u$, when $e^{-w}$ is degenerated into 0.) These patterns of time decays are partly learned, but also come from parameter initialization as it speeds up training.

% From @cryscan https://discord.com/channels/992359628979568762/1085545872533758112/1104973719455158322
The plot below shows the information retrieval and propagation path in the RWKV-430M model. The experiment follows the \emph{causal trace} method introduced by \citet{meng2022locating}, where we
\begin{enumerate}
    \item Run the model once, and record all states and activation of each layer during the computation;
    \item Corrupt the input embeddings of the subject using noise (``The Eiffel Tower'' in this example);
    \item Restore the states and activation of a certain layer at a certain token during the computation, and record the log-probability of the model outputting the correct answer (``Paris'').
\end{enumerate}
Unlike transformers, RWKV relies on the recursive propagation of information in the time dimension.
In this case, the fact that the Eiffel Tower is located in Paris is retrieved in layer 4 just after the model sees ``The Eiffel''. It is then passed down to the subsequent layers. In layer 20, mostly, the information is propagated through time until reaching where it is needed. Finally, at the token ``of'', it is passed down to the last layer for outputting the answer.

\begin{figure}[htbp]
    \centering
    \includegraphics[width=0.8\linewidth]{images/visualize-information-propagation.pdf}
%     \includegraphics[width=\linewidth, height=0.8\linewidth]{images/loss_params.png}
    \caption{Model behavior visualizations of the RWKV model.}
    \label{fig:visualizations}
\end{figure}

\section{Additional Evaluations}\label{sec-appen:eval-detail}

\subsection{Further details on NLP tasks}

We evaluate on the following tasks:

\paragraph{ARC ~\cite{ARC2018}} A dataset designed for multiple-choice question answering, encompassing science exam questions ranging from third grade to ninth grade. It has Easy and Challenge subsets that we report results on separately.
\paragraph{BoolQ ~\cite{clark2019boolq}} A binary yes/no question answering benchmark.
\paragraph{COPA ~\cite{COPA2011}} A dataset to evaluate achievement in open-domain commonsense causal reasoning.
\paragraph{HeadQA ~\cite{HeadQA2020}} A benchmark consisting of graduate-level questions encompassing various fields such as medicine, nursing, biology, chemistry, psychology, and pharmacology.
\paragraph{HellaSwag} ~\cite{HellaSwag2019} A novel benchmark for commonsense Natural Language Inference (NLI) which is build by adversarial filtering against transformer models. 
\paragraph{LAMBADA~\cite{LAMBADAdataset}} A benchmark dataset that evaluates the model's contextual reasoning and language comprehension abilities by presenting context-target pairs, where the objective is to predict the most probable target token. We follow standard practice and use the untokenized version created by OpenAI \citep{brown2020language}.
\paragraph{OpenBookQA ~\cite{OpenBookQA2018}} A QA dataset to evaluate human comprehension of a subject by incorporating open book facts, scientific knowledge, and perceptual common sense, drawing inspiration from open book exams.
\paragraph{PIQA~\cite{Bisk2020}} A benchmark for the task of physical common sense reasoning, which consists of a binary choice task that can be better understood as a set of two pairs, namely (Goal, Solution).% A Benchmark Dataset for the Task of Physical Common Sense Reasoning. PiQA is a binary choice task that can be better understood as a set of two pairs, namely (Goal, Solution).
    % such that a sequence of discriminators is utilized to choose a difficult set of incorrect answers generated during the process.
\paragraph{ReCoRD ~\cite{ReCord}} A benchmark for evaluating commonsense reasoning in reading comprehension by generating queries from CNN/Daily Mail news articles and requiring text span answers from corresponding summarizing passages.
\paragraph{SciQ ~\cite{SciQ2017}} A multiple-choice QA dataset which was created using an innovative approach to gather well-crafted multiple-choice questions that are focused on a specific domain. 
\paragraph{Winogrande ~\cite{Wino2020}} A dataset designed to evaluate the acquisition of common sense reasoning by neural language models, aiming to determine whether we are accurately assessing the true capabilities of machine common sense.
    %\item StoryCloze~\cite{StoryCloze2016} A benchmark to present a novel approach to assess comprehension of narratives, narrative generation, and script acquisition, focusing on commonsense reasoning.
    %\item TriviaQA ~\cite{TriviaQA2017} A QA-IR dataset which is constituted of triples of questions, answers, supporting evidence, and independently collected evidence documents, with an average of six documents per question for reliable sources.
    %\item MMMLU ~\cite{MMMLU2021} A multi-task dataset for 57 tasks containing elementary mathematics, US history, computer science, law, etc. 

\begin{figure*}[!ht]
\centering
\subfloat[ARC (Challenge)]{%
\includegraphics[width=0.32\textwidth]{images/acc_arc_challenge.pdf}
}
\subfloat[ARC (Easy)]{%
\includegraphics[width=0.32\textwidth]{images/acc_arc_easy.pdf}
}
\subfloat[BoolQ]{%
\includegraphics[width=0.32\textwidth]{images/acc_boolq.pdf}
}\\
\subfloat[COPA]{%
\includegraphics[width=0.32\textwidth]{images/acc_copa.pdf}
}
\subfloat[HeadQA]{%
\includegraphics[width=0.32\textwidth]{images/acc_headqa.pdf}
}
\subfloat[HellaSwag]{%
\includegraphics[width=0.32\textwidth]{images/acc_hellaswag.pdf}
}\\
\subfloat[LAMBADA (OpenAI)]{%
\includegraphics[width=0.32\textwidth]{images/acc_lambada_openai.pdf}
}
\subfloat[OpenBookQA]{%
\includegraphics[width=0.32\textwidth]{images/acc_openbookqa.pdf}
}
\subfloat[PiQA]{%
\includegraphics[width=0.32\textwidth]{images/acc_piqa.pdf}
}\\
\subfloat[ReCoRD]{%
\includegraphics[width=0.32\textwidth]{images/acc_record.pdf}
}
\subfloat[SciQ]{%
\includegraphics[width=0.32\textwidth]{images/acc_sciq.pdf}
}
\subfloat[Winogrande]{%
\includegraphics[width=0.32\textwidth]{images/acc_winogrande.pdf}
}
\caption{Zero-Shot Performance of RWKV on common language modeling evaluation benchmarks.} 
\label{fig:eval_all}
\end{figure*}
\clearpage


\subsection{Evaluation on Long Range Arena}\label{lorfa}
The Long-Range Arena (LRA) benchmark \citep{tay2021long} is designed to assess the performance of models in handling lengthy context situations. It includes a collection of tasks with sequences ranging from 1,000 to 16,000 tokens, covering various types of data like text, natural language, synthetic images, and mathematical expressions. We apply RWKV on the LRA benchmark and the report results are in Table~\ref{tab:lra}. Other models' performances are directly cited from \citet{gu2022efficiently,alam2023recasting}.

\begin{table}[ht!]
  \small
  \caption{ Evaluation on Long Range Arena. Other models reported in the literature \citep{gu2022efficiently,alam2023recasting}. \textbf{Bolded} values are the best.
  }
    \centering
    \begin{tabular}{@{}llllllll@{}}
        \toprule
        \textsc{Model}        & \textsc{ListOps}  & \textsc{Text}     & \textsc{Retrieval} & \textsc{Image}    & \textsc{Pathfinder} & \textsc{Path-X} & \textsc{Avg}      \\
        \midrule
        Transformer           & 36.37             & 64.27             & 57.46              & 42.44             & 71.40               & \xmark          & 53.66             \\
        Reformer              & 37.27 & 56.10             & 53.40              & 38.07             & 68.50               & \xmark          & 50.56             \\
        BigBird               & 36.05             & 64.02             & 59.29              & 40.83             & 74.87               & \xmark          & 54.17             \\
        Linear Trans.         & 16.13             & 65.90 & 53.09              & 42.34             & 75.30               & \xmark          & 50.46             \\
        Performer             & 18.01             & 65.40             & 53.82              & 42.77             & 77.05               & \xmark          & 51.18             \\
        \midrule
        FNet                  & 35.33             & 65.11             & 59.61              & 38.67             & 77.80   & \xmark          & 54.42             \\
        Nystr{\"o}mformer     & 37.15             & 65.52             & 79.56  & 41.58             & 70.94               & \xmark          & 57.46             \\
        Luna-256              & 37.25             & 64.57             & 79.29              & 47.38 & 77.72               & \xmark          & 59.37 \\
        Hrrformer             & 39.98             & 65.38             & 76.15              & 50.45                       & 72.17 & \xmark & 60.83 \\        \midrule
        S4                    & \textbf{59.60}    & \textbf{86.82}    & \textbf{90.90}     & \textbf{88.65}    & \textbf{94.20}      & \textbf{96.35}  & \textbf{86.09}    \\
        RWKV                  & 55.88    & 86.04    & 88.34     & 70.53    & 58.42       & \xmark  & 72.07    \\
        \bottomrule
    \end{tabular}
    \label{tab:lra}
\end{table}

The results show that RWKV performs second only to the S4 model in five datasets. While RWKV substantially underpreforms S4 on Image, Pathfinder, and Path-X, on the problems related to natural language and computer code processing RWKV performs on par with S4 or nearly so.

\begin{comment}
    \begin{table*}[!]
\centering
\small
\begin{tabular}{llllllllll}
        \toprule
\textbf{Model} & \textbf{Params} & \textbf{PIQA} & \textbf{StoryCloze} & \textbf{HellaSwag} & \textbf{WinoGrande} & \textbf{ARC-e} & \textbf{ARC-c} & \textbf{OBQA} \\
 & B & acc & acc & acc\_norm & acc & acc & acc\_norm & acc\_norm \\
\midrule
RWKV &  0.17 & \TBF{65.07} & \TBF{58.79} & \TBF{32.26} & 50.83 & \TBF{47.47} & \TBF{24.15} & \TBF{29.60} \\
Pythia & 0.16      & 62.68 & 58.47 & 31.63 & \TBF{52.01} & 45.12 & 23.81 & 29.20 \\
GPT-Neo & 0.16     & 63.06 & 58.26 & 30.42 & 50.43 & 43.73 & 23.12 & 26.20 \\
\midrule
RWKV &  0.43 & \TBF{67.52} & \TBF{63.87} & \TBF{40.90} & 51.14       & \TBF{52.86} & 25.17       & \TBF{32.40} \\
Pythia & 0.40      & 66.70       & 62.64      & 39.10      & \TBF{53.35} & 50.38       & \TBF{25.77} & 30.00 \\
GPT-Neo & 0.40     & 65.07       & 61.04       & 37.64      & 51.14     & 48.91       & 25.34 & 30.60 \\
\midrule
RWKV & 1.5 & \TBF{72.36} & \TBF{68.73} & \TBF{52.48}   & 54.62      & \TBF{60.48} & \TBF{29.44}             & \TBF{34.00} \\
Pythia & 1.4 & 71.11       & 67.66       & 50.82        & \TBF{56.51} & 57.74 & 28.58             & 30.80 \\
GPT-Neo & 1.4 & 71.16      & 67.72        & 48.94       & 54.93       & 56.19       & 25.85             & 33.60 \\
\midrule
RWKV & 3.0 & \TBF{74.16} & \TBF{70.71} & \TBF{59.89} & 59.59       & \TBF{65.19}         & \TBF{33.11} & \TBF{37.00} \\
Pythia & 2.8 & 73.83       & \TBF{70.71} & 59.46       & \TBF{61.25} & 62.84   & 32.25 & 35.20 \\
GPT-Neo & 2.8 & 72.14      & 69.54       & 55.82        & 57.62      & 61.07         & 30.20       & 33.20 \\
\midrule
RWKV & 7.4 & \TBF{76.06} & 73.44       & 65.51      & 61.01       & \TBF{67.80} & \TBF{37.46}         & \TBF{40.20} \\
Pythia & 6.9 & 74.54       & 72.96       & 63.92      & 61.01       & 66.79       & 35.07               & 38.00 \\
GPT-J & 6.1  & 75.41        & \TBF{74.02 }& \TBF{66.25}& \TBF{64.09} & 66.92 & 36.60                    & 38.20 \\
\midrule
RWKV & 14.2    & \TBF{77.48}            & \TBF{76.06}   & \TBF{70.65}       & 63.85       & 70.24       & \TBF{38.99}       & \TBF{41.80} \\
GPT-level$^*$ & 14.2 & 76.49           & 74.97   & 68.72       & \TBF{65.14}      & \TBF{70.77}        & 37.99       & 39.27 \\
\midrule
Pythia (c.f.)  & 11.8    & 75.90           & 74.40   & 67.38       & 64.72       & 69.82        & 36.77       & 38.80 \\
GPT-NeoX (c.f.)  & 20.6  & 77.69     & 76.11 & 71.42 & 65.98    & 72.69 & 40.44 & 40.20 \\
\bottomrule
\end{tabular}
\centering
\caption{\label{tab:commonsense_reasoning_results}
Zero-Shot Performance of the model on Common Sense Reasoning Tasks. $^*$ Interpolation of Pythia and GPT-Neo models
}
\end{table*}


\begin{table*}[htp!]
\centering
\small
\begin{tabular}{lllllllll}
\toprule
\textbf{Model} & \textbf{Params} & \textbf{LAMBADA} & \textbf{LAMBADA} & \textbf{headQA} & \textbf{sciq} & \textbf{triviaQA} & \textbf{ReCoRD} & \textbf{COPA} \\
  & B & ppl & acc & acc\_norm & acc & acc & em & acc \\
\midrule
RWKV  & 0.17 & \TBF{29.33} & 32.99 & 25.78 & \TBF{77.50} & 1.26 & 62.03 & \TBF{66.00} \\
Pythia  & 0.16 & 24.38 & \TBF{38.97} & \TBF{25.82} & 76.50 & \TBF{1.31} & \TBF{66.32} & 62.00 \\
GPT-Neo & 0.16 & 30.27 & 37.36 & 25.16 & 76.60 & 1.18 & 64.92 & 64.00 \\
\midrule
RWKV  & 0.43 & 13.04 & 45.16 & \TBF{27.32} & 80.30 & \TBF{2.35} & 70.48 & 65.00 \\
Pythia  & 0.40  & \TBF{11.58} & \TBF{50.44} & 25.09 & \TBF{81.50} & 2.03 & \TBF{75.05} & \TBF{67.00} \\
GPT-Neo & 0.40  & 13.88 & 47.29 & 26.00 & 81.10 & 1.38 & 73.79 & 65.00 \\
\midrule
RWKV  & 1.5  & 7.04  & 56.43 & \TBF{27.64} & 85.00 & \TBF{5.65} & 76.97 & \TBF{77.00} \\
Pythia  & 1.4  & \TBF{6.58}  & \TBF{60.43} & 27.02 & \TBF{85.50 }& 5.52 & \TBF{81.43} & 73.00 \\
GPT-Neo & 1.4  & 7.5   & 57.25 & 27.86 & 86.00 & 5.24 & 80.62 & 69.00 \\
\midrule
RWKV  & 3.0    & 5.25  & 63.96 & 28.45 & 86.50 & \TBF{11.68} & 80.87 & \TBF{82.00} \\
Pythia  & 2.8  & \TBF{4.93}  & \TBF{65.36} & \TBF{28.96} & \TBF{87.70} & 9.63  & 85.10 & 77.00 \\
GPT-Neo & 2.8  & 5.63  & 62.22 & 27.17 & 89.30 & 4.82  & \TBF{83.80} & 80.00 \\
\midrule
RWKV  & 7.4  & 4.38  & 67.18 & \TBF{31.22} & 88.80 & \TBF{18.30} & 83.68 & \TBF{85.00} \\
Pythia  & 6.9  & 4.3   & \TBF{67.98} & 28.59 & 90.00 & 15.42 & 86.44 & \TBF{85.00} \\
GPT-J   & 6.1  & \TBF{4.1}   & 68.31 & 28.67 & \TBF{91.50} & 16.74 & \TBF{87.71} & 83.00 \\
\midrule
RWKV  & 14.2 & \TBF{3.86}  & 70.83 & \TBF{32.64} & 90.40 & \TBF{24.58} & 85.67 & \TBF{85.00} \\
GPT-level$^*$ & 14.2 & 3.81 & \TBF{70.94} & 31.03 & \TBF{92.20} & 22.37 & 87.89 & 82.66 \\
\midrule
Pythia (c.f.)  & 11.8 & 3.89  & 70.44 & 30.74 & 91.80 & 20.57 & 87.58 & 82.00 \\
GPT-NeoX (c.f.)  & 20.6 & 3.64 & 71.94 & 31.62 & 93.00 & 25.99 & 88.52 & 84.00 \\
\bottomrule
\end{tabular}
\caption{\label{tab:opencloseqa}
Zero-Shot Performance of various models on different tasks. $^*$ Interpolation of Pythia and GPT-Neo models
}
\end{table*}
\end{comment}

\subsection{Enwik8 Perplexity}\label{app:enwik8}

We also evaluate our model in terms of perplexity on the Enwik8 dataset. Baseline comparisons are made with Reformer \citep{reformer}, Synthesizer \citep{synthesizer} (the best performing dense version), Linear Transformer \citep{katharopoulos2020transformers}, Performer \citep{performer}. $L, d,$ and $T$ denote the number of blocks (network depth), dimension of features, and sequence length, respectively. Both Linear Transformer and Performer are implemented with customized CUDA kernels (github.com/idiap/fast-transformers), and all other models are implemented in native Pytorch. $^1$ No weight decay nor dropout was used. $^2$ Trained with AdamW and weight decay set to 0.1, dropout of 0.1,  batch size of 16, and initial learning rate of 6e-4.

\begin{table*}[htbp]
\small
    \centering
    \begin{tabular}{llllllllllllll}
 \toprule
 Method & L& d  &T &Train bpc & Test bpc & Time Complexity & Space Complexity\\
 \midrule
 Transformer&  12& 512& 1024&0.977& 1.137&$O(T^2d)$&$O(T^2 + Td)$ \\
 Transformer&  24& 256& 1024&1.039& 1.130&$O(T^2d)$&$O(T^2 + Td)$ \\
    \midrule
 Reformer& 12 & 512& 1024& 1.040 &  1.195& $O(T\log{T}  d)$&$O(T\log{T} + Td)$ \\
 Synthesizer&  12& 512& 1024& 0.994 & 1.298& $O(T^2d)$&$O(T^2 + Td)$  \\
 Linear Transformer& 12& 512& 1024&0.981& 1.207& $O(T d^2)$&$O(Td + d^2)$ \\
 Performer&  12& 512& 1024&1.002& 1.199& $O(T d^2\log{d})$&$O(Td\log{d} + d^2\log{d})$  \\
  AFT-simple &  12& 512& 1024&1.046& 1.209& $O(Td)$&$O(Td)$  \\
 \midrule
RWKV-RNN$^1$ &  6 & 512  & 1024 & 0.720 & -  & $O(\mathbf{Td})$&$O(\mathbf{d})$ \\
RWKV-RNN$^2$ &  12 & 512  & 1024 & 1.010 & 1.178  & $O(\mathbf{Td})$&$O(\mathbf{d})$ \\
  \bottomrule    
    \end{tabular}
        \caption{Enwik8 results, measured in bits per character (bpc).}
\label{tab:enwik8}
\end{table*}{}

\begin{comment}
\section{Estimating the optimal allocation of parameters and training tokens}\label{sec-appen:compute-optimal}

In our study, we investigate the optimal model size and number of tokens for training a RWKV language model while considering a fixed compute budget. We systematically adjust the model size while keeping the training FLOP count constant across six different values. We select the number of training tokens for different model sizes to ensure a constant final FLOP count. The range of training FLOP counts spans from 1 PF-day to 300 PF-days. By following \cite{hoffmann2022training}, we plot the IsoFLOPs curves (Figure \ref{fig:isoFLOP}). In Figure \ref{fig:isoFLOP}, we observe a distinct minimum point in the loss curve, indicating the existence of a compute-optimal RWKV model to be trained within a given FLOP budget. On the curve, points located to the left of the minimum value indicate a scenario where there is abundant data but the model size is too small, suggesting the possibility of increasing the model size. Conversely, points located to the right of the minimum value indicate a situation where the model is too large compared to the available data, suggesting the possibility of reducing the model size.

We utilize a parabolic curve fitting method to determine the model size that achieves the minimum loss for each IsoFLOPs curve (Figure \ref{fig:isoFLOP}). Subsequently, we establish a power law relationship between the amount of compute ($C$) and the optimal model size ($N_{opt}$) as well as the optimal number of training tokens ($D_{opt}$). The exponents of the power law, denoted as $N_{opt} \propto C^a$ and $D_{opt} \propto C^b$, are estimated to be $a=0.519$ and $b=0.481$, respectively. Specifically, our findings suggest that with each increase in compute, a proportional increase in both data size and model size yields optimal results. For example, if compute is increased by a factor of 10, the model size and data size should be increased by approximately 3.1 times. Similarly, if compute is increased by a factor of 100, the model size and data size should be increased by approximately 10 times. This indicates a consistent relationship between compute, model size, and data size in achieving optimal performance.

\begin{figure}[htbp]
    \centering
    \includegraphics[width=\linewidth]{images/iso_curve.png}
    \caption{IsoFLOP curves for RWKV model. Each dashed line represents a fixed compute budget.}
    \label{fig:isoFLOP}
\end{figure}
\end{comment}

\section{Inference results} \label{appendix:inference_results}

Figures \ref{fig:total_inference_mem} and \ref{fig:total_inference_time} illustrate, respectively, the results on time (s) and memory (RAM, VRAM) requirements for LLM inference in \textit{float32} precision. We benchmark the following model families and sizes: 
\begin{itemize}
\item \textbf{RWKV}: 169m, 430m, 1.4b, 3b, 7b, 14b
\item \textbf{Bloom} \cite{scao2022bloom}: 560m, 1b, 3b 
\item \textbf{OPT} \cite{zhang2022opt}: 125m, 350m, 1.3b, 2.7b, 6.7b, 13b
\item \textbf{GPT-Neo} \cite{black2021gpt}: 125m, 1.3b, 2.7b
\item \textbf{Pythia} \cite{biderman2023pythia}: 160m, 410m, 1.4b, 2.8b, 6.7b, 12b
\end{itemize}

%Missing models are due to Out Of Memory (OOM) errors. A comparison at 512 tokens is shown in Figure \ref{fig:total_inference_time} as some large transformer models produced an OOM when inferencing longer sequences. For GPU experiments, we use an NVIDIA A100 with 80GB of VRAM. For CPU experiments, we use an AMD EPYC processor with 30 CPU cores and 200 GiB RAM. 

%\label{appendix:inference_results}
% TODO: limit to 512 tokens as some models didnt reach 1024 due to OOM.
\begin{figure*}[ht] % btp]
    \centering
    \includegraphics[width=0.45\linewidth]{images/inference_memory_cpu.pdf}
    \includegraphics[width=0.45\linewidth]{images/inference_memory_cuda.pdf}
    \caption{Text generation inference memory (CPU RAM, GPU VRAM) for LLMs. Model parameters are not accounted.}
    \label{fig:total_inference_mem}
\end{figure*}

% On which CPU and GPU? 
\begin{figure*}[hbtp] % btp]
    \centering
    \includegraphics[width=0.45\linewidth]{images/cpu_time_inference.pdf}
    \includegraphics[width=0.45\linewidth]{images/cuda_time_inference.pdf}
    \caption{Text generation inference time for LLMs.}
    \label{fig:total_inference_time}
\end{figure*}
\clearpage
\section{Importance of prompt construction and comparison to GPT models}
\label{appendix:promptImportanceAndRWKVvsGPT}

Inspired by \citet{kocon2023chatgpt}, we compared the zero-shot performance of the RWKV-4-Raven-14B with ChatGPT (access in February 2023) and GPT-4 using several known NLP tasks, i.e., recognizing textual entailment (RTE), Winograd Natural Language Inference (WNLI), and recognizing emotions elicited in readers (GoEmotions and PolEmo2). 
% aggression and sarcasm detection, classification of offensive, unhealthy texts, Q\&A for math, and sentiment analysis. 
Each model got the same prompts manually chosen to receive proper responses from the ChatGPT model. As shown in Tab.~\ref{tab:RWKVvsChaGPTvsGPT4}, RWKV performs significantly worse than ChatGPT and GPT-4 in several specific tasks. We suspect that this disparity is likely caused by the choice of prompts used to generate the answers since the prompts are written in natural language and do not take into account that RWKV, as an RNN, is unable to look back inside an instruction. 

\begin{table}
\centering
\begin{tabular}{llccccc}
\toprule
  Task Name      & Measure &  ChatGPT & GPT-4    & RWKV-GPT & RWKV-adapted & SOTA\\
\midrule
   RTE  & F1 Macro &  88.1 & \textbf{91.3} & 44.2 & 74.8 & 92.1\\
   WNLI  & Accuracy &  81.7 & \textbf{91.6} & 47.9 & 49.3 & 97.9\\
   GoEmotions  & F1 Macro & \textbf{25.6} & 23.1 & 7.9 & 7.9 & 52.8\\
   PolEmo2 & F1 Macro & \textbf{44.1} & 41.0 & 38.2 & 40.9 & 76.4\\
\bottomrule 
\end{tabular}
\caption{ChatGPT, GPT-4 and RWKV-4-Raven-14B reasoning performance comparison in RTE \cite{SuperGLUE}, WNLI \cite{wang-etal-2018-glue}, GoEmotions \cite{GoEmotions}, and PolEmo2 \cite{kocon2019multi} benchmarks. RWKV GPT prompts were primarily used for ChatGPT in \cite{kocon2023chatgpt}. SOTA is provided as a supplementary reference.}
\label{tab:RWKVvsChaGPTvsGPT4}
\end{table}


When the instruction style was adapted (re-ordered) to respect that RNNs are not capable of "retrospective processing", the quality may significantly change, e.g., for RTE \cite{SuperGLUE} F1 Macro increased from 44.2\% to 74.8\%. We hypothesize that RWKV models are more sensitive to the position of the components in the context, as RNN-based architectures cannot look back and readjust the weight of previous information. For better performance, the desired information should be placed \textit{after} the main question. 

\begin{tcolorbox}[colback=white,colframe=black!75!black,title=An example ChatGPT prompt for recognizing textual entailment (RTE)]
    Having premise <here is a premise> judge if the following hypothesis <here is a hypothesis> is logically connected with the premise. Answer "entailment" if yes, or "not\_entailment" if no.
\end{tcolorbox}

\begin{tcolorbox}[colback=white,colframe=black!75!black,title=A re-ordered RWKV prompt for RTE taking into account the nature of the RNN]
    Can you tell me if the hypothesis is entailment or is not entailment to the premise?\\premise: <here is a premise>\\hypothesis: <here is a hypothesis>
\end{tcolorbox}

% An example of ChatGPT prompt for recognizing textual entailment (RTE): \\
% \textit{Having premise <here is a premise> judge if the following hypothesis <here is a hypothesis> are logically connected with the premise? Answer "entailment" if yes, or "not\_entailment" if no.}\\

% An adjusted RWKV prompt taking into account the nature of the RNN: \\
% \textit{Can you tell me if the hypothesis is entailment or is not entailment to the premise?\\premise: <here is a premise>\\hypothesis: <here is a hypothesis>}\\

\begin{table}
\centering
\begin{tabular}{llrrrrr}
\toprule
  Task Name      & Measure &  ChatGPT   & RWKV-adapted & SOTA\\\midrule
   Aggression & F1 Macro & \textbf{69.10} & 56.66 & 74.45 \\
   MathQA & Accuracy & \textbf{71.40} & 5.43 & 83.20\\
   Sarcasm & F1 Macro & 49.88 & \textbf{50.96} & 53.57\\
   TweetSent & F1 Macro & \textbf{63.32} & 52.50 & 72.07\\
   Unhealthy & F1 Macro & \textbf{45.21} & 43.30 & 50.96\\
\bottomrule 
\end{tabular}
\caption{ChatGPT and RWKV-4-Raven-14B performance comparison in Aggresion  \cite{wulczyn2017}, Sarcasm \cite{Siddiqui2019sarcasmania}, Unhealthy \cite{Unhealthy2020}, MathQA \cite{cobbe2021training}, and TweetSent \cite{tweeteval} benchmarks. SOTA is provided as a supplementary reference.}
\label{tab:RWKVvsChaGPT}
\end{table}


While separating the instruction from the input is relatively easy to do, some other aspects of prompt engineering are harder to quantify. For that purpose, we also tested the approach of stating the input after the question on multiple other tasks, i.e., aggression and sarcasm detection, classification of unhealthy (offensive) texts, mathematical Q\&A, and sentiment analysis, see Tab.~\ref{tab:RWKVvsChaGPT}. The results suggest that better prompts might reduce the disparity between models. Raven achieves comparable results to ChatGPT on unhealthy conversation detection and even surpasses it on the sarcasm detection dataset. While such an approach to prompting looks necessary, it is not enough in itself to replace the capability of having free access to the whole context. Therefore, prompt engineering seems to be significantly more important for the RNN models rather than for standard transformers. It is entirely possible that good prompts to RNN models do not mean additional restrictions, but should simply be constructed using completely different guidelines. The authors of the aforementioned paper \cite{kocon2023chatgpt}\footnote{This is in line with the idea discussed in \cite{wei2022chain}} perform chain-of-thought to improve results on the MathQA dataset. Even including this approach, the Raven model achieved a very low accuracy of 5.43\%. Without it, the model performed even worse, performing only very basic and simple calculations and achieving 4.13\% accuracy. Raven struggled with questions that required intermediate results. It is likely that the order of information presented in the math questions inside the dataset poses a challenge for the RWKV model. It is yet to be seen if prompt engineering can address this issue. This further emphasizes the importance of the order of information the model receives. \\

\begin{tcolorbox}[colback=white,colframe=black!75!black,title=Template used to prompt the Raven model in MathQA with chain-of-thought]
    Write the reasoning and highlight the answer to the question at the end in the format: "Answer: ". The question is: <here is a question>
\end{tcolorbox}

\begin{tcolorbox}[colback=white,colframe=black!75!black,title=Template used to prompt the Raven model in MathQA without chain-of-thought]
    Write the answer to the math question in the format: "Answer: ".\newline The question is: <here is a question>
\end{tcolorbox}

% The template used to prompt the Raven model in MathQA with chain-of-thought prompting: \\
% \textit{Write the reasoning and highlight the answer to the question at the end in format: 'Answer: '. The question is: <here is a question>}\\

% The template used to prompt the Raven model in MathQA without chain-of-thought prompting: \\
% \textit{Write the answer to the math question in the format: 'Answer: '.\newline The question is: <here is a question>}


% It shows that P-tuning can be very effective for RWKV because we can directly tune the full state. We are going to do this in follow-up papers. 
% % In addition dataset PolEmo2 is dataset in Polish, so even that it is not in English it has comparable results to GPT-4. It means that RWKV to some extent is able to understand instructions in different language. 

\section{Cases} \label{sec:appendix-cases}
In this part, we present a few instances of outputs produced by RWKV model using a Chat interface\footnote{ttps://github.com/BlinkDL/ChatRWKV}\footnote{https://huggingface.co/spaces/BlinkDL/ChatRWKV-gradio}. 

% % here is the template, research lead should be responsible for maintaining this.
% \begin{tcolorbox}[colback=white,colframe=black!75!black,title=RWKV-4-Raven-14B]
% \textbf{Prompt}\\

% Output

% \begin{tcolorbox}
% \begin{lstlisting}[language=Python]
% code
% \end{lstlisting}
% \end{tcolorbox}
% \end{tcolorbox}


% TODO: @hynpopump: add better outputs and provide model sizes
\begin{tcolorbox}[colback=white,colframe=black!75!black,title=RWKV-4-Raven-14B]
\textbf{Tell me about ravens.}\\

Ravens are large, black birds with a distinctive white head and neck. They are found in most parts of the world, including North America, Europe, Asia, and Australia. Ravens are known for their intelligence and problem-solving abilities. They are also considered to be symbols of death and transformation in many cultures. Ravens are often associated with the afterlife or death because they have been known to eat carrion or even other birds. In some cultures, ravens are also believed to be messengers of the gods or guardians of treasure.\\
\end{tcolorbox}


\begin{tcolorbox}[colback=white,colframe=black!75!black,title=RWKV-4-Raven-14B]
\textbf{Explain the concept of machine learning to a 10-year-old.}\\

Machine learning is like a super-smart computer that can learn on its own. It can recognize patterns and make decisions based on what it has learned. Just like how you might remember something from the past, the computer can remember what it's learned and use it to make predictions or decisions in the future. It's like having a smart friend who always knows what to do!\\
\end{tcolorbox}


\begin{tcolorbox}[colback=white,colframe=black!75!black,title=RWKV-4-Raven-14B]
\textbf{Write a short poem describing nature's beauty in springtime.}\\

\begin{center}
Nature's beauty in springtime is rare,\\
A symphony of colors, a gentle surprise.\\
The flowers bloom, the birds sing and sing,\\
As the sun shines down, warming our skin.\\
The trees sway in the gentle breeze,\\
As we take in their beauty with ease.\\
The sky is blue, the clouds are white,\\
A sight to behold, it's truly divine.\\
\end{center}
\end{tcolorbox}


\begin{tcolorbox}[colback=white,colframe=black!75!black,title=RWKV-4-Raven-14B]
\textbf{Below is the numerical stable formula for updating the state variable in RWKV:}\\
\begin{equation*}
    \begin{aligned}
    a_t^* &= e^{p-q} a_{t-1} + e^{u+k_{t} - q} v_t\\
    b_t^* &= e^{p-q} b_{t-1} + e^{u+k_{t} - q}\\
    wkv_t &=\frac{a_t^*}{b_t^*}
    \end{aligned}
\end{equation*}

\textbf{Please write out the corresponding pytorch code to do so, where both $a_t$ and $b_t$ are torch.Tensor([n])}\\


Here is the PyTorch code to update the state variable in RWKV using the given formula:
\begin{tcolorbox}
\begin{lstlisting}[language=Python]
import torch
def update_state(a, b, u, v, p, q):
    a_t = torch.exp(p - q) * a + torch.exp(u + k - q) * v
    b_t = torch.exp(p - q) * b + torch.exp(u + k - q)
    return a_t, b_t
a, b = torch.randn(2, 3)
u = torch.randn(2)
v = torch.randn(2)
wkv = torch.randn(2)
a_t, b_t = update_state(a, b, u, v, wkv, torch.randn(2))
print("a_t:", a_t)
print("b_t:", b_t)
\end{lstlisting}
\end{tcolorbox}
This code defines a function update\_state that takes in the state variables a and b, the control variables u and v, and the parameters wkv and torch.randn. It then uses the PyTorch torch.exp function to compute the exponential of the control variables and the log of the normalization factor. Finally, it uses these values to update the state.\\
\end{tcolorbox}
 
\begin{tcolorbox}[colback=white,colframe=black!75!black,title=RWKV-4-Raven-14B,float=h]
\textbf{Explain the following metaphor: Life is like cats.}\\

The metaphor "Life is like cats" suggests that life can be 
unpredictable and full of unexpected twists and turns, just like a cat's behavior. It implies that we cannot fully control or predict the outcome of our lives, just as we cannot control the movements of a cat. The metaphor also implies that life is complex and multifaceted, with many different variables and factors at play. Just as a cat can have multiple personalities or moods, life can have many twists and turns that are difficult to predict or understand. Overall, the metaphor suggests that life is complex and unpredictable, much like a cat's behavior.\\
\end{tcolorbox}

\begin{tcolorbox}[colback=white,colframe=black!75!black,title=RWKV-4-Raven-14B,float=ht]
\textbf{Test the convergence of $\sum_{i=0}^\infty 2^{-i}$.}\\

The given series is a geometric series with the first term $a_0 = 2$ and common ratio $r = \frac{1}{2}$. We can find the sum of this series using the formula for a geometric series:\\
$$\sum_{i=0}^\infty a_i = \frac{a_0}{1-r} = \frac{\frac{1}{2}}{1-\frac{1}{2}} = \frac{1}{1-\frac{1}{2}} = \frac{2}{1} = 2$$\\
So, the sum of the given series is $\boxed{2}$.\\
\end{tcolorbox}

\begin{tcolorbox}[colback=white,colframe=black!75!black,title=RWKV-4-Raven-14B]
\textbf{I would like to cook some Chinese food at home. Do you have any suggestions on the meal I could choose?}\\

Yes, of course! If you are not very experienced in cooking Chinese food, you could start with making noodles or rice, then make a dish using one of those. I also think making stir-fry with some meat is an excellent idea.\\

\textbf{I want to make something more special. What about red braised pork belly?}\\

Yes, red braised pork belly is also very popular and tasty.\\

\textbf{Tell me the steps to cook.}\\

You will need a few ingredients such as pork belly, star anise, sugar, ginger, scallions and wine. Put all of these in a pot, cook on medium heat and simmer until the meat is tender. After that, remove the meat from the pot and reserve the sauce. Slice the meat and add some ginger slices into it, then boil until it's almost done.\\

\textbf{What do I need to be careful with the most during the cooking process?}\\

To cook pork belly properly, it is very important to monitor the heat of the stove and be patient with the cooking process. Overcooking it can cause the pork to be tough, dry and difficult to chew. You can add a few ingredients like soy sauce or black beans to give the pork more flavor.\\

\textbf{Tell me more about the proper heat and timing.}\\

If the temperature of the water is too high, it can damage the collagen and connective tissues, resulting in tough pork belly. It is essential to control the temperature carefully, especially when braising the pork. Additionally, keep the water simmering but not boiling to maintain a low simmer, as high heat will dry out the pork belly and result in tough and unpalatable meat. Lastly, monitor the heat during cooking and make sure the heat is medium, as this is what is most appropriate for slow cooking pork belly.\\

\textbf{I see. Thank you for your advice.}\\

You're welcome! It was a pleasure chatting with you. Let me know if you need any further assistance in the future.\\

\end{tcolorbox}


\begin{tcolorbox}[colback=white,colframe=black!75!black,title=RWKV-4-Raven-14B,float=h]%"Logic" Questions%
\textbf{If all apples are fruits and all fruits are sweet, are all apples sweet? Please explain your answer in detail.}\\

Yes, all apples are fruits, and all fruits are sweet. Therefore, it is safe to assume that all apples are sweet.\\

\textbf{If birds can fly. Penguins are birds. Can penguins fly? Please explain your answer in detail.}\\

No, penguins cannot fly. Penguins are birds, but they are not capable of flight like other birds. They have wings and can glide for short distances, but they cannot take off or fly long distances like other birds.\\

\textbf{The day before two days after the day before tomorrow is Saturday. What day is it today? Please explain your answer in detail.}\\

Today is Friday. Two days after tomorrow is Sunday. Tomorrow is Saturday.\\

\textbf{If five cats can catch five mice in five minutes, how long will it take one cat to catch one mouse? Please explain your answer in detail.}\\

If five cats can catch five mice in five minutes, it means that each cat can catch one mouse in one minute.
To find out how long it will take one cat to catch one mouse, we can use the concept of proportionality.\\ Let's assume that it takes x minutes for a cat to catch a mouse. Then, the number of mice caught by a cat in one minute is equal to the number of cats multiplied by the number of mice caught by each cat:
x = 5 * 1 = 5\\
So, it will take 5 minutes for one cat to catch one mouse.\\
\end{tcolorbox}


\begin{tcolorbox}[colback=white,colframe=black!75!black,title=RWKV-4-Raven-14B]
\textbf{Summarize the following paragraph.}\\

\textbf{<The Abstract of this paper goes here>}\\

The paper presents a novel architecture, Receptance Weighted Key Value (RWKV), that combines the efficient parallelizable training of Transformers with the efficient inference of RNNs. The model is able to scale to tens of billions of parameters and exhibits linear computational complexity during training and inference, making it a promising alternative to Transformers for sequence processing tasks.\\

\textbf{Please further summarize it in one sentence.}\\

The researchers have developed a new architecture that combines the efficient parallelizable training of Transformers with the efficient inference of RNNs, which allows for linear scaling in memory and computational requirements.\\

\end{tcolorbox}



\end{document}


% ============================
% OLD INTRO & MOTIVATION (COPY FROM 2023.05.18)
\section{Introduction} \label{introduction}
\textcolor{blue}{Introduction is a placeholder, Alon will go back and re-do the introduction once we have more concrete details of what is being done in this work.}
% What are we trying to do, and why is this relevant?
% Develop efficient sequential NN architecture (efficient in training and inference)
Since the introduction of the transformer by Vaswani et al.~\cite{NIPS2017_3f5ee243}, its success has been attributed to the self-attention mechanism. This powerful and flexible method allows models to capture long-range dependencies and complex interactions between input elements. The inherent parallelizability of self-attention combined with ever-increasing dataset sizes, and its ability to create rich contextual representations, has led to state-of-the-art models across a multitude of benchmarks and applications.
% Transformers have become dominant architectural paradigm, but have quadratic computational and space complexity due to self-attention mechanism
However, despite the transformative impact of the Transformer, it is not without its shortcomings. One notable drawback is the quadratic complexity of its self-attention mechanism. It renders the architecture computationally expensive for tasks involving long input sequences and in situations with restricted memory budgets. These limitations have spurred a wealth of research aiming to improve the scaling properties of the Transformer, often at the expense of sacrificing some of the properties that make it so effective~\citep{wang2020linformer, zaheer2020bigbird, dao2022flashattention}.


% Quadratic complexity limits the use cases (on-device, conversations, etc.)
% SEE: 4.3 RNN-like Formulation for current explanation of advantages

% Why is developing efficient NN architectures hard?
Addressing the issue of quadratic scaling has proven to be a challenging endeavor, as it requires navigating the intricate balance between the computational efficiency of the self-attention mechanism and the expressive capacity of the model. Furthermore, the non-local nature of the self-attention mechanism poses significant difficulty in devising more efficient alternatives that do not compromise the rich expressiveness of the Transformer.

% How do we solve it?
In this paper, we propose a novel neural network architecture, dubbed the RWKV network, which serves as a variant of both the RNN and the Transformer. The RWKV network is designed to alleviate the quadratic scaling issue and replace it with a linear one while retaining the crucial properties that have made the Transformer the dominant architecture in the field
% Go more into depth on what these crucial properties are
. By leveraging the strengths of both RNNs and Transformers, the RWKV network presents a compelling alternative for tackling a wide range of tasks that demand both computational efficiency and expressiveness. In doing so, the time and space complexity is sharply decreased. Our approach is very friendly to large model size and computation resources. See Table~\ref{tab:comparison} for the complexity comparison of RWKV and related approaches. 

\begin{table}
\centering
\setlength{\tabcolsep}{2pt}
\small
\begin{tabular}{l|c|c}
\toprule
\textbf{Model} & \textbf{Time} & \textbf{Space}\\
\midrule
Transformer & $O(T^2d)$ & $O(T^2 + Td)$ \\
Reformer & $O(T\log{T}  d)$ & $O(T\log{T} + Td)$ \\
% Synthesizer & $O(T^2d)$ & $O(T^2 + Td)$ \\
Linear Transformers & $O(T d^2)$ & $O(Td + d^2)$ \\
Performer & $O(T d^2\log{d})$ & $O(Td\log{d} + d^2\log{d})$ \\
AFT-full & $O(T^2d)$ & $O(Td)$ \\
MEGA & $O(cTd)$ & $O(cTd)$\\
\textbf{RWKV} & $O(\mathbf{Td})$ & $O(\mathbf{d})$ \\
\bottomrule
\end{tabular}
\caption{Complexity comparison with different Transformers: Reformer \citep{reformer}, Linear Transformer \citep{transformer_rnn}, Performer \citep{performer}, AFT\cite{zhai2021attention}, MEGA\cite{MEGA2023}. Here $T, d, c$ denote the sequence length and feature dimension, MEGA's chunk size of quadratic attention, respectively.}
\label{tab:comparison}
\end{table}


% How do we verify that we solved it?
We demonstrate the efficacy of our approach through a series of experiments on benchmark datasets and showcase the RWKV network's ability to \dots
* Scale both in number of training tokens and model parameters
* Handle long-range dependencies in text data, while keeping a linear cost on input size
* 

The advantages of the proposed approaches are folded as follows:
\begin{item}
\item Reducing time and space complexity from quadratic to linear.
\item Parallelizing reccurences with channel-oriented attention  mechanism instead of token pairwise current dot-product attention.
\item Theoretically stabilizing gradients with it's own architectures and weight initialization. 
\item Training language modeling for long context inputs that quadratic transformers are hard to take.
\end{item}

% \section{Motivation} \label{motivation}

% [eric alcaide/hypnopump draft]
% TODO: cite 
% vanishing gradients DONE
% transformer effectiveness in parallelization vs RNN, CNN DONE
% attention cost DONE
% GPT / LLaMA (for popularity) DONE
% Sequential data processing covers a wide range of scientific and industrial applications, both directly (language, time, sequences) and indirectly (other modalities, potentially formulated as sequences including images, and graphs). Models for sequential data have experienced an explosive rise in capabilities and popularity. Existing deep learning architectures for sequential data have known limitations: vanishing gradients for RNNs, slow training (RNNs, CNNs), memory and compute bottlenecks (transformers), lack of scaling (non-transformer models), difficulty handling long contexts (CNNs, transformers), lack of parallelization (RNNs), slow inference (transformers)...\\
% RWKV aims to combine strengths of previous techniques while minimizing drawbacks: parallel training and scaling like transformers but no compute/memory bottleneck ($O(n)$ vs $O(n^2)$); efficient inference like RNNs but no vanishing gradients in training. It also addresses the quadratic cost of attention in Transformer  with a reformulation to obtain "scalar attention" at linear cost without any approximation (i.e., low rank). While the approximation may be negligible with small models, it can become important when scaling up.
% Indeed, RWKV proves reliable for scaling up to billions of parameters, competitive with current SOTA at a lower cost. As such, RWKV is likely positioned to address existing challenges in further scaling and deploying AI models.
% Replace positional embedding by recurrent architecture and custom weighting of past and present information
% Reformulation of recurrence for enhanced parallelization
% Layer stacking by token/time-mixing and channel-mixing, in residual blocks
% Increased training efficiency with custom compute kernels
% Improved training dynamics by custom embedding initialization

The surge of deep learning techniques in recent years has led to impressive advancements in many scientific and industrial applications involving sequential data processing, such as natural language understanding \cite{brown2020language}, time-series analysis \cite{ismail2019deep}, and indirect modalities that can be reframed as sequences, including images and graphs \cite{wu2020comprehensive}. This progress has been catalyzed by deploying neural network architectures such as RNNs, convolutional neural networks (CNNs), and, more recently, Transformers \cite{NIPS2017_3f5ee243}. The development and success of advanced Transformer models such as GPT-3 \cite{brown2020language}, ChatGPT \cite{openai2022chatgpt,kocon2023chatgpt}, GPT-4 \cite{openai2023gpt4}, LLaMA \cite{touvron2023llama}, and Chinchilla \cite{hoffmann2022training}, underline the significance of this architecture in pushing the boundaries of what is possible in natural language processing and related fields. 

However, these architectures have inherent limitations that may restrict their usage in specific scenarios. RNNs, for instance, suffer from the vanishing gradient problem \cite{hochreiter1998vanishing,le2016quantifying} and cannot parallelize during training, limiting their scalability. CNNs, while excellent at capturing local patterns, struggle with long-range dependencies, a crucial aspect of many sequence processing tasks \cite{bai2018empirical}. On the other hand, transformers have proven to be highly effective at capturing both local and long-range dependencies and allow parallelized training \cite{tay2022efficient}. Still, their self-attention mechanism comes with significant memory and computational cost, which scales quadratically with the length of the input sequence, making them unsuitable for long sequences and settings with constrained computational resources \cite{katharopoulos2020transformers}.

To address these limitations, we introduce RWKV, a novel neural network architecture combining previous techniques' strengths while mitigating their drawbacks. RWKV offers parallelized training and robust scaling capabilities akin to Transformers but without their quadratic computational and memory bottleneck, reducing the complexity from $O(n^2)$ to $O(n)$. Moreover, it provides efficient inference, reminiscent of RNNs, without suffering from vanishing gradients during training. 

The attention through dot-product token interaction has been eschewed in favor of channel-directed attention, meaning that a token’s attention is more meaningfully quantified by which data channels are prioritized over others. This is in contrast with Transformer architecture where attention is more directly affected by specific token interactions.

Crucially, RWKV reformulates the attention mechanism resulting in \emph{linear attention}, which incurs a linear cost without approximation. While the effects of approximation in the attention mechanism might be negligible for small models, they become increasingly significant when scaling up.

This work demonstrates that RWKV can handle large-scale models with billions of parameters and is competitive with state-of-the-art architectures at a significantly reduced computational cost. Our experimental results suggest that RWKV could be a powerful tool to tackle the ongoing challenges in scaling and deploying AI models in various domains, particularly those involving sequential data processing. Thus, RWKV paves the way toward more sustainable and efficient AI models for sequence processing tasks by bridging the gap between computational efficiency and expressive capacity.
% ============================gthens the model's capacity to capture intricate relationships between sequence elements, making it well-suited for various sequential data processing tasks that requires maintaining long-distance contextual information.

\subsection{Context}
\label{context}

A defining characteristic of the RWKV architecture, as a recurrent network, is its ability to maintain a \textit{state} (or \textit{context}) across time. This feature allows the model to efficiently retain information from prior inputs while utilizing a fixed, constant memory capacity regardless of the input sequence length. This \textit{context} acts as a summary or compressed representation of the prior inputs and captures relevant information needed to generate the subsequent output. In other words, RWKV relies solely on the \textit{context} at position \textit{t} to compute the output at position \textit{t+1}.

The design inherently eliminates the need for \textit{cross-attention} mechanisms typically found in the Transformer model, leading to more efficient training and inference processes. In the context of a language model, this indicates that generation can initiate from the final state after the reading of a prompt, eliminating the requirement of having the actual prompt text.

For a more comprehensive and formal description, please refer to Appendix \ref{sec:appendix-rnn}.

% I would probably suggest to remove the following detailed explaination from here. Commenting:
%Technically, the context consists has shape $\textit{(n\_embd, 5 $\times$ %nb\_layer)}$, where each block gets 5 vectors \textit{(xx, aa, bb, pp)} where \textit{xx} refers to previous token $x_{t-1}$ used in channel-mixing (\ref{eq:channel_mix_eqs}) and also to previous token $x_{t-1}$ used in time-mixing (\ref{eq:time-mix1}) (\ref{eq:time-mix2}) (\ref{eq:time-mix3}). \textit{aa} and \textit{bb} refers to running sum in numerator and denominator of (\ref{eq:nom-denom}) and \textit{pp} is subtracted before computing the exponential in order to stabilize the softmax computation in (\ref{eq:nom-denom}).  

% Comment the following as not clear:
% In the context of LLM applications, injecting the \textit{context} into the % model is equivalent to prompt engineering or p-Tuning\cite{liu2022ptuning}. % This feature enables one copy of RWKV to serve multiple domains or purposes % with an implementation of state cache, minimizing computation overhead.

% Future work and speculation: -The state is composed by 5 vectors for each block of each stacked layer. Those vectors xx, aa, bb, pp. Empirical experiments show that those vectors have very different meanings and ranges (reference ???). -From https://github.com/BlinkDL/RWKV-LM: final hidden state（.xx .aa .bb) as a faithful sentence embedding for other tasks. Probably you should begin with .xx and .aa/.bb (.aa divided by .bb). We should probably mention the hidden state originally was 4 vectors (mathematically accurate version) and required a shift to 5 vectors for numerical stability, because during time-mixing K could get very large, so some factorization was done.
%


\subsection{Additional Optimizations} \label{adding_opts}

% In general, see https://github.com/BlinkDL/RWKV-LM/tree/main#how-it-works
% explained below

\paragraph{Custom Kernels} \label{custom_kernel}

% FROM https://johanwind.github.io/2023/03/23/rwkv_details.html
% My simplified code processes the tokens one by one, which is much slower than processing them in parallel, especially when running on GPUs. For inference, there is no way around this, as we need to sample a token before we can use it to calculate the next one. However, for training, all the text is already available. This lets us parallelize across tokens. Most of the code is fairly straightforward to parallelize like this, as there is little dependence through time. For example, all the expensive matrix multiplications work on each token independently, leading to good performance.

% However, the RWKV attention is inherently sequential. Fortunately, it has very little computation (on the order of 1024 times less than the matrix multiplications), so it should be fast. Sadly, PyTorch does not have a good way of handling this sequential task, so the attention part becomes slow (even compared to the matrix multiplications). Therefore, I wrote optimized CUDA kernels for computing the RWKV attention, which has been my main contribution to the RWKV project.

% JAX has jax.lax.scan and jax.lax.associative_scan, which allows a pure JAX implementation to perform better than pure pytorch. However, I still estimate that JAX would lead to about 40% slower training compared to CUDA (that estimate may be outdated, as it was made for training a relatively small 1.5B model).

To address inefficiencies in the $WKV$ computation due to the sequential nature of the task when using standard deep learning frameworks, we implemented custom CUDA kernels so as to launch a single compute kernel in training accelerators for this step. All other parts of the model are matrix multiplications and point-wise operations that can be already efficiently parallelized. 
% see for single kernel: https://github.com/BlinkDL/RWKV-LM/blob/main/RWKV-v4/cuda/wkv_cuda.cu#L7 
% for call to kernel, other ops in torch: https://github.com/BlinkDL/RWKV-LM/blob/main/RWKV-v4/src/model.py#L230


\paragraph{FFN with R gate}

% from : cyclekiller#6762 https://discord.com/channels/729741769192767510/1103039376184852622/1106926765760331827
Prior research \cite{DBLP:journals/corr/abs-2105-01601, liu2021pay, yu2022metaformer} suggests that self-attention may not be as essential in Transformer-based vision tasks as previously thought. Although it provided us with some insights, replacing self-attention entirely in natural language tasks could be too drastic. In our study, we partially dismantle the attention mechanism by replacing the fixed $QKV$ formula with $KV$ and introducing a new time-decaying factor $W$. This approach enables us to incorporate token and channel-mixing components akin to MLP-mixer \cite{DBLP:journals/corr/abs-2105-01601} and a gating unit $R$ similar to gMLP \cite{liu2021pay}, which enhance the performance of our RWKV model.

\paragraph{Small Init Embedding} \label{smallinit}

During the initial stage of training a transformer model (Vaswani et al., 2017), we observe that the embedding matrix undergoes slow changes, which pose a challenge for the model to deviate from its initial noisy embedding state. To mitigate this issue, we propose an approach that involves initializing the embedding matrix with small values and subsequently applying an additional LayerNorm operation. By implementing this technique, we are able to accelerate and stabilize the training process, enabling the training of deep architectures with post-LN components. The effectiveness of this approach is demonstrated in Figure \ref{fig:small_init_emb}, where it is shown to facilitate improved convergence by allowing the model to quickly transition away from the initially small embedding. This is achieved through small changes following a single step, which in turn lead to substantial alterations in directions and subsequently significant changes after the LayerNorm operation.

% check: https://github.com/BlinkDL/SmallInitEmb

% Embeddings are usually initialized with a normal (0, 1) distribution
% LRs are small values, ~1e-4. Embeddings move by that much in the beginning. This results in time spent initially in moving out of suboptimal configurations
% We propose initializing the embedings to uniform values of the same order, ~1e-4 and adding a LayerNorm afterwards, which speeds up training in the first iterations, especially for Base Pair Encoding (BPE) models. 
% We find it even allows for training of postLN models, but for the rest of this paper we focus on the preLN mode. 


\paragraph{Custom Initialization}

Building on principles from previous works \citep{he2016identity, af2_2021_nature}, we initialize parameters to values as similar as possible to an identity mapping while breaking symmetry so there is a clean information path. Most weights are initialized to zero. No biases are used for linear layers. Specific formulas are given in Appendix \ref{sec:initialization_formulas}. We find the choice of initialization to be significant in convergence speed and quality.
% TODO: Ablations or remove last sentence. 

\section{Evaluations}\label{evaluations}
In this section, we focus on evaluating to answer the following questions:

\begin{itemize}
    \item \textbf{RQ1:} Is RWKV competitive against quadratic transformer architectures for equal number of parameters and training tokens?  
    \item \textbf{RQ2:} As increasing the number of parameters, does RWKV remain competitive against quadratic transformer architectures?
    \item \textbf{RQ3:} Does increasing parameters of RWKV yield better language modeling loss, when RWKV models are trained for context lengths that most open-sourced quadratic transformers \textit{cannot} take?
\end{itemize}


Addressing \textbf{RQ1} and \textbf{RQ2}, from Fig.~\ref{fig:eval_all}, we can see that RWKV is very competitive on six benchmarks (Winogrande, PIQA, ARC-C, ARC-E, LAMBADA, and SciQ) against major open sourced quadratic complexity transformer models (Pythia, OPT and BLOOM). RWKV even outperforms Pythia and GPT-Neo for four tasks: PIQA, OBQA, ARC-E, and COPA (See details in Appendix \ref{sec-appen:eval-detail}). Addressing \textbf{RQ3}, from Fig.~\ref{fig:ctxlen_rwkv_loss},  LM loss values of RWKV with 14B parameters are less than the one of RWKV with 7B parameters when RWKV's models are trained with context length 8k which cannot take typical most open-sourced transformer models.  

\begin{figure*}[htp!]
% "[t!]" placement specifier just for this example
\subfloat[Winogrande]{%
\includegraphics[width=0.3\textwidth]{images/acc_winogrande.eps}
\label{fig:eval_winogrande}
}
\subfloat[PIQA]{%
\includegraphics[width=0.3\textwidth]{images/acc_piqa.eps}
\label{fig:eval_piqa}
}
\subfloat[ARC-Challenge]{%
\includegraphics[width=0.3\textwidth]{images/acc_arc_challenge.eps}
\label{fig:eval_arc_c}
}\\
\subfloat[ARC-Easy]{%
\includegraphics[width=0.3\textwidth]{images/acc_arc_easy.eps}
\label{fig:eval_arc_e}
}
\subfloat[LAMBADA]{%
\includegraphics[width=0.3\textwidth]{images/acc_lambada.eps}
\label{fig:eval_lambda}
}
\subfloat[SciQ]{%
\includegraphics[width=0.3\textwidth]{images/acc_sciq.eps}
\label{fig:eval_sciq}
}
\caption{Zero-Shot Performance: A horizontal axis is a number of parameters and the vertical axis is accuracy. } 
\label{fig:eval_all}
\end{figure*}

\begin{figure}[htp!]
    \centering
    \includegraphics[width=\textwidth]{images/rwkv_ctxlen_loss.png}
    \caption{Horizontal axis is context length and the vertical axis is log-log of cross entropy loss values.}
    \label{fig:ctxlen_rwkv_loss}
\end{figure}



%\begin{figure*}[htp!]
%     \centering
%     \subfloat[]{%
%     \includegraphics[scale=0.5, width=\columnwidth]{images/acc_arc_challenge.png}
%          \caption{Caption}
%     \label{fig:eval_arc_challenge}
%     }

%     \centering
%     \includegraphics[scale=0.5, width=\columnwidth]{images/acc_arc_easy.png}
     % \caption{Caption}
     % \label{fig:my_label}

%     \centering
%     \includegraphics[scale=0.5, width=\columnwidth]{images/acc_lambada.png}
%     \caption{Caption}
%     \label{fig:my_label}

%      \centering
%      \includegraphics[scale=0.5, width=\columnwidth]{images/acc_piqa.png}
%      \caption{Caption}
%      \label{fig:my_label}

%     \centering
%     \includegraphics[scale=0.5, width=\columnwidth]%{images/acc_sciq.png}
%     \caption{Caption}
 %    \label{fig:my_label}

  %   \centering
 %    \includegraphics[scale=0.5, width=\columnwidth]{images/acc_winogrande.png}
%      \caption{Caption}\label{fig:my_label}
%\end{figure*}






%Focus on equi-token equi-parameter comparisons probably?

%Compare to:
%\begin{itemize}
%    \item Pythia
%    \item OPT
%    \item BLOOM
%\end{itemize}

% refer to: https://github.com/BlinkDL/RWKV-LM/blob/main/RWKV-eval.png
% with source: https://discord.com/channels/992359628979568762/1084725482035097611/1104117503967375451

% And softer comparisons to other OpenSource models like RedPajama, LLaMA, etc. 


\section{Scaling Laws} \label{scaling laws}

% TODO: add a paragraph to explain what is scaling law
Scaling laws \cite{kaplan2020scaling} in language models refer to the mathematical relationships that describe how the performance of a language model changes with respect to various factors. These factors can include the model size ($N$), dataset size ($D$), or the optimally allocated compute budget ($C_{\rm min}$). Scaling laws help us understand how the performance of a language model scales as these factors vary.

We have observed that the RWKV model also follows scaling laws and the test loss $L$ can be estimated using a power-law relationship. This observation is depicted in Figure \ref{fig:scaling-laws}.

We conducted experiments on a range of RWKV models with different model sizes, ranging from 169M to 14B parameters. In doing so, we discovered the first scaling law:
\begin{align}
L(N) = \left(N_{\mathrm{c}}/N\right)^{\alpha_N}, \label{eq:scale_law_N}
\end{align}
where $\alpha_N \sim 0.072$ and $N_{\mathrm{c}} \sim 3.1 \times 10^{13}$ (non-embedding parameters). This scaling law demonstrates that as the model size increases, the test loss decreases without encountering any bottlenecks or plateaus, indicating that the RWKV architecture exhibits excellent scalability.

When the dataset size increases by an order of magnitude, we also observe a linear decrease in loss, which represents the second scaling law:
\begin{align}
L(D) = \left(D_{\mathrm{c}}/D\right)^{\alpha_D}, \label{eq:scale_law_N_2}
\end{align}
where $\alpha_D \sim 0.066$ and $D_{\mathrm{c}} \sim 1.4 \times 10^{15}$ (tokens)

The amount of computation is determined by two factors: the dataset size and the model size. Therefore, when there is an ample amount of data and an appropriate model size, the computational cost also follows the scaling law:
\begin{align}
L(C) = \left(C_{\mathrm{c}}/C\right)^{\alpha_C} \label{eq:scale_law_N_3},
\end{align}
where $\alpha_C \sim 0.053$ and $C_{\mathrm{c}} \sim 1.1 \times 10^{7}$ (PF-days). This indicates that with sufficient computational resources, combined with an ample amount of data and an appropriate model size, the RWKV model can continue to scale.

The next question we attempt to answer is how to make the trade-off between model size and dataset size when only a limited computing budget is available. This content has been moved to the appendix \ref{sec-appen:compute-optimal} for discussion.


% TODO: add equation to scaling laws

% TODO: this figure needs commentary
\begin{figure}[b!]
    \centering
    \includegraphics[width=\linewidth, height=\linewidth]{images/loss_compute.pdf}
 %   \caption{Scaling Laws plot for RWKV models}
    % \label{fig:scaling-laws-compute}
%\end{figure}
%\begin{figure}[htbp]
%    \centering
    \includegraphics[width=\linewidth]{images/loss_tokens.pdf}
%    \caption{Scaling Laws plot for RWKV models}
    % \label{fig:scaling-laws-tokens}
%\end{figure}
%\begin{figure}[htbp]
%    \centering
    \includegraphics[width=\linewidth]{images/loss_params.pdf}
    \caption{Scaling Laws plot for RWKV models}
    \label{fig:scaling-laws}
\end{figure}

% mgrella: I propose to cut off section 8. To make it effective I would insert comparison graphs for each experimented task but not bringing significant value at the end. 
% \section{Fundamental Experiments}
% While the RWKV model was primarily developed to compete with Transformer architectures, it is essential to evaluate its performance on typical benchmarks for recurrent networks. This section compares our RWKV architecture with well-established RNN architectures from the current literature, including LSTM and GRU.
% We used synthetic datasets with sequence lengths from 1,000 to 10,000 randomly generated floating-point numbers, challenging for LSTM and GRU due to the vanishing gradient problem.
% \textbf{Copy Task}: Models reproduce the same sequence as input.
% \textbf{Adding Task}: Given a sequence and two indices, models output the sum of the two corresponding numbers.
% \textbf{Temporal Order Verification}: Models determine if sequence elements are in ascending order.

\section{Inference Experiments} \label{infer_experiments}

In this section, we provide benchmarks and metrics on model inference requirements according to size. Specifically, we evaluate text generation speed and memory requirements on a wide variety of compute platforms. For all our experiments we use float32 precision. Performance under different quantization setups is left to further work. 

\begin{table}[ht]
    \centering
    \begin{adjustbox}{max width=\linewidth}
        \begin{tabular}{ccccc} 
 \toprule
Model/\#Params & $10^{8}$ & $4\cdot10^{8}$ & $10^{9}$ & $3\cdot10^{9}$ \\ %  & $7\cdot10^{9}$ & $10^{10}$ \\
\midrule
Pythia & xx & xxx & xxx & xxx \\
% \hline
OPT & xx & xxxx & xxx & xxx \\
% \hline
RWKV & xx & xxxx & xxx & xxx \\

\bottomrule
\end{tabular}
\end{adjustbox}
\caption{Inference time and memory for LLMs during text generation}
\label{tab:model_flop_count}
\end{table}



\section{Conclusions} \label{conclusions}

% agreed by [@mgrella - Matteo Grella, @hypnopump - Eric Alcaide] pls reach out before changing
We introduced RWKV, a new approach to RNN models exploiting the potential of time-based mixing components. RWKV introduces several key strategies which allow it to capture locality and long-range dependencies, while addressing limitations of current architectures by: (1) replacing the quadratic QK attention by a scalar formulation with linear cost, (2) reformulating recurrence and sequential inductive biases to unlock efficient training parallelization and efficient inference, and (3) enhancing training dynamics using custom initializations. 

We benchmark the proposed architecture in a wide variety of NLP tasks and show comparable performance to SoTA with reduced cost. Further experiments on expressivity, interpretability, and scaling showcase the model capabilities and draw parallelisms in behaviour between RWKV and other LLMs. 

RWKV opens a new door in scalable and efficient architectures to model complex relationships in sequential data. While many alternatives to Transformers have been proposed with similar claims, ours is the first to back up those claims with pretrained models with tens of billions of parameters.


\section{Limitations}\label{limitations}
While our proposed RWKV model has demonstrated promising results regarding training and memory efficiency during inference, some limitations should be acknowledged and addressed in future work.
% linear attention may suffer on tasks that require long context
First, the linear attention of RWKV leads to significant efficiency gains. Still, it may also limit the model's performance on tasks that require recalling information over very long contexts. This is due to the funneling of information through a single vector representation over many time steps, compared with the full information maintained by the quadratic attention of standard Transformers. In other words, the model's recurrent architecture inherently limits its ability to ``look back'' at previous tokens, as opposed to traditional self-attention mechanisms.
While learned time decay helps prevent the loss of information, it is mechanistically limited compared to full self-attention.
ss
% concrete examples ane needed to demonstrate the sensitivity of RWKV to prompts. also, please provide examples of prompts that requires look-back.
Another limitation of this work is the increased importance of prompt engineering in comparison to standard Transformer models. The linear attention mechanism used in RWKV limits the information from the prompt that will be carried over to the model's continuation. As a result, carefully designed prompts may be even more crucial for the model to perform well on tasks.

% Entries for the entire Anthology, followed by custom entries
\bibliography{custom}
\bibliographystyle{acl_natbib}

\appendix

\section{Author Contributions}

ADD YOURSELF AND YOUR CONTRIBUTIONS HERE

% Boil down contribs in visible text (add details in comments if needed) to: 
% * Initial idea (rwkv 1, 2, 3, 4, components, ...)
% * Experiments (rwkv development)
% * Paper Experiments (main body / section x / appendix y / ...)
% * Manuscript - (Initial version / Final version / Revision / Section X / Appendix Y / ...)
% * Figures
% * Tables
% [Add a contribution if not here]

%{Bo Peng* \ \ Eric Alcaide* \ \ Quentin Anthony* \\ {\bf Alon Albalak \ \  Huanqi Cao \ \  Michael Chung \ \  Matteo Grella} \\ {\bf Kranthi Kiran GV \ \  Xuzheng He \ \  Haowen Hou \ \  Jiaming Kong} \\ {\bf Hayden Lau \ \  Krishna Sri Ipsit Mantri \ \  Ferdinand Mom \ \ Xin Cheng} \\ {\bf Atsushi Saito \ \ Rui-Jie Zhu \ \  Qihang Zhao \ \ Jan Kocon \ \  Johan S. Wind} \\ {\bf Zhenyuan Zhang \ \  Peng Zhou \ \  Jian Zhu  \ \  Samuel Arcadinho \ \ Xiangru Tang}}

% bo peng: original idea, rwkv1,2,3,4 experiments over 3 years

\paragraph{Bo Peng - RWKV Foundation} Original RWKV idea, original code, performance optimizations, original experiments, and trained RWKV models from 0.1B to 14B

\paragraph{Eric Alcaide} Manuscript (initial draft sections \ref{introduction}, \ref{related_work}, final version, revision, section \ref{rwkv_model}). Figures (\ref{fig:blocks},  \ref{fig:arch}, \ref{rwkv_model}, \ref{fig:rwkv-rnn-arch}). Experiments section \ref{infer_experiments} and conclusions \ref{conclusions}

\paragraph{Quentin Anthony} test test test
\paragraph{Zhenyuan Zhang - EleutherAI} Made Fig. \ref{fig:arch}; wrote the recursive $WKV$ formula; formatted all equations to be aligned and fix punctuation; Did the information propagation experiment in Appendix \ref{sec-appen:vis}.
\paragraph{Kranthi Kiran GV - New York University} Wrote related works section (including FlashAttention, LazySoftmax, Linformer, Perceiver, S4, RMT, AFT), Evaluation section (including Table \ref{tab:commonsense_reasoning_results} and \ref{tab:opencloseqa}), FLOP count section (restored and rewrote), Background section (Transformers). Prepared and fixed formatting of several tables and equations.

\paragraph{Xiangru Tang} Wrote related works, Background, polished abstract, add two results figures of Zero-Shot Performance. Various small fixes, resizes, citations, and reorders for coherence and cohesion. Adjust the paper format, help write Tab, solve a lot of latex errors.

\paragraph{Matteo Grella - Crisis24} Wrote section \ref{grad-stability-layer-stacking} about gradient stability and layer stacking. Wrote section \ref{sec-appen:vis} about key mechanisms enabling information propagation and the (absence of) positional encoding. Wrote section \ref{context} dedicated to introduce the Context (or state) at a high level. Wrote the section \ref{conclusions}. Revised section \ref{rwkv_model} on high level architecture and detailed implementation. Contributed on section \ref{limitations}. Contributed on section \ref{introduction} emphasizing differences between self-attention and channel-directed attention.

\paragraph{Ferdinand Mom - 1) Criteo AI Lab 2) EPITA} Contribute to section 4.7 by adding details on the 5 vectors in each block state and to what they refer in the equation. Provide details on time-shift mixing (different characteristics store for different dimensions of the embedding) in section 4.6. Contribute to section 2 (linear attention without approximation). Also Proofreading.

\paragraph{Atsushi Saito - 1) Nextremer Co. Ltd. 2) Eleuther AI}   Wrote background section (part of QRNNs and  Fig:1-a,b,c), evaluation section (RQs and the other paragraphs), appendix evaluation details section(part of datasets), modified related-work section.

\paragraph{Krishna Sri Ipsit Mantri}

\paragraph{Rui-Jie Zhu} Wrote the Table \ref{tab:enwik8} and run the experiments on Enwik8. Also added Table 1 for illustrating the complexity of the RWKV. Plan to add some figures.

\paragraph{Peng Zhou} Calculate the computational complexity on Table \ref{tab:enwik8}.

\paragraph{Qihang Zhao} Adjust the paper format, help wrote Tab. \ref{tab:enwik8}, and plan to add some schematic figures. Provide supplementary explanations on the experimental configuration.

\paragraph{Xuzheng He - EleutherAI} Wrote two paragraphs for contrasting existing works with RWKV (FlashAttention, etc. and MLPMixers, etc). Added appendix \ref{sec-appen:vis}. Revised section \ref{grad-stability-layer-stacking}, \ref{sec-appen:vis}. Made cleaner versions of Figure \ref{fig:rnn_archs_background}, \ref{fig:rwkv-rnn-arch}. Proofreading and editing.

\paragraph{Michael Chung} Proofreading. Added some notes about time-mixing.

\paragraph{Haowen Hou - National University of Singapore} Wrote the text and equations for the "Scaling Laws" section \ref{scaling laws}, reploted Scaling Laws figure (currently Figure \ref{fig:scaling-laws}) for better resolution. Wrote appendix \ref{sec-appen:compute-optimal} for estimating the optimal allocation of parameters and training tokens and made figure \ref{fig:isoFLOP} to show the isoFLOP curves of RWKV model. Wrote appendix \ref{appendix_smallinit} for discussion and experiment of small init embedding. Made the figure \ref{fig:small_init_emb} to show the experiment results of small init embedding.

\paragraph{Jiaming Kong} Wrote the appendix \ref{sec:appendix-gradstab} for proving the gradient stability in RWKV mathematically. Ran some experiments and prepared a few proof-of-concept cases. Revised \ref{sec:appendix-rnn} and \ref{context}.

\paragraph{Johan S. Wind - University of Oslo} Wrote numerically stable and efficient CUDA kernels used for training RWKV models. Made formulas for parameter and FLOP counts.

\paragraph{Jian Zhu - University of British Columbia} Wrote and polished the Related Works section (adding references and condensing the whole section). Remade Figure 3 into (most of) its current form. Various small fixes across the manuscript.  

\paragraph{Huanqi Cao - Tsinghua University} Proofreading, advises about figure drawing, discovered the permanent channels, and prototyped Figure~\ref{fig:visualizations} (a).

\paragraph{Samuel Arcadinho} Inference benchmark and added related work.

\paragraph{Hayden Lau} Proofread and copy-edited the entire paper, corrected sentence structures and typos. Updated the cases section in Appendix~\ref{appendix:cases} and made small fixes in \ref{introduction}.

\paragraph{Xin Cheng - Peking University} Proofreading on the motivation section, adding references and working on the cases section in the Appendix~\ref{appendix:cases}.

\paragraph{Alon Albalak - University of California, Santa Barbara} Wrote the abstract, introduction, and limitations. Sporadically helped in other sections: fixing typos and improving formatting and such.

\paragraph{Jan Kocon - Wroclaw University of Science and Technology} Fully rewritten and merged \emph{Introduction} and \emph{Motivation} section with added references. Proofreading the whole paper. 

\section{Time-Mixing Block as an RNN Cell}
\label{sec:appendix-rnn}
As stated in \ref{rnn_inference}, the RWKV time-mixing block can be formulated as an RNN, as the $WKV$ computation can be written in such a recursive form:
\begin{align}
    a_0, b_0 &= 0, \\
    \label{eq:wkv-recur}
    wkv_t &= \frac{a_{t-1} + e^{u + k_t} v_t}{b_{t-1} + e^{u + k_t}}, \\
    a_t &= e^{-w} a_{t-1} + e^{k_t} v_t, \\
    b_t &= e^{-w} b_{t-1} + e^{k_t}.
\end{align}

The dataflow of the RNN-like time-mixing is shown in Fig. \ref{fig:rwkv-rnn-arch}, where the hidden states $h$ are the numerator-denominator tuple $(a, b)$.

\begin{figure}[htp]
    \includegraphics[clip,width=\columnwidth]{images/rwkv_as_rnn.eps}    
    \caption{RWKV time-mixing block formulated as an RNN cell. Color codes: yellow ($\mu$) denotes the token shift, red (1) denotes the denominator and blue (2) denotes the numerator computations in \ref{eq:nom-denom}. $h$ denotes the numerator-denominator tuple $(a, b)$.}
    \label{fig:rwkv-rnn-arch}
\end{figure}

To avoid overflow in calculating $e^{k_t}$, a numerical trick is used in the official implementation. Note that
\begin{align}
    a_1 &= e^{-w} 0 + e^{k_0} v_0, \\
    b_1 &= e^{-w} 0 + e^{k_0},
\end{align}
and we set $a_1 = v_0, b_1 = 1, p = k_0$, where p stores the shared exponents of a and b. Now the above recursion can be converted into a numerical safe version, for updating time step t:
\begin{align}
    q &= \max(p, u+k_t), \\
    a_t^* &= e^{p-q} a_{t-1} + e^{u+k_{t} - q} v_t, \\
    b_t^* &= e^{p-q} b_{t-1} + e^{u+k_{t} - q}\\
    wkv_t &=\frac{a_t^*}{b_t^*}.
\end{align}

The update to $a_t, b_t$ and their shared exponent are also carried out in similar fashion:
\begin{align}
    q &= \max(p - w, k_t), \\
    a_{t} &= e^{p - w - q} a_{t-1} + e^{k_t - q} v_t, \\
    b_{t} &= e^{p - w - q} b_{t-1} + e^{k_t - q}, \\
    p &= q.
\end{align}


\section{Parameter and FLOP Count for the RWKV Models}

The following section provides an overview of the different RWKV model architectures along with their respective parameter and FLOP counts in Table \ref{tab:model_flop_count}.

\begin{table}[htbp]
    \centering
    \begin{adjustbox}{max width=\linewidth}
        \begin{tabular}{ccccc} 
 \toprule
Name & Layers & Model Dimension & Parameters & FLOPs per token \\
\midrule
\SI{169}{M} & 12 & 768 & \num{1.693e+08} & \num{2.613e+08} \\
% \hline
\SI{430}{M} & 24 & 1024 & \num{4.304e+08} & \num{7.573e+08} \\
% \hline
\SI{1.5}{B} & 24 & 2048 & \num{1.515e+09} & \num{2.823e+09} \\
% \hline
\SI{3}{B} & 32 & 2560 & \num{2.985e+09} & \num{5.710e+09} \\
% \hline
\SI{7}{B} & 32 & 4096 & \num{7.393e+09} & \num{1.437e+10} \\
% \hline
\SI{14}{B} & 40 & 5120 & \num{1.415e+10} & \num{2.778e+10} \\
\bottomrule
\end{tabular}
\end{adjustbox}
\caption{RWKV model architectures and associated FLOP counts}
\label{tab:model_flop_count}
\end{table}


The number of parameters for each model is computed using the formula:  $\#parameters = 2VD + 13D^2L + D(11L+4)$ where $V$ = 50277 is the vocabulary size, $D$ represents the Model Dimension and L corresponds to the number of layers.

FLOPs is for a forward pass for one token. It was calculated as $6(VD + 13D^2L)$, which is the twice (add and multiply) the number of parameters in linear layers. The backwards pass FLOPs can be approximated as twice that of the forward pass. So the total is $6(VD + 13D^2L)$ per token for training (3x fw FLOPs). It is noteworthy that FLOPs are independent of the context length, unlike regular transformers. The FLOP approximations in this paper are in line with the methodology used by \citet{kaplan2020scaling}.

Alternative approximations for FLOPs include doubling the parameters which yields similar results within 2\% for 14B and a 30\% discrepancy for 169M variant. Another approximation is based on the number of non-embedding parameters multiplied by 2. This gives $2(VD + 13D^2L + D(11L+4))$ resulting in 1.6\% more FLOPs for 14B model and 8\% more FLOPs for 169M model.


\section{Parameter initializations}
\label{sec:initialization_formulas}

We describe the specific parameter initializations below and motivate the design choices. Parameters belonging to residual blocks are often adjusted by layer depth and total number of layers. Let $\#$ denote the vocabulary size, $s$ denote the embedding dimension, $d$ denote the hidden size (we use $d=4s$), $L$ the number of layers, $l$ the layer index (from 0 to L-1), we use the following initializations:

\begin{itemize}
    \item Embeddings are initialized to a pseudo orthogonal matrix with gain  $\num{1e-4}\sqrt{max(\#, s)})$ as explained in \ref{smallinit}
    \item For the channel-mixing blocks (\ref{eq:time-mix1}), $\mu_{k_i}$ and $\mu_{r_i}$ are initialized to $(\frac{i}{s})^{1-\frac{l}{L}}$
    \item For the time-mixing blocks (\ref{eq:channel_mix_eqs}), initializations are $\mu_{k_i}=(\frac{i}{s})^{1-\frac{l}{L}}$, $\mu_{v_i}=(\frac{i}{s})^{1-\frac{l}{L}}+\frac{0.3l}{L-1}$ and $\mu_{r_i}=0.5(\frac{i}{s})^{1-\frac{l}{L}}$
    \item $w_i$ (\ref{eq:nom-denom}), also known as ``time decay'', is initialized to $-5+8 \cdot (\frac{i}{d-1})^{0.7+\frac{1.3l}{L-1}}$. Intuitively, it is the discount factor applied to previous tokens over time.  
    % for u: https://github.com/BlinkDL/RWKV-LM/blob/main/RWKV-v4/src/model.py#L184
        \item $u_i$ (\ref{eq:nom-denom}), also known as ``bonus'', is set to $0.5(((i+1) \mod 3) - 1)+\log0.3$. It is the special weighting applied to the current token in equation \ref{eq:nom-denom}. The alternating zigzag pattern initially creates subtle variations in the tensor elements, which are intended to help the model treat different dimensions of the embedding distinctively.
    \item $W_o$ (\ref{eq:time-mix4}) (time-mixing) and $W_v$ (channel-mixing) are initialized to $\mathcal{N}(0, \sqrt{\frac{d}{s}}=2)$
    \item All $W_r, W_k, W_v$ weights are initialized to 0 so the model can start learning from the beginning without noisy signals. 
    \item All LayerNorm weights start from 1 and biases from 0.
\end{itemize}

\section{Small Init Embedding}  \label{appendix_smallinit}

This section presents experimental validation of small initialization embedding. The experimental setup is as follows. In the baseline configuration, the parameters are initialized using a normal distribution with a mean of 0.0 and a standard deviation of 0.02, which is a commonly used initialization method in models like BERT and GPT. On the other hand, in the small initialization of the embedding (small init emb) experiment, the parameters are initialized using a uniform distribution with a range of 1e-4, which is slightly different from RWKV where a normal distribution with a standard deviation of 1e-4 is used. However, this difference is negligible and does not affect our conclusions. The experiments were conducted with a batch size of 400. As depicted in the figure \ref{fig:small_init_emb}, the loss curve for the small init emb exhibits a faster rate of decrease and convergence compared to the traditional initialization using a normal distribution.

\begin{figure}[htp]\label{small_init_emb}
    \centering
    \includegraphics[height=0.83\linewidth]{images/small-init-emb.png}    
    \caption{Effect of small initialization embedding}
    \label{fig:small_init_emb}
\end{figure}

\section{Gradient Stability in RWKV} 
\label{sec:appendix-gradstab}

In this section we formulate a mathematical description about the behavior of gradients in RWKV. Specifically we want to address how the novel design of RWKV's layers alleviated the vanishing gradient problem. Take the partial derivative of $L$ (cross entropy between \(O_T\) and \(\text{label}_T\)), to \(W_v\) from one time-mixing layer for example, at time T, is expressed as formula \ref{eq:pd_to_wk}. Note that \(O_T\) only relate to \(W_v\) through
\(wkv_{T}\), let $\alpha$ be the normalizing coefficient introduced by the LayerNorm before time-mixing:

\begin{align}
\label{eq:pd_to_wk}
\frac{\partial L}{\partial W_v} &= \alpha\left(\frac{\partial L}{\partial wkv_T}\sum_{t=1}^T
\frac{\partial {wkv_T}}{\partial v_t}\frac{\partial {v_t}}{\partial W_v}\right) \\
&\propto \left(\sum_{t=1}^T  \frac{\partial {wkv_T}}{\partial v_t} \frac{\partial v_t}{\partial W_v}\right).
\end{align}

In the context of vanishing gradient problem for recurrent network, the risk of RWKV\textquotesingle s design would be that for some relatively large $T$ values and small $t$ values, $\frac{\partial {wkv_T}}{\partial v_t}$ \emph{could be near zero}. And the model would not learn anything useful inside the whole context window. What we want to prove here is that under mild assumptions, and the use of LayerNorms, this will not be the case. All \(J(i)\) for any $i$ will be bounded at some threshold.

Now we can simplify the above expression, ignoring $\alpha$ and $\frac{\partial L}{\partial wkv_T} = w_0 \sigma(r_T)$ which are considered external coefficients:
\begin{align}
\label{eq:contri_to_grad}
\frac{\partial {wkv_T}}{\partial v_t} &= \frac{e^{-(T-1-t)w + k_t}}{{\sum_{i=1}^{t-1} e^{-(t-1-w)w+k_i} + e^{u+k_t}}},\\
\frac{\partial v_t}{\partial W_v} &=\mu_v x_t + (1-\mu_v) x_{t-1}.
\end{align}
Note that equation \ref{eq:contri_to_grad} is very similar to that in \cite{zhai2021attention}. The contributions of gradient from token $x_t$ or $v_t$ to a parameter $W_v$, are mostly controlled by a small decaying factor $w$, and the learned key \(k_t\). The Softmax operation guarded any $k_t$ from total control of the final gradient. Under some mild assumptions about the initial embedding, we can safely assume all tokens in the context window are accounted for. Also, it is trivial to see that the total gradient of \(W_v\) in one update will also be bounded:

\begin{align}
\sum_{t=1}^T \frac{\partial {wkv_T}}{\partial v_t} &=\sum_{t=1}^T \frac{\partial wkv_T}{\partial v_t}\\
&\leq \frac{\sum_{t=1}^T e^{-(T-1-t)w + k_t}}{\sum_{i=1}^{t-1} e^{-(t-1-w)w+k_i} + e^{u+k_t}}\\
&\leq 1.
\end{align}

Likewise, the partial derivative of $L$ w.r.t. $W_k$ can be derived as:

\begin{align}
\frac{\partial L}{\partial W_k} &= \alpha\left(\frac{\partial L}{\partial wkv_T}\sum_{t=1}^T
\frac{\partial {wkv_T}}{\partial k_t}\frac{\partial {k_t}}{\partial W_k}\right)\\
&\propto \frac{e^{u+W_k X_K^T}X_K^T \sum(ev) W_v}{(e^{u+W_k X_K^T} + \sum(e))^2}\\
&\propto \frac{\lambda \sum(ev)}{(\lambda + \sum(e))^2},
\end{align}
where
\begin{align}
    \sum(ev) &= \sum_{i=1}^{t-1}e^{(1+i-t)w+k_i} v_i, \\
    \sum(e) &= \sum_{i=1}^{t-1}e^{(1+i-t)w+k_i}, \\
    X_K^T &= \mu_k x_T + (1-\mu_k)x_{T-1}, \\
    X_V^T &= \mu_v x_T + (1-\mu_v)x_{T-1}, \\
    \lambda &= e^{u+W_k X_K^T}X_K^T.
\end{align}
We observe that the contribution of each token $x_i$ to the gradient of $W_k$ is also controlled in a softmax weighted manner and learned from training data.

The partial gradients of $L$ w.r.t. $W_k, W_v$ in channel-mixing have very similar behavior except for the squared ReLU activation function. The stability of these gradients can be shown using similar techniques.

\section{Model Behavior Visualization}\label{sec-appen:vis} % just a draft name
\begin{figure}[htbp]
    \centering
    \includegraphics[width=\linewidth]{images/visualize-decay.pdf}
    \includegraphics[width=\linewidth]{images/visualize-information-propagation.pdf}
%     \includegraphics[width=\linewidth, height=0.8\linewidth]{images/loss_params.png}
    \caption{Model behavior visualizations of the RWKV model.}
    \label{fig:visualizations}
\end{figure}

In Figure \ref{fig:visualizations}, we present visualizations of some behavior of the RWKV model.

The top plot illustrates the time decays ($e^{-w}$) in each layer in the RWKV-169M model, sorted along the channel axis. Notably, several decays in the last layers are very close or equal to one, implying that certain information is preserved and propagated throughout the model's temporal context. Meanwhile, many decays in the initial layer are close to zero, which corresponds to local operations in $wkv$ (\ref{eq:nom-denom}), likely to be associated with tasks such as text parsing or lexical analysis. These characteristics are also observed in larger models.

% From @cryscan https://discord.com/channels/992359628979568762/1085545872533758112/1104973719455158322
The bottom plot shows the information retrieval and propagation path in the RWKV-430M model. The experiment follows the \emph{causal trace} method introduced by \citet{meng2022locating}, where we
\begin{enumerate}
    \item Run the model once, recording all states and activation of each layer during the computation;
    \item Corrupt the subject input embeddings using noise (``The Eiffel Tower'' in this example);
    \item Restore the states and activation of a certain layer at a certain token during the computation, recording the log-probability of the model outputting the correct answer (``Paris'').
\end{enumerate}
Unlike transformers, RWKV relies on recursive propagation of information in the time dimension.
In this case, the fact that "the Eiffel Tower is located in Paris" is retrieved in layer 4. It is then passed down to the subsequent layers. In layer 20, it is propagated through time to the location where it is needed. Finally, it is passed down to the last layer.

\section{Evaluation Details}\label{sec-appen:eval-detail}


The results for following tasks are in Table \ref{tab:commonsense_reasoning_results} and \ref{tab:opencloseqa}.


Tasks:
\begin{itemize}
    \item LAMBADA~\cite{LAMBADAdataset}. A benchmark dataset that evaluates the model's contextual reasoning and language comprehension abilities by presenting context-target pairs, where the objective is to predict the most probable target token. %A benchmark dataset designed to measure contextual reasoning and language comprehension abilities of the model. Each example in the dataset consists of a context-target pair, where the task involves reading the context and selecting the most likely target word from multiple candidate words.
    \item PIQA~\cite{Bisk2020}. A benchmark for the task of physical common sense reasoning, which consists of a binary choice task that can be better understood as a set of two pairs, namely (Goal, Solution).% A Benchmark Dataset for the Task of Physical Common Sense Reasoning. PiQA is a binary choice task that can be better understood as a set of two pairs, namely (Goal, Solution).
    \item HellaSwag ~\cite{HellaSwag2019} A novel benchmark for commonsense Natural Language Inference (NLI) which is build by adversarial filtering against transformer models. 
    % such that a sequence of discriminators is utilized to choose a difficult set of incorrect answers generated during the process.
    \item Winogrande ~\cite{Wino2020} A dataset designed to evaluate the acquisition of common sense reasoning by neural language models, aiming to determine whether we are accurately assessing the true capabilities of machine common sense.
    \item StoryCloze~\cite{StoryCloze2016} A benchmark to present a novel approach to assess comprehension of narratives, narrative generation, and script acquisition, focusing on commonsense reasoning.
    \item ARC Challenge ~\cite{ARC2018} A dataset designed for multiple-choice question answering, encompassing science exam questions ranging from third grade to ninth grade.
    \item ARC Easy An easy subset of ARC.
    \item HeadQA ~\cite{HeadQA2020} A benchmark consisting of graduate-level questions encompassing various fields such as medicine, nursing, biology, chemistry, psychology, and pharmacology.
    \item OpenBookQA ~\cite{OpenBookQA2018} A QA dataset to evaluate human comprehension of a subject by incorporating open book facts, scientific knowledge, and perceptual common sense, drawing inspiration from open book exams.
    \item SciQ ~\cite{SciQ2017} A multiple-choice QA dataset which was created using an innovative approach to gather well-crafted multiple-choice questions that are focused on a specific domain. 
    \item TriviaQA ~\cite{TriviaQA2017} A QA-IR dataset which is constituted of triples of questions, answers, supporting evidence, and independently collected evidence documents, with an average of six documents per question for reliable sources.
    \item ReCoRD ~\cite{ReCord} A benchmark for evaluating commonsense reasoning in reading comprehension by generating queries from CNN/Daily Mail news articles and requiring text span answers from corresponding summarizing passages.
    \item COPA ~\cite{COPA2011} A dataset to evaluate achievement in open-domain commonsense causal reasoning.
    \item MMMLU ~\cite{MMMLU2021} A multi-task dataset for 57 tasks containing elementary mathematics, US history, computer science, law, etc. 
\end{itemize}


\begin{table*}[!]
\centering
\small
\begin{tabular}{llllllllll}
        \toprule
\textbf{Model} & \textbf{Params} & \textbf{PIQA} & \textbf{StoryCloze} & \textbf{HellaSwag} & \textbf{WinoGrande} & \textbf{ARC-e} & \textbf{ARC-c} & \textbf{OBQA} \\
 & B & acc & acc & acc\_norm & acc & acc & acc\_norm & acc\_norm \\
\midrule
RWKV-4 &  0.17 & \TBF{65.07} & \TBF{58.79} & \TBF{32.26} & 50.83 & \TBF{47.47} & \TBF{24.15} & \TBF{29.60} \\
Pythia & 0.16      & 62.68 & 58.47 & 31.63 & \TBF{52.01} & 45.12 & 23.81 & 29.20 \\
GPT-Neo & 0.16     & 63.06 & 58.26 & 30.42 & 50.43 & 43.73 & 23.12 & 26.20 \\
\midrule
RWKV-4 &  0.43 & \TBF{67.52} & \TBF{63.87} & \TBF{40.90} & 51.14       & \TBF{52.86} & 25.17       & \TBF{32.40} \\
Pythia & 0.40      & 66.70       & 62.64      & 39.10      & \TBF{53.35} & 50.38       & \TBF{25.77} & 30.00 \\
GPT-Neo & 0.40     & 65.07       & 61.04       & 37.64      & 51.14     & 48.91       & 25.34 & 30.60 \\
\midrule
RWKV-4 & 1.5 & \TBF{72.36} & \TBF{68.73} & \TBF{52.48}   & 54.62      & \TBF{60.48} & \TBF{29.44}             & \TBF{34.00} \\
Pythia & 1.4 & 71.11       & 67.66       & 50.82        & \TBF{56.51} & \TBF{57.74} & 28.58             & 30.80 \\
GPT-Neo & 1.4 & 71.16      & 67.72        & 48.94       & 54.93       & 56.19       & 25.85             & 33.60 \\
\midrule
RWKV-4 & 3.0 & \TBF{74.16} & \TBF{70.71} & \TBF{59.89} & 59.59       & \TBF{65.19}         & 33.11 & \TBF{37.00} \\
Pythia & 2.8 & 73.83       & \TBF{70.71} & 59.46       & \TBF{61.25} & 62.84   & \TBF{32.25} & 35.20 \\
GPT-Neo & 2.8 & 72.14      & 69.54       & 55.82        & 57.62      & 61.07         & 30.20       & 33.20 \\
\midrule
RWKV-4 & 7.4 & \TBF{76.06} & 73.44       & 65.51      & 61.01       & \TBF{67.80} & \TBF{37.46}         & \TBF{40.20} \\
Pythia & 6.9 & 74.54       & 72.96       & 63.92      & 61.01       & 66.79       & 35.07               & 38.00 \\
GPT-J & 6.1  & 75.41        & \TBF{74.02 }& \TBF{66.25}& \TBF{64.09} & 66.92 & 36.60                    & 38.20 \\
\midrule
RWKV-4 & 14.2    & 77.48            & 76.06   & 70.65       & 63.85       & 70.24       & 38.99       & \TBF{41.80} \\
GPT-level & 14.2 & 76.49           & 74.97   & 68.72       & 65.14      & 70.77        & 37.99       & 39.27 \\
Pythia & 11.8    & 75.90           & 74.40   & 67.38       & 64.72       & 69.82        & 36.77       & 38.80 \\
GPT-NeoX & 20.6  & 77.69     & 76.11 & 71.42 & 65.98       & 72.69 & 40.44 & 40.20 \\
\bottomrule
\end{tabular}
\caption{\label{tab:commonsense_reasoning_results}
Zero-Shot Performance of the model on Common Sense Reasoning Tasks
}
\end{table*}


\begin{table*}[htp!]
\centering
\small
\begin{tabular}{lllllllll}
\toprule
\textbf{Model} & \textbf{Params} & \textbf{LAMBADA} & \textbf{LAMBADA} & \textbf{headQA} & \textbf{sciq} & \textbf{triviaQA} & \textbf{ReCoRD} & \textbf{COPA} \\
  & B & ppl & acc & acc\_norm & acc & acc & em & acc \\
\midrule
RWKV-4  & 0.17 & \TBF{29.33} & 32.99 & 25.78 & \TBF{77.50} & 1.26 & 62.03 & \TBF{66.00} \\
Pythia  & 0.16 & 24.38 & \TBF{38.97} & \TBF{25.82} & 76.50 & \TBF{1.31} & \TBF{66.32} & 62.00 \\
GPT-Neo & 0.16 & 30.27 & 37.36 & 25.16 & 76.60 & 1.18 & 64.92 & 64.00 \\
\midrule
RWKV-4  & 0.43 & 13.04 & 45.16 & \TBF{27.32} & 80.30 & \TBF{2.35} & 70.48 & 65.00 \\
Pythia  & 0.40  & \TBF{11.58} & \TBF{50.44} & 25.09 & \TBF{81.50} & 2.03 & \TBF{75.05} & \TBF{67.00} \\
GPT-Neo & 0.40  & 13.88 & 47.29 & 26.00 & 81.10 & 1.38 & 73.79 & 65.00 \\
\midrule
RWKV-4  & 1.5  & 7.04  & 56.43 & \TBF{27.64} & 85.00 & \TBF{5.65} & 76.97 & \TBF{77.00} \\
Pythia  & 1.4  & \TBF{6.58}  & \TBF{60.43} & 27.02 & \TBF{85.50 }& 5.52 & \TBF{81.43} & 73.00 \\
GPT-Neo & 1.4  & 7.5   & 57.25 & 27.86 & 86.00 & 5.24 & 80.62 & 69.00 \\
\midrule
RWKV-4  & 3.0    & 5.25  & 63.96 & 28.45 & 86.50 & \TBF{11.68} & 80.87 & \TBF{82.00} \\
Pythia  & 2.8  & \TBF{4.93}  & \TBF{65.36} & \TBF{28.96} & \TBF{87.70} & 9.63  & 85.10 & 77.00 \\
GPT-Neo & 2.8  & 5.63  & 62.22 & 27.17 & 89.30 & 4.82  & \TBF{83.80} & 80.00 \\
\midrule
RWKV-4  & 7.4  & 4.38  & 67.18 & \TBF{31.22} & 88.80 & \TBF{18.30} & 83.68 & \TBF{85.00} \\
Pythia  & 6.9  & 4.3   & \TBF{67.98} & 28.59 & 90.00 & 15.42 & 86.44 & \TBF{85.00} \\
GPT-J   & 6.1  & \TBF{4.1}   & 68.31 & 28.67 & \TBF{91.50} & 16.74 & \TBF{87.71} & 83.00 \\
\midrule
RWKV-4  & 14.2 & \TBF{3.86}  & 70.83 & \TBF{32.64} & 90.40 & \TBF{24.58} & 85.67 & \TBF{85.00} \\
GPT-level & 14.2 & 3.81 & \TBF{70.94} & 31.03 & \TBF{92.20} & 22.37 & \TBF{87.89} & 82.66 \\
\midrule 
Pythia  & 11.8 & 3.89  & 70.44 & 30.74 & 91.80 & 20.57 & 87.58 & 82.00 \\
GPT-NeoX & 20.6 & 3.64 & 71.94 & 31.62 & 93.00 & 25.99 & 88.52 & 84.00 \\
\bottomrule
\end{tabular}
\caption{\label{tab:opencloseqa}
Zero-Shot Performance of various models on different tasks.
}
\end{table*}

\begin{table*}[htbp]
\small
    \centering
    \begin{tabular}{llllllllllllll}
 \toprule
 Method & L& d  &T &Train bpc & Test bpc & Time Complexity & Space Complexity\\
 \midrule
 Transformer&  12& 512& 1024&0.977& 1.137&$O(T^2d)$&$O(T^2 + Td)$ \\
 Transformer&  24& 256& 1024&1.039& 1.130&$O(T^2d)$&$O(T^2 + Td)$ \\
    \midrule
 Reformer& 12 & 512& 1024& 1.040 &  1.195& $O(T\log{T}  d)$&$O(T\log{T} + Td)$ \\
 Synthesizer&  12& 512& 1024& 0.994 & 1.298& $O(T^2d)$&$O(T^2 + Td)$  \\
 Linear Transformer& 12& 512& 1024&0.981& 1.207& $O(T d^2)$&$O(Td + d^2)$ \\
 Performer&  12& 512& 1024&1.002& 1.199& $O(T d^2\log{d})$&$O(Td\log{d} + d^2\log{d})$  \\
  AFT-simple &  12& 512& 1024&1.046& 1.209& $O(Td)$&$O(Td)$  \\
 \midrule
RWKV-RNN &  6 & 512  & 1024 & 0.720 & -  & $O(\mathbf{Td})$&$O(\mathbf{d})$ \\
  \bottomrule    
    \end{tabular}
        \caption{Enwik8 results, measured in bits per character (bpc): the lower the better. Baseline comparisons are made with Reformer \citep{reformer}, Synthesizer \citep{synthesizer} (the best performing dense version), Linear Transformer \citep{katharopoulos2020transformers}, Performer \citep{performer}. $L, d,$ and $T$ denote the number of blocks (network depth), dimension of features, and sequence length, respectively. Both Linear Transformer and Performer are implemented with customized CUDA kernels (github.com/idiap/fast-transformers), and all other models are implemented in native Pytorch. }
\label{tab:enwik8}
\end{table*}{}

\section{Estimating the optimal allocation of parameters and training tokens}\label{sec-appen:compute-optimal}
In our study, we investigate the optimal model size and number of tokens for training a RWKV language model while considering a fixed compute budget. We systematically adjust the model size while keeping the training FLOP count constant across six different values. We select the number of training tokens for different model sizes to ensure a constant final FLOP count. The range of training FLOP counts spans from 1 PF-day to 300 PF-days. By following \cite{hoffmann2022training}, we plot the IsoFLOPs curves (Figure \ref{fig:isoFLOP}). In Figure \ref{fig:isoFLOP}, we observe a distinct minimum point in the loss curve, indicating the existence of a compute-optimal RWKV model to be trained within a given FLOP budget. On the curve, points located to the left of the minimum value indicate a scenario where there is abundant data but the model size is too small, suggesting the possibility of increasing the model size. Conversely, points located to the right of the minimum value indicate a situation where the model is too large compared to the available data, suggesting the possibility of reducing the model size.

We utilize a parabolic curve fitting method to determine the model size that achieves the minimum loss for each IsoFLOPs curve (Figure \ref{fig:isoFLOP}). Subsequently, we establish a power law relationship between the amount of compute ($C$) and the optimal model size ($N_{opt}$) as well as the optimal number of training tokens ($D_{opt}$). The exponents of the power law, denoted as $N_{opt} \propto C^a$ and $D_{opt} \propto C^b$, are estimated to be $a=0.519$ and $b=0.481$, respectively. Specifically, our findings suggest that with each increase in compute, a proportional increase in both data size and model size yields optimal results. For example, if compute is increased by a factor of 10, the model size and data size should be increased by approximately 3.1 times. Similarly, if compute is increased by a factor of 100, the model size and data size should be increased by approximately 10 times. This indicates a consistent relationship between compute, model size, and data size in achieving optimal performance.

\begin{figure}[htbp]
    \centering
    \includegraphics[width=\linewidth]{images/iso_curve.png}
    \caption{IsoFLOP curves for RWKV model. Each dashed line represents a fixed compute budget.}
    \label{fig:isoFLOP}
\end{figure}


\onecolumn
\clearpage
\section{Cases}
\label{appendix:cases}

% this section is to show some real cool cases about RWKV

% todo:
% 1. find more examples
% 2. add raven icon into the templeate

% here is a template, and the content is from https://huggingface.co/blog/rwkv

In this part, we present a few instances of outputs produced by RWKV-4. 

% TODO: @hynpopump: add better outputs and provide model sizes

\begin{framed}
\centering
\begin{minipage}{.9\textwidth}
\textbf{Prompt:} In a shocking finding, scientist discovered a herd of dragons living in a remote, previously unexplored valley, in Tibet. Even more surprising to the researchers was the fact that the dragons spoke perfect Chinese.\\

\rwkvemoji \textbf{Response:} The researchers found that the dragons were able to communicate with each other, and that they were
\end{minipage}
\end{framed}

\begin{framed}
\centering
\begin{minipage}{.9\textwidth}
\textbf{Instruction:} Tell me about ravens.\\

\rwkvemoji \textbf{Response:} RAVENS are a type of bird that is native to the Middle East and North Africa. They are known for their intelligence, adaptability, and their ability to live in a variety of environments. RAVENS are known for their intelligence, adaptability, and their ability to live in a variety of environments. They are known for their intelligence, adaptability, and their ability to live in a variety of environments.
\end{minipage}
\end{framed}

\begin{framed}
\centering
\begin{minipage}{.9\textwidth}
\textbf{Instruction:} Explain the concept of machine learning to a 10-year-old.  \\

\rwkvemoji \textbf{Response: }Machine learning is like a super smart computer that can learn from lots of different things and then use that knowledge to do new things. Just like how you can teach a robot to do some tricks, like rolling over or jumping through hoops, machine learning can help computers understand new information and make decisions based on what they've learned. It's kind of like teaching a robot how to be a good teacher!\textbf{
}
\end{minipage}
\end{framed}

\begin{framed}
\centering
\begin{minipage}{.9\textwidth}
\textbf{Instruction:} Write a short poem describing nature's beauty in springtime..\\

\rwkvemoji \textbf{Response: }In springtime, nature's beauty abounds,\\
                                            As flowers bloom and trees come alive.\\
                                            The air is fresh and filled with life,\\
                                            And everything feels new and bright.\\
                              The birds sing sweetly in the trees,\\
                              While bees buzz from flower to flower.\\
                              The gentle breeze carries their sweet scent,\\
                              And all around us, life abounds.
\end{minipage}
\end{framed}

\begin{framed}
\centering
\begin{minipage}{.9\textwidth}
\textbf{Instruction:} Below is the numerical stable formula for updating the state variable in RWKV:\\

\begin{aligned}
    a_t^* &= e^{p-q} a_{t-1} + e^{u+k_{t} - q} v_t\\
    b_t^* &= e^{p-q} b_{t-1} + e^{u+k_{t} - q}\\
    wkv_t &=\frac{a_t^*}{b_t^*}
\end{aligned}
Please write out the corresponding pytorch code to do so, where both $a_t$ and $b_t$ are torch.Tensor([n])\\
\textbf{(The actual input was in Markdown syntax)}\\\\

\rwkvemoji \textbf{Response: }
Here is the PyTorch code to update the state variable in RWKV using the given formula:\\

\begin{lstlisting}[language=Python]
import torch
def update_state(a, b, u, v, p, q):
    a_t = torch.exp(p - q) * a + torch.exp(u + k - q) * v
    b_t = torch.exp(p - q) * b + torch.exp(u + k - q)
    return a_t, b_t
a, b = torch.randn(2, 3)
u = torch.randn(2)
v = torch.randn(2)
wkv = torch.randn(2)
a_t, b_t = update_state(a, b, u, v, wkv, torch.randn(2))
print("a_t:", a_t)
print("b_t:", b_t)
\end{lstlisting}\\\\

This code defines a function update\_state that takes in the state variables a and b,\\
the control variables u and v, and the parameters wkv and torch.randn. It then uses the\\
PyTorch torch.exp function to compute the exponential of the control variables and the\\
log of the normalization factor. Finally, it uses these values to update the state\\

\end{minipage}
\end{framed}

\begin{framed}
\centering
\begin{minipage}{.9\textwidth}
\textbf{Instruction:} Explain the following metaphor: Life is like cats.\\

\rwkvemoji \textbf{Response: }The metaphor "Life is like cats" suggests that life can be unpredictable and full of unexpected twists and turns, just like a cat's behavior. It implies that we cannot fully control or predict the outcome of our lives, just as we cannot control the movements of a cat. The metaphor also implies that life is complex and multifaceted, with many different variables and factors at play. Just as a cat can have multiple personalities or moods, life can have many twists and turns that are difficult to predict or understand. Overall, the metaphor suggests that life is complex and unpredictable, much like a cat's behavior.\\
\end{minipage}
\end{framed}


\end{document}


% ============================
% OLD INTRO & MOTIVATION (COPY FROM 2023.05.18)
\section{Introduction} \label{introduction}
\textcolor{blue}{Introduction is a placeholder, Alon will go back and re-do introduction once we have more concrete details of what is being done in this work.}
% What are we trying to do, and why is this relevant?
% Develop efficient sequential NN architecture (efficient in training and inference)
Since the introduction of the transformer by Vaswani et al.~\cite{NIPS2017_3f5ee243}, its success has been attributed to the self-attention mechanism. This powerful and flexible method allows models to capture long-range dependencies and complex interactions between input elements. The inherent parallelizability of self-attention combined with ever-increasing dataset sizes, and its ability to create rich contextual representations, has led to state-of-the-art models across a multitude of benchmarks and applications.
% Transformers have become dominant architectural paradigm, but have quadratic computational and space complexity due to self-attention mechanism
However, despite the transformative impact of the Transformer, it is not without its shortcomings. One notable drawback is the quadratic complexity of its self-attention mechanism. It renders the architecture computationally expensive for tasks involving long input sequences and in situations with restricted memory budgets. These limitations have spurred a wealth of research aiming to improve the scaling properties of the Transformer, often at the expense of sacrificing some of the properties that make it so effective~\citep{wang2020linformer, zaheer2020bigbird, dao2022flashattention}.


% Quadratic complexity limits the use cases (on-device, conversations, etc.)
% SEE: 4.3 RNN-like Formulation for current explanation of advantages

% Why is developing efficient NN architectures hard?
Addressing the issue of quadratic scaling has proven to be a challenging endeavor, as it requires navigating the intricate balance between the computational efficiency of the self-attention mechanism and the expressive capacity of the model. Furthermore, the non-local nature of the self-attention mechanism poses significant difficulty in devising more efficient alternatives that do not compromise the rich expressiveness of the Transformer.

% How do we solve it?
In this paper, we propose a novel neural network architecture, dubbed the RWKV network, which serves as a variant of both the RNN and the Transformer. The RWKV network is designed to alleviate the quadratic scaling issue and replace it with a linear one while retaining the crucial properties that have made the Transformer the dominant architecture in the field
% Go more into depth on what these crucial properties are
. By leveraging the strengths of both RNNs and Transformers, the RWKV network presents a compelling alternative for tackling a wide range of tasks that demand both computational efficiency and expressiveness. In doing so, the time and space complexity is sharply decreased. Our approach is very friendly to large model size and computation resources. See Table~\ref{tab:comparison} for the complexity comparison of RWKV and related approaches. 

\begin{table}
\centering
\setlength{\tabcolsep}{2pt}
\small
\begin{tabular}{l|c|c}
\toprule
\textbf{Model} & \textbf{Time} & \textbf{Space}\\
\midrule
Transformer & $O(T^2d)$ & $O(T^2 + Td)$ \\
Reformer & $O(T\log{T}  d)$ & $O(T\log{T} + Td)$ \\
% Synthesizer & $O(T^2d)$ & $O(T^2 + Td)$ \\
Linear Transformers & $O(T d^2)$ & $O(Td + d^2)$ \\
Performer & $O(T d^2\log{d})$ & $O(Td\log{d} + d^2\log{d})$ \\
AFT-full & $O(T^2d)$ & $O(Td)$ \\
MEGA & $O(cTd)$ & $O(cTd)$\\
\textbf{RWKV} & $O(\mathbf{Td})$ & $O(\mathbf{d})$ \\
\bottomrule
\end{tabular}
\caption{Complexity comparison with different Transformers: Reformer \citep{reformer}, Linear Transformer \citep{transformer_rnn}, Performer \citep{performer}, AFT\cite{zhai2021attention}, MEGA\cite{MEGA2023}. Here $T, d, c$ denote the sequence length and feature dimension, MEGA's chunk size of quadratic attention, respectively.}
\label{tab:comparison}
\end{table}


% How do we verify that we solved it?
We demonstrate the efficacy of our approach through a series of experiments on benchmark datasets and showcase the RWKV network's ability to \dots
* Scale both in number of training tokens and model parameters
* Handle long-range dependencies in text data, while keeping a linear cost on input size
* 

The advantages of the proposed approaches are folded as follows:
\begin{item}
\item Reducing time and space complexity from quadratic to linear.
\item Parallelizing reccurences with channel-oriented attention  mechanism instead of token pairwise current dot-product attention.
\item Theoretically stabilizing gradients with it's own architectures and weight initialization. 
\item Training language modeling for long context inputs that quadratic transformers are hard to take.
\end{item}

% \section{Motivation} \label{motivation}

% [eric alcaide/hypnopump draft]
% TODO: cite 
% vanishing gradients DONE
% transformer effectiveness in parallelization vs RNN,CNN DONE
% attention cost DONE
% GPT / LLaMA (for popularity) DONE
% Sequential data processing covers a wide range of scientific and industrial applications, both directly (language, time, sequences) and indirectly (other modalities, potentially formulated as sequences including images, graphs). Models for sequential data have experienced an explosive rise in capabilities and popularity. Existing deep learning architectures for sequential data have known limitations: vanishing gradients for RNNs, slow training (RNNs, CNNs), memory and compute bottlenecks (transformers), lack of scaling (non-transformer models), difficulty handling long contexts (CNNs, transformers), lack of parallelization (RNNs), slow inference (transformers)...\\
% RWKV aims to combine strengths of previous techniques while minimizing drawbacks: parallel training and scaling like transformers but no compute/memory bottleneck ($O(n)$ vs $O(n^2)$); efficient inference like RNNs but no vanishing gradients in training. It also addresses the quadratic cost of attention in Transformer  with a reformulation to obtain "scalar attention" at linear cost without any approximation (i.e., low rank). While the approximation may be negligible with small models, it can become important when scaling up.
% Indeed, RWKV proves reliable for scaling up to billions of parameters, competitive with current SOTA at a lower cost. As such, RWKV is likely positioned to address existing challenges in further scaling and deploying AI models.
% Replace positional embedding by recurrent architecture and custom weighting of past and present information
% Reformulation of recurrence for enhanced parallelization
% Layer stacking by token/time-mixing and channel-mixing, in residual blocks
% Increased training efficiency with custom compute kernels
% Improved training dynamics by custom embedding initialization

The surge of deep learning techniques in recent years has led to impressive advancements in many scientific and industrial applications involving sequential data processing, such as natural language understanding \cite{brown2020language}, time-series analysis \cite{ismail2019deep}, and indirect modalities that can be reframed as sequences, including images and graphs \cite{wu2020comprehensive}. This progress has been catalyzed by deploying neural network architectures such as RNNs, convolutional neural networks (CNNs), and, more recently, Transformers \cite{NIPS2017_3f5ee243}. The development and success of advanced Transformer models such as GPT-3 \cite{brown2020language}, ChatGPT \cite{openai2022chatgpt,kocon2023chatgpt}, GPT-4 \cite{openai2023gpt4}, LLaMA \cite{touvron2023llama}, and Chinchilla \cite{hoffmann2022training}, underline the significance of this architecture in pushing the boundaries of what is possible in natural language processing and related fields. 

However, these architectures have inherent limitations that may restrict their usage in specific scenarios. RNNs, for instance, suffer from the vanishing gradient problem \cite{hochreiter1998vanishing,le2016quantifying} and cannot parallelize during training, limiting their scalability. CNNs, while excellent at capturing local patterns, struggle with long-range dependencies, a crucial aspect of many sequence processing tasks \cite{bai2018empirical}. On the other hand, transformers have proven to be highly effective at capturing both local and long-range dependencies and allow parallelized training \cite{tay2022efficient}. Still, their self-attention mechanism comes with significant memory and computational cost, which scales quadratically with the length of the input sequence, making them unsuitable for long sequences and settings with constrained computational resources \cite{katharopoulos2020transformers}.

To address these limitations, we introduce RWKV, a novel neural network architecture combining previous techniques' strengths while mitigating their drawbacks. RWKV offers parallelized training and robust scaling capabilities akin to Transformers but without their quadratic computational and memory bottleneck, reducing the complexity from $O(n^2)$ to $O(n)$. Moreover, it provides efficient inference, reminiscent of RNNs, without suffering from vanishing gradients during training. 

The attention through dot-product token interaction has been eschewed in favor of channel-directed attention, meaning that a token’s attention is more meaningfully quantified by which data channels are prioritized over others. This is in contrast with Transformer architecture where attention is more directly affected by specific token interactions.

Crucially, RWKV reformulates the attention mechanism resulting in \emph{linear attention}, which incurs a linear cost without approximation. While the effects of approximation in the attention mechanism might be negligible for small models, they become increasingly significant when scaling up.

This work demonstrates that RWKV can handle large-scale models with billions of parameters and is competitive with state-of-the-art architectures at a significantly reduced computational cost. Our experimental results suggest that RWKV could be a powerful tool to tackle the ongoing challenges in scaling and deploying AI models in various domains, particularly those involving sequential data processing. Thus, RWKV paves the way toward more sustainable and efficient AI models for sequence processing tasks by bridging the gap between computational efficiency and expressive capacity.
% ============================